\documentclass[a4paper,10pt]{amsart}
\usepackage{style}
\title{Trivializing the $S^1$ action}

\begin{document}

\maketitle

\tableofcontents

Let $\Top$ be the category of compactly generated hasudorff topological spaces equipped with the standard model structure i.e.  with the weak equivalences given by weak homotopy equivalences and fibrations given by Serre fibrations. Let $\redTopSeq$ denote the category $\Top^{\ZZ \times \ZZ \setminus (0,0)}$ of bi-sequences of topological spaces , with a model structure induced from $\Top$. Here `reduced' refers to the fact these sequences do not contain a $(0,0)$ component.


%section
\section{Review of Reduced Properads}
In this section we present a brief overview of reduced properads in the category of topological spaces. For a detailed treatment see \cite{vallette2007koszul}. We shall define properads as algebras over a certain moand $\GG$. This point of view will be convenient when we discuss the model structure and cofibrant resolutions of properads below.

%We start by recalling the definition of a monad and an algebra over it:
\subsection{Monads and algebras over them}
\begin{definition}\label{tripledef}
A \textbf{monad} (also called a \textbf{triple}) $T$ on a category $\catc$, is a functor $T \colon \catc \to \catc$ along with natural transformations $\mu \colon T\circ T \to T$ and $\eta \colon \mathbbm{1}_{\catc} \to T$ such that the following diagrams commute:

\[
\begin{tikzcd}
T \circ T \circ T \ar{r}{\mu \circ 1} \ar{d}{\mu \circ 1} &T \circ T \ar{d}{\mu} \\
T \circ T \ar{r}{\mu} & T\\
\end{tikzcd} \mbox{ \	 \	}
\begin{tikzcd}
\mathbbm{1}_{\catc} \circ T\ar{r}{\eta \circ 1} \ar[equals]{rd} &T \circ T \ar{d}{\mu} &T \circ \mathbbm{1}_{\catc} \ar{l}{1 \circ \eta} \ar[equals]{ld}\\
 &T & \\
\end{tikzcd}
\]
The natural transformations $\mu$ and $\eta$ are referred to as the composition and the unit of $T$, respectively. The first commutative diagram describes the associativity of $\mu$, whereas the second one describes the unitality of the $\eta$. 

An \textbf{algebra over }$T$ is an object $X \in \catc$ equipped with a map $m \colon TX \to X \in \catc$ such that the following diagrams commute: 
\[
\begin{tikzcd}
T T X \ar{r}{\mu (X)} \ar{d}{Tm} &TX\ar{d}{m} \\
TX  \ar{r}{m} & T\\
\end{tikzcd} \mbox{ \	 \	}
\begin{tikzcd}
X \ar{r}{\eta (X)}\ar[equals]{rd} &T X \ar{d}{m}\\
 &X & \\
\end{tikzcd}
\]
\end{definition}
We denote the category of $T$-algebras in $\catc$ by \textbf{T-Alg}. For any $X \in \catc$, $TX$ is canonically a $T$-algebra. This defines the free $T$-algebra functor $T \colon \catc \to$ \textbf{T-Alg} left adjoint to the forgetful functor \textbf{T-Alg} $\to \catc$.


\subsection{Graphs}
To define the monad $\GG$, we need to fix some conventions regarding graphs. We follow the terminology from \cite{getzler1998modular}.

By a \textbf{graph} $G$ we mean the data of of a set $Flag(G)$ together with an involution $\sigma$ and a partition $\lambda$. 

Elements of $Flag(G)$ are referred to as \textbf{flags}. An equivalence class of $\lambda$ is called a \textbf{vertex} of $G$. The set of vertices of $G$ is denoted by $V(G)$. An orbit of $\sigma$ containing $2$ distinct elements is called an \textbf{edge}, whereas orbit consisting of a single element is called a \textbf{leg} of $G$. The set of edges and legs of $G$ are denoted by $E(G)$ and $L(G)$ respectively. 

An isomorphism between graphs $G$ and $G^\prime$ is a bijection $\phi \colon Flag(G) \to Flag(G^\prime)$ compatible with the partitions and the involution.

The \textbf{topological graph} associated to a graph $G$ is a $1$ dimensional CW-complex, obtained as follows: Take a copy of $[0,1]$ for every flag. Identify the points $0 \in [0,1]$ for all flags in an equivalence class of the partition $\lambda$, and identify the points $1 \in [0,1]$ for every pair of flags exchanged by the involution $\sigma$.

We say that a graph $G$ is \textbf{connected} if the topological graph associated to it is connected.

By a \textbf{directed graph} we mean a \emph{connected} graph $G$ along with the data of an orientation on all the edges and legs of the graph. The set of edges and legs directed into (respectively, out of) a vertex $v$ are denoted $in(v)$ (respectively, $out(v)$) and is referred to as \emph{inputs (respectively, outputs)}of $v$. The set of legs directed towards (respectively, away from) a vertex of the graph is denoted $in(G)$ (respectively, $out(G)$) and is referred to as \emph{inputs (respectively, outputs)} 	of $G$. In addition, we assume that $in(G), out(G)$ as well as $in(v),out(v)$ for every $v \in V(G)$, are all equipped with orderings.

An isomorphism of directed graphs is an isomorphism of the underlying graphs which preserves the orientations on the edges and legs, as well as the orderings on the inputs and outputs of the vertices and the graph.

A directed graph $G$ is called \textbf{acylic} if it contains no directed cycles. 

Finally, we will say that a directed graph is \textbf{reduced} if either $in(G) \neq \varnothing$ or $out(G) \neq \varnothing$. Below we will refer to reduced reduced acyclic graph as \emph{rda-graphs}.

%%%%%%%%%%%%%%%%%%%%%%%%%%%%%%%%%%%%%%%%%%%%%%%%%%%%%%
%%%%%%%%%%%%%%%%                      Subsection              %%%%%%%%%%%%%%%%%%%%%%%

\subsection{The monad $\GG$}\label{GGtripsubsec}
For a da-graph $G$ and $P \in \redTopSeq$, let 
\[ P(G) := \prod_{v \in V(G)} P(|in(v)|,|out(v)|) \]
denote the set of $P$-decoration of $G$. Denote by $\mathcal{G}_{n_-,n_+}$ the collection of graphs one from each isomorphism class of rda-graphs with $n_-$ input and $n_+$ output leaves. Set 
\[ \GG_{n_-,n_+}(P) := \coprod_{G \in \mathcal{G}_{n_-,n_+}}P(G)\]
Denote by $\GG(P) \in \redTopSeq$ the bi-sequence given by $\GG(P) (n_-,n_+) = \GG_{n_-,n_+}(P)$. 


We claim that there is a natural structure of a monad on the functor $\GG \colon \redTopSeq \to \redTopSeq$. To see this we need to define natural transformations $\mu \colon \GG \GG \to \GG$ and $\eta \colon \mathbbm{1} \to \GG$, satisfying the condition mentioned in Definition	 \ref{tripledef}.\\
For $P \in \redTopSeq$, $\GG\GG(P)$ is given by
\[ \GG\GG(P)(n_-,n_+) = \coprod_{G\in \mathcal{G}_{n_-,n_+}} \GG(P) (G) \]
Thus an element in $\GG\GG(P)$ is represented by a $\GG(P)$-decoration of an rda-graph. Such a decorated graph can be identified with a `nested' rda-graph. Forgetting the nesting gives an rda-graph with $P$-decorations. This defines the natural transformation $\mu \colon \GG\GG \to \GG$ (see ?? for details). It is not difficult to verify that this satisfies the associativity condition in Definition \ref{tripledef}.\\
Let us now turn to the unit. Given $P \in \redTopSeq$, consider the maps $P(n_-,n_+) \to \GG_{n_-,n_+}(P)$ which identify $P(n_-,n_+)$ with the $P$-decorations of $(n_-,n_+)$-corolla, the unique graph in $\mathcal{G}_{n_-,n_+}$ having a single vertex and no internal edges. This defines a functor $\eta \colon \mathbbm{1} \to \GG$. Again, it is not difficult to verify that $\eta$ satisfies the unitality conditions in Definition \ref{tripledef}.\\
We thus conclude, $(\GG,\mu,\eta)$ is a monad.

%%%%%%%%%%%%%%%%%%%%%%%%%%%%%%%%%%%%%%%%%%%%%%%%%%%%%%
%%%%%%%%%%%%%%%%                      Section                 %%%%%%%%%%%%%%%%%%%%%%%


\subsection{Reduced Properads as monad algebras}
\begin{definition}
	A \textbf{Reduced Properad} in $\Top$ is an algebra over the monad $(\GG, \mu, \eta)$ in the category $\redTopSeq$.
\end{definition}



%%%%%%%%%%%%%%%%%%%%%%%%%%%%%%%%%%%%%%%%%%%%%%%%%%%%%%
%%%%%%%%%%%%%%%%                      Figure                 %%%%%%%%%%%%%%%%%%%%%%%

\begin{figure} 
	\centering
	\begin{subfigure}[b]{0.3\textwidth}
		\includegraphics[width=\textwidth]{../img/perm_corr1.png}
		\caption{a permuted corolla}
		\label{perm_corr1}
	\end{subfigure}
	\begin{subfigure}[b]{0.35\textwidth}
		\includegraphics[width=\textwidth]{../img/perm_corr2.png}
		\caption{a partially grafted corolla}
		\label{perm_corr2}
	\end{subfigure}
	\caption{}
\end{figure}

%%%%%%%%%%%%%%%%%%%%%%%%%%%%%%%%%%%%%%%%%%%%%%%%%%%%%%
%%%%%%%%%%%%%%%%%%%%%%%%%%%%%%%%%%%%%%%%%%%%%%%%%%%%%%

\begin{remark}
	Therefore if $P$ is a properad, for any da-graph $G$ we have a map
	\[ P(G) \to P(in(G),out(G))\]
	we will refer to it as the \textbf{$\GG$-algebra composition indexed by $G$} and denote it by $\mu_G$. In particular, we have the following structure maps:
\begin{enumerate}
\item \label{proprmk1} Using the graphs as in Figure \ref{perm_corr1}), we obtain a left $\Sigma_{n_-}$ and a right $\Sigma_{n_+}$ action on $P(n_-,n_+)$.
\item \label{proprmk2} Using the graphs as in Figure \ref{perm_corr2}), for any $n_-,n_+,m_-,m_+ \geq 0$ and \emph{non-empty} subsets $s_-$ and $s_+$ of $[n_+]$ and $[m_-]$ respectively, with an identification $\phi \colon s_- \to s_+$, we get a composition map
 \[ \circ_{\phi} \colon P(n_-,n_+) \times P(m_-,m_+) \to P(n_-+m_--|s_+|,n_++m_+-|s_-|)\]
provided $n_-,m_- - |s_+|, n_+- |s_-|$, and $m_+$ are not all simultaneously $0$. This slightly unnatural constraint is imposed by the fact that we work with reduced properads which have no $(0,0)$ component.\\
These compositions are associative and are equivariant with respect to the $\Sigma \times \Sigma^{op}$-actions, described in point \eqref{proprmk1}, in a suitable sense.
\end{enumerate}
Conversely, it can be shown that $\Sigma \times \Sigma^{op}$-actions as in point \eqref{proprmk1} and (associative, $\Sigma \times \Sigma^{op}$-equivariant) composition maps as in point \eqref{proprmk2} define the structure of a $(\GG,\mu,\eta)$-algebra over $P$. We will not use this view point and hence omit the details.
\end{remark}

\section{Model Structure on Reduced Properads}
We now construct a model structure on the category of reduced properads in $\Top$. We will use a criteria descried in \cite{schwede2000algebras} (see Proposition \ref{modprop}) to transfer the model structure from $\redTopSeq$ to the category of $\GG$-algebras in $\redTopSeq$:

\subsection{Transferred model structure}%The Criterion}
Let $(T,\mu,\eta)$ be a monad	 acting on a closed model category $\catc$ and $X \in \catc$ be a $T$-alg . Then, we define a \textbf{path object} of $X$ to be an $T$-alg $X^{I}$ together with maps 
\[ X \xrightarrow{\sim} X^I \xrightarrow{} X \times X\]
factoring the diagonal $X \xrightarrow{\Delta} X \times X$, such that the first map is a weak equivalence and the second map is a fibration in $\catc$.\\
Also, recall that if $\catc$ is in addition cocomplete and if $V$ is an object and $\mathcal{D}$ a subcategory of $\catc$, we say that $V$ is \textbf{small relative to $\mathcal{D}$} if there exists a regular cardinal $\kappa$ such that for every $\lambda \geq \kappa$ and every functor $W \colon \lambda \to \mathcal{D}$, 
\[ \varinjlim_{\lambda} Hom_{\catc}(V,W_{\lambda}) \xrightarrow{\sim} Hom_{\catc}(V,\varinjlim_{\lambda}W_{\lambda}) \]
is an isomorphism.\\ % \textcolor{red}{When, $\catc=\Top$ and hence when $\catc = \redTopSeq$, it suffices to consider sequential limits.}
Finally, if $I$ is a collection of morphisms in $\catc$, by \textbf{regular $I$-cofibrations}, denoted $I-cof_{reg}$, we mean the class of (possibly transfinite) compositions of pushouts of maps in $I$.

\begin{proposition}[{\cite[Lemma 2.3]{schwede2000algebras}}]\label{modprop}
	Suppose that $\catc$ is a cofibrantly generated category with generating cofibrations $I$ and generating acyclic cofibrations $J$. Let $(T,\mu,\eta)$ be a monad acting on $\catc$ and let $I_T, J_T$ be the images of $I,J$ respectively under the free $T$-alg functor. Now suppose that
	\begin{enumerate}
		\item $T \colon \catc \to \catc$ commutes with filtered colimits
		\item domains of $I_T,J_T$ are small relative to $I_T-cof_{reg}, J_T-cof_{reg}$ respectively
		\item atleast one of the following conditions is satisfied: 
			\begin{itemize}
				\item every regular $J_T$-cofibration is a weak equivalence
				\item every object of $\catc$ is fibrant and every $T$-algebra has a path object
			\end{itemize}
	\end{enumerate} 
Then the category $T-Alg$ is a cofibrantly generated model category with $I_T$ a generating set of cofibrations and $J_T$ a generating set of acyclic cofibrations.
\end{proposition}

We will apply this to the case of the functor $\GG$ acting on the category $\redTopSeq$. \\ 
Recall that the model structure on $\redTopSeq$ cofibrantly generated, with generating cofibrations $I=\{\{\partial \DD^{k_{n_-,n_+}} \to \DD^{k_{n_-,n_+}} \}_{n_-,n_+}\}$ and generating acyclic cofibrations $J=\{ \{\DD^{k_{n_-,n_+}}\times \{0\} \to \DD^{k_{n_-,n_+}} \times \DD^1  \}_{n_-,n_+}\}$.\\
 The condition on filtered colimits follows directly from the definition of $\GG$.\\
Also, note that every object in $\redTopSeq$ is fibrant and that for any $\GG$-Alg $X = \{X_{n_-n_+}\} \in \redTopSeq$, $X^I = \{ Hom_{\Top}(I,X_{n_-,n_+}	)\}$ has a natural structure of a $\GG$-Alg and is thus a path object for $X$. It remains to verify the smallness condition on the domains of $I_\GG$ and $J_{\GG}$. The key technical input towards this is 

the following Lemma:

\begin{lemma}\label{lempush}
Let $f \colon X \to Y \in cofib(\redTopSeq)$ and $A$ be a $\Galg$ which is cofibrant in $\redTopSeq$. Let $\alpha \colon X \to A$ be a map of topological sequences. Denote the induced map $\GG(X) \to A$ by $\widetilde{\alpha}$. Let $A^\prime$ be the pushout
\[
	\begin{tikzcd}
	\GG(X) \ar{r}{\widetilde{\alpha}} \ar{d} & A \ar{d} \\
	\GG(Y) \ar{r} & A^\prime
	\end{tikzcd}
\]
Then the map $A \to A^\prime$ is an inclusion of topological sequences.
\end{lemma}
  
Recall that $f \colon X \to Y \in \Top$ is called an \textbf{inclusion} if it is injective and any $U \subset X$ is open if and only if there exists $V \subset Y$ such that $U = f^{-1}(V)$. In other words, $f$ maps $X$ homeomorphically onto its image in $Y$. A map of topological sequences is called an \textbf{inclusion} if it is an inclusion in each component.

It will be convenient to set up some notation before we start the proof of Lemma \ref{lempush}.
\begin{notation}
	Let $m \geq 0$ and $\nbar \in (\ZZ_{\geq 0} \times \ZZ_{\geq 0})^m$ be an $m$-term sequence of input-output profiles $(n_-^{(1)},n_+^{(1)}),....,(n_-^{(m)},n_+^{(m)})$. Let $\Sigma_{\nbar}$ be the group of permutations $\sigma$ of $[m]=\{1,...,m\}$ for which $(n_-^{(i)},n_+^{(i)})$ and $(n_-^{(\sigma(i))},n_+^{(\sigma(i))})$ coincide for all $i$.
	
	For $Z \in \redTopSeq$, let $Z^{\nbar}$ denote the product $\prod_{i=1}^m Z(n_-^{(i)},n_+^{(i)})$.
	
	Let $Cube^m$ denote the $m$-cube category i.e. the poset category of the subsets of $[m]$. Consider the functor 
	\begin{align*}
		Q_{\nbar} \colon Cube^m & \to \redTopSeq\\
		W &\mapsto Q_{\nbar}(W) := X^{\nbar_W} \times Y^{\nbar_{[m] \setminus W}}
	\end{align*}
	(here $\nbar_{W}$ and $\nbar_{[m] \setminus W}$ denote the subsequences of $\nbar$ consisting of terms indexed by $W$ and $[m] \setminus W$ respectively). For inclusions of subsets $W_1 \subset W_2$, the maps $Q_{\nbar}(W_1) \to Q_{\nbar}(W_2)$ are defined using the components of $f \colon X \to Y$. Set
	\[ Q_{\nbar}(X,Y) := \varinjlim_{W \in Cube^m \setminus \{[m]\}} Q_{\nbar}(W) \]
	Note that there is a canonical map $Q_{\nbar}(X,Y) \to Y^{\nbar}$.
	
	Now, denote by $\mathcal{G}_{\nbar}$ the collection of isomorphism classes of pairs $(G,\gamma)$ where $G$ is an rda-graph and $\gamma \colon [m] \to V(G)$ such that $\gamma(i)$ is a vertex of type $(n_-^{(i)},n_+^{(i)})$.
%	groupoid whose objects are rda-graphs $G$ together with a map $\gamma \colon [m] \to V(G)$ such that $\gamma(i)$ is a vertex of type $(n_-^{(i)},n_+^{(i)})$. The maps of $\mathcal{G}_{\nbar}$ are the isomorphisms of such objects. For $(G,\gamma) \in \mathcal{G}_{\nbar}$,
	 Set 
	\[ (G,\gamma)(Z) = \prod_{v\in V(G) \setminus \gamma([m])} Z(|in(v)|,|out(v)|) \]
	Finally define functors 
	\[ \GG_{\nbar}, \BB_{\nbar} \colon \redTopSeq \to \redTopSeq \]
	as follows:
	\begin{align*} 
	\GG_{\nbar}(Z) &= \coprod_{(G,\gamma) \in \mathcal{G}_{\nbar}} (G,\gamma)(Z) \\
	\BB_{\nbar}(Z) &= coeq^{\redTopSeq}\Big[ \GG_{\nbar}(\GG Z) \rightrightarrows \GG_{\nbar}(Z) \Big]
	\end{align*}
\end{notation}
\begin{proof}[Proof of Lem. \ref{lempush}]
	We prove the lemma by expressing $A^\prime$ as a colimit
	\[A^\prime = \varinjlim ( A^\prime_{(0)} \to A^\prime_{(1)} \to \hdots A^\prime_{(m)} \to \hdots) \]
	such that $A^\prime_{(0)} =A$ and each $A^\prime_{(i)} \to A^\prime_{(i+1)}$ is an inclusion.
	The sequence of morphisms is constructed inductively as follows:
	\begin{itemize}
		\item $A^\prime_{(0)} = A$
		\item $A^\prime_{(m)}$ is defined as a pushout as follows:
		\[
			\begin{tikzcd}\label{lempushdig}
			\coprod_{\nbar \in (\ZZ_{\geq 0} \times \ZZ_{\geq 0})^m} \BB_{\nbar}(A) \times_{\Sigma_{\nbar}} Q_{\nbar} (X,Y) \ar{r}\ar{d} & \coprod_{\nbar \in (\ZZ_{\geq 0} \times \ZZ_{\geq 0})^m} \BB_{\nbar}(A) \times_{\Sigma_{\nbar}} Y^{\nbar} \ar{d}\\
			A^\prime_{(m-1)} \ar{r} &A^\prime_{(m)}
			\end{tikzcd}
		\]
	\end{itemize}
	Using the fact that $X \to Y$ is an inclusion, it follows that the upper horizontal arrows in each of the diagrams \ref{lempushdig} are inclusions. Since the class of inclusions is closed under pushouts, it follows that $A^\prime_{(m-1)} \to A^\prime_{(m)}$ is an inclusion for every $m$.
	
	This completes the proof of the Lemma.
\end{proof}

To prove the smallness hypothesis in Prop. \ref{modprop} for the case of the monad $\GG$ acting on $\redTopSeq$, we require two more easy Lemmas. The proof of these statements are straightforward and the details are left to the reader:
\begin{lemma}\label{modproplem1}
Let $\lambda$ be a regular cardinal and $W \colon \lambda \to \Galg$ be a functor. Then, the $\redTopSeq$-colimit $\varinjlim_{\lambda}^{\redTopSeq} W_{\lambda}$ has a natural $\Galg$-structure and hence coincides with the $\Galg$-colimit $\varinjlim_{\lambda}^{\Galg} W_{\lambda}$
\end{lemma}
\begin{lemma}\label{modproplem2}
Let $f \colon V \to W$ be a $\Galg$ map such that the underlying map of topological sequences is an inclusion. Then the square
\[	
		\begin{tikzcd}
			Hom^{\Galg}(A,V) \ar{r}\ar{d} & Hom^{\Galg}(A,W) \ar{d} \\
			Hom^{\TopSeq}(A,V) \ar{r} & Hom^{\TopSeq}(A,W)
		\end{tikzcd}
	\]
is a fiber pullback (in sets).
\end{lemma}

%We are now ready to verify the smallness hypothesis in \ref{modprop} and prove the existence of transferred model structure on $\Galg$:
%
%We now return to verifying the smallness hypothesis in the Prop. \ref{modprop} for the case of the monad $\GG$ acting on $\redTopSeq$. Observe that, since the class of inclusions is closed under (transfinite) compositions and pushouts, it follows from Lem. \ref{lempush} that the $\redTopSeq$-maps underlying $I_{\GG}$ and $J_{\GG}$ are inclusions. Now, we have the following fact:
%\begin{fact} 
%	Every object in $\Top$ is small relative to the inclusions \textcolor{red}{what about objects in $\Galg$? $\GG$ commutes with filtered colimits, so $\Top$ filtered colimit has a natural $\GG$-algebra structure, so filtered colimits in $\Galg$ and $\Top$ coincide. Filtered colimits commute with finite limits and for $A, X ,Y$ $\GG$-algebras and $X \to Y$ map of $\Galg$, the induced diagram}
%	\[	
%		\begin{tikzcd}
%			Hom^{\Galg}(A,X) \ar{r}\ar{d} & Hom^{\Galg}(A,Y) \ar{d} \\
%			Hom^{\TopSeq}(A,X) \ar{r} & Hom^{\TopSeq}(A,Y)
%		\end{tikzcd}
%	\]
%\textcolor{red}{is a fiber diagram (in sets).}
%\end{fact}
%It follows that the smallness conditions on the domains of $I_{\GG}$ and $J_{\GG}$ are satisfied \emph{in $\redTopSeq$}. Towards showing the smallness condition in $\Galg$, let $\mathcal{D}$ denote the subcategory of $\Galg$ corresponding to either $I_{\GG}$ or $J_{\GG}$. Let $A$ be a $\GG$ algebra and $\kappa$ be the cardinal corresponding to $A$ given by the fact that $A$ is small relative to $\mathcal{D}$ in $\redTopSeq$. Let $\lambda \geq \kappa$ and let $W: \lambda \to \mathcal{D}$ be a functor. From the $\redTopSeq$ smallness we know that 
%\[ \varinjlim_{\lambda} Hom^{\redTopSeq}(A,W_{\lambda}) \to  Hom^{\redTopSeq}(A,\varinjlim^{\redTopSeq}_{\lambda} W_{\lambda}) \]
%is an isomorphism. The filtered colimit $\varinjlim^{\redTopSeq}_{\lambda} W_{\lambda})$ a natural $\Galg$ structure and hence it follows that
%\[  \varinjlim^{\redTopSeq}_{\lambda} W_{\lambda}) \simeq \varinjlim^{\Galg}_{\lambda} W_{\lambda})\]


Thus, applying Prop. \ref{modprop} to $\GG$ acting on $\redTopSeq$, we have
\begin{corollary}
	The category $\Galg$ has a cofibrantly generated model structure, with the generating cofibrations (respectively, acyclic cofibrations) $I_{\GG}$ (respectively, $J_{\GG}$). The weak equivalences (respectively, fibrations) are give by the weak equivalence (respectively, fibrations) of the underlying topological sequences.
\end{corollary}
\begin{proof}
First, observe that since the class of inclusions is closed under (transfinite) compositions and pushouts, it follows from Lemma \ref{lempush} that the $\redTopSeq$-maps underlying $I_{\GG}$ and $J_{\GG}$ are inclusions. Now, we have the following fact:
\begin{fact} 
	Every object in $\Top$ is small relative to the inclusions 
\end{fact}
Fix a $\Galg$ $A$. Using the fact above, it follows that $A$ is small relative to  $I_{\GG}$ \emph{in $\redTopSeq$}. Let $\kappa$ be the cardinal guaranteed by this corresponding to $A$. Let $\lambda \geq \kappa$ and $W \colon \lambda \to I_{\GG}$ be a functor. Using Lemma \ref{modproplem1} it follows that 
\[ \varinjlim_{\lambda} Hom^{\redTopSeq}(A, W_{\lambda}) \xrightarrow{\sim} Hom^{\Galg}(A, \varinjlim^{\Galg}_{\lambda} W_{\lambda}) \]
is an isomorphism. Since $\redTopSeq$-maps underlying $I_{\GG}$ are inclusions, we know that the $\redTopSeq$-map underlying 
\[ W_{\lambda} \to  \varinjlim^{\Galg}_{\lambda} W_{\lambda}\]
is an inclusion as well. Then from Lemma \ref{modproplem2}, it follows that we have a pullback diagram
\[	
		\begin{tikzcd}
			Hom^{\Galg}(A, W_{\lambda}) \ar{r}\ar{d} & Hom^{\Galg}(A,\varinjlim^{\Galg}_{\lambda} W_{\lambda}) \ar{d} \\
			Hom^{\TopSeq}(A, W_{\lambda}) \ar{r} & Hom^{\TopSeq}(A,\varinjlim^{\Galg}_{\lambda} W_{\lambda})
		\end{tikzcd}
\]
Using the fact that filtered colimits commute with limits we now conclude that 
\[ \varinjlim_{\lambda} Hom^{\Galg}(A, W_{\lambda}) \xrightarrow{\sim} Hom^{\Galg}(A, \varinjlim_{\lambda}^{\Galg} W_{\lambda})\]
is an isomorphism. This proves that $A$ is small relative to $I_{\GG}$ \emph{in $\Galg$}. Exactly the same argument shows that similar conclusion holds for $J_{\GG}$. This completes the proof of the fact that monad $\GG$ acting on $\redTopSeq$ satisfies the hypothesis of Proposition \ref{modprop}. The conclusion of the corollary now follows from the Proposition.
\end{proof}



\section{Model Structure on Reduced Properads}
We now construct a model structure on the category of reduced properads in $\Top$. We will use a criteria descried in \cite{schwede2000algebras} (see Proposition \ref{modprop}) to transfer the model structure from $\redTopSeq$ to the category of $\GG$-algebras in $\redTopSeq$:

\subsection{Transferred model structure}%The Criterion}
Let $(T,\mu,\eta)$ be a monad	 acting on a closed model category $\catc$ and $X \in \catc$ be a $T$-alg . Then, we define a \textbf{path object} of $X$ to be an $T$-alg $X^{I}$ together with maps 
\[ X \xrightarrow{\sim} X^I \xrightarrow{} X \times X\]
factoring the diagonal $X \xrightarrow{\Delta} X \times X$, such that the first map is a weak equivalence and the second map is a fibration in $\catc$.\\
Also, recall that if $\catc$ is in addition cocomplete and if $V$ is an object and $\mathcal{D}$ a subcategory of $\catc$, we say that $V$ is \textbf{small relative to $\mathcal{D}$} if there exists a regular cardinal $\kappa$ such that for every $\lambda \geq \kappa$ and every functor $W \colon \lambda \to \mathcal{D}$, 
\[ \varinjlim_{\lambda} Hom_{\catc}(V,W_{\lambda}) \xrightarrow{\sim} Hom_{\catc}(V,\varinjlim_{\lambda}W_{\lambda}) \]
is an isomorphism.\\ % \textcolor{red}{When, $\catc=\Top$ and hence when $\catc = \redTopSeq$, it suffices to consider sequential limits.}
Finally, if $I$ is a collection of morphisms in $\catc$, by \textbf{regular $I$-cofibrations}, denoted $I-cof_{reg}$, we mean the class of (possibly transfinite) compositions of pushouts of maps in $I$.

\begin{proposition}[{\cite[Lemma 2.3]{schwede2000algebras}}]\label{modprop}
	Suppose that $\catc$ is a cofibrantly generated category with generating cofibrations $I$ and generating acyclic cofibrations $J$. Let $(T,\mu,\eta)$ be a monad acting on $\catc$ and let $I_T, J_T$ be the images of $I,J$ respectively under the free $T$-alg functor. Now suppose that
	\begin{enumerate}
		\item $T \colon \catc \to \catc$ commutes with filtered colimits
		\item domains of $I_T,J_T$ are small relative to $I_T-cof_{reg}, J_T-cof_{reg}$ respectively
		\item atleast one of the following conditions is satisfied: 
			\begin{itemize}
				\item every regular $J_T$-cofibration is a weak equivalence
				\item every object of $\catc$ is fibrant and every $T$-algebra has a path object
			\end{itemize}
	\end{enumerate} 
Then the category $T-Alg$ is a cofibrantly generated model category with $I_T$ a generating set of cofibrations and $J_T$ a generating set of acyclic cofibrations.
\end{proposition}

We will apply this to the case of the functor $\GG$ acting on the category $\redTopSeq$. \\ 
Recall that the model structure on $\redTopSeq$ cofibrantly generated, with generating cofibrations $I=\{\{\partial \DD^{k_{n_-,n_+}} \to \DD^{k_{n_-,n_+}} \}_{n_-,n_+}\}$ and generating acyclic cofibrations $J=\{ \{\DD^{k_{n_-,n_+}}\times \{0\} \to \DD^{k_{n_-,n_+}} \times \DD^1  \}_{n_-,n_+}\}$.\\
 The condition on filtered colimits follows directly from the definition of $\GG$.\\
Also, note that every object in $\redTopSeq$ is fibrant and that for any $\GG$-Alg $X = \{X_{n_-n_+}\} \in \redTopSeq$, $X^I = \{ Hom_{\Top}(I,X_{n_-,n_+}	)\}$ has a natural structure of a $\GG$-Alg and is thus a path object for $X$. It remains to verify the smallness condition on the domains of $I_\GG$ and $J_{\GG}$. The key technical input towards this is the following Lemma:

\begin{lemma}\label{lempush}
Let $f \colon X \to Y \in cofib(\redTopSeq)$ and $A$ be a $\Galg$ which is cofibrant in $\redTopSeq$. Let $\alpha \colon X \to A$ be a map of topological sequences. Denote the induced map $\GG(X) \to A$ by $\widetilde{\alpha}$. Let $A^\prime$ be the pushout
\[
	\begin{tikzcd}
	\GG(X) \ar{r}{\widetilde{\alpha}} \ar{d} & A \ar{d} \\
	\GG(Y) \ar{r} & A^\prime
	\end{tikzcd}
\]
Then the map $A \to A^\prime$ is an inclusion of topological sequences.
\end{lemma}
  
Recall that $f \colon X \to Y \in \Top$ is called an \textbf{inclusion} if it is injective and any $U \subset X$ is open if and only if there exists $V \subset Y$ such that $U = f^{-1}(V)$. In other words, $f$ maps $X$ homeomorphically onto its image in $Y$. A map of topological sequences is called an \textbf{inclusion} if it is an inclusion in each component.

It will be convenient to set up some notation before we start the proof of Lemma \ref{lempush}.
\begin{notation}
	Let $m \geq 0$ and $\nbar \in (\ZZ_{\geq 0} \times \ZZ_{\geq 0})^m$ be an $m$-term sequence of input-output profiles $(n_-^{(1)},n_+^{(1)}),....,(n_-^{(m)},n_+^{(m)})$. Let $\Sigma_{\nbar}$ be the group of permutations $\sigma$ of $[m]=\{1,...,m\}$ for which $(n_-^{(i)},n_+^{(i)})$ and $(n_-^{(\sigma(i))},n_+^{(\sigma(i))})$ coincide for all $i$.
	
	For $Z \in \redTopSeq$, let $Z^{\nbar}$ denote the product $\prod_{i=1}^m Z(n_-^{(i)},n_+^{(i)})$.
	
	Let $Cube^m$ denote the $m$-cube category i.e. the poset category of the subsets of $[m]$. Consider the functor 
	\begin{align*}
		Q_{\nbar} \colon Cube^m & \to \redTopSeq\\
		W &\mapsto Q_{\nbar}(W) := X^{\nbar_W} \times Y^{\nbar_{[m] \setminus W}}
	\end{align*}
	(here $\nbar_{W}$ and $\nbar_{[m] \setminus W}$ denote the subsequences of $\nbar$ consisting of terms indexed by $W$ and $[m] \setminus W$ respectively). For inclusions of subsets $W_1 \subset W_2$, the maps $Q_{\nbar}(W_1) \to Q_{\nbar}(W_2)$ are defined using the components of $f \colon X \to Y$. Set
	\[ Q_{\nbar}(X,Y) := \varinjlim_{W \in Cube^m \setminus \{[m]\}} Q_{\nbar}(W) \]
	Note that there is a canonical map $Q_{\nbar}(X,Y) \to Y^{\nbar}$.
	
	Now, denote by $\mathcal{G}_{\nbar}$ the collection of isomorphism classes of pairs $(G,\gamma)$ where $G$ is an rda-graph and $\gamma \colon [m] \to V(G)$ such that $\gamma(i)$ is a vertex of type $(n_-^{(i)},n_+^{(i)})$.
%	groupoid whose objects are rda-graphs $G$ together with a map $\gamma \colon [m] \to V(G)$ such that $\gamma(i)$ is a vertex of type $(n_-^{(i)},n_+^{(i)})$. The maps of $\mathcal{G}_{\nbar}$ are the isomorphisms of such objects. For $(G,\gamma) \in \mathcal{G}_{\nbar}$,
	 Set 
	\[ (G,\gamma)(Z) = \prod_{v\in V(G) \setminus \gamma([m])} Z(|in(v)|,|out(v)|) \]
	Finally define functors 
	\[ \GG_{\nbar}, \BB_{\nbar} \colon \redTopSeq \to \redTopSeq \]
	as follows:
	\begin{align*} 
	\GG_{\nbar}(Z) &= \coprod_{(G,\gamma) \in \mathcal{G}_{\nbar}} (G,\gamma)(Z) \\
	\BB_{\nbar}(Z) &= coeq^{\redTopSeq}\Big[ \GG_{\nbar}(\GG Z) \rightrightarrows \GG_{\nbar}(Z) \Big]
	\end{align*}
\end{notation}
\begin{proof}[Proof of Lem. \ref{lempush}]
	We prove the lemma by expressing $A^\prime$ as a colimit
	\[A^\prime = \varinjlim ( A^\prime_{(0)} \to A^\prime_{(1)} \to \hdots A^\prime_{(m)} \to \hdots) \]
	such that $A^\prime_{(0)} =A$ and each $A^\prime_{(i)} \to A^\prime_{(i+1)}$ is an inclusion.
	The sequence of morphisms is constructed inductively as follows:
	\begin{itemize}
		\item $A^\prime_{(0)} = A$
		\item $A^\prime_{(m)}$ is defined as a pushout as follows:
		\[
			\begin{tikzcd}\label{lempushdig}
			\coprod_{\nbar \in (\ZZ_{\geq 0} \times \ZZ_{\geq 0})^m} \BB_{\nbar}(A) \times_{\Sigma_{\nbar}} Q_{\nbar} (X,Y) \ar{r}\ar{d} & \coprod_{\nbar \in (\ZZ_{\geq 0} \times \ZZ_{\geq 0})^m} \BB_{\nbar}(A) \times_{\Sigma_{\nbar}} Y^{\nbar} \ar{d}\\
			A^\prime_{(m-1)} \ar{r} &A^\prime_{(m)}
			\end{tikzcd}
		\]
	\end{itemize}
	Using the fact that $X \to Y$ is an inclusion, it follows that the upper horizontal arrows in each of the diagrams \ref{lempushdig} are inclusions. Since the class of inclusions is closed under pushouts, it follows that $A^\prime_{(m-1)} \to A^\prime_{(m)}$ is an inclusion for every $m$.
	
	This completes the proof of the Lemma.
\end{proof}

To prove the smallness hypothesis in Prop. \ref{modprop} for the case of the monad $\GG$ acting on $\redTopSeq$, we require two more easy Lemmas. The proof of these statements are straightforward and the details are left to the reader:
\begin{lemma}\label{modproplem1}
Let $\lambda$ be a regular cardinal and $W \colon \lambda \to \Galg$ be a functor. Then, the $\redTopSeq$-colimit $\varinjlim_{\lambda}^{\redTopSeq} W_{\lambda}$ has a natural $\Galg$-structure and hence coincides with the $\Galg$-colimit $\varinjlim_{\lambda}^{\Galg} W_{\lambda}$
\end{lemma}
\begin{lemma}\label{modproplem2}
Let $f \colon V \to W$ be a $\Galg$ map such that the underlying map of topological sequences is an inclusion. Then the square
\[	
		\begin{tikzcd}
			Hom^{\Galg}(A,V) \ar{r}\ar{d} & Hom^{\Galg}(A,W) \ar{d} \\
			Hom^{\TopSeq}(A,V) \ar{r} & Hom^{\TopSeq}(A,W)
		\end{tikzcd}
	\]
is a fiber pullback (in sets).
\end{lemma}

%We are now ready to verify the smallness hypothesis in \ref{modprop} and prove the existence of transferred model structure on $\Galg$:
%
%We now return to verifying the smallness hypothesis in the Prop. \ref{modprop} for the case of the monad $\GG$ acting on $\redTopSeq$. Observe that, since the class of inclusions is closed under (transfinite) compositions and pushouts, it follows from Lem. \ref{lempush} that the $\redTopSeq$-maps underlying $I_{\GG}$ and $J_{\GG}$ are inclusions. Now, we have the following fact:
%\begin{fact} 
%	Every object in $\Top$ is small relative to the inclusions \textcolor{red}{what about objects in $\Galg$? $\GG$ commutes with filtered colimits, so $\Top$ filtered colimit has a natural $\GG$-algebra structure, so filtered colimits in $\Galg$ and $\Top$ coincide. Filtered colimits commute with finite limits and for $A, X ,Y$ $\GG$-algebras and $X \to Y$ map of $\Galg$, the induced diagram}
%	\[	
%		\begin{tikzcd}
%			Hom^{\Galg}(A,X) \ar{r}\ar{d} & Hom^{\Galg}(A,Y) \ar{d} \\
%			Hom^{\TopSeq}(A,X) \ar{r} & Hom^{\TopSeq}(A,Y)
%		\end{tikzcd}
%	\]
%\textcolor{red}{is a fiber diagram (in sets).}
%\end{fact}
%It follows that the smallness conditions on the domains of $I_{\GG}$ and $J_{\GG}$ are satisfied \emph{in $\redTopSeq$}. Towards showing the smallness condition in $\Galg$, let $\mathcal{D}$ denote the subcategory of $\Galg$ corresponding to either $I_{\GG}$ or $J_{\GG}$. Let $A$ be a $\GG$ algebra and $\kappa$ be the cardinal corresponding to $A$ given by the fact that $A$ is small relative to $\mathcal{D}$ in $\redTopSeq$. Let $\lambda \geq \kappa$ and let $W: \lambda \to \mathcal{D}$ be a functor. From the $\redTopSeq$ smallness we know that 
%\[ \varinjlim_{\lambda} Hom^{\redTopSeq}(A,W_{\lambda}) \to  Hom^{\redTopSeq}(A,\varinjlim^{\redTopSeq}_{\lambda} W_{\lambda}) \]
%is an isomorphism. The filtered colimit $\varinjlim^{\redTopSeq}_{\lambda} W_{\lambda})$ a natural $\Galg$ structure and hence it follows that
%\[  \varinjlim^{\redTopSeq}_{\lambda} W_{\lambda}) \simeq \varinjlim^{\Galg}_{\lambda} W_{\lambda})\]


Thus, applying Prop. \ref{modprop} to $\GG$ acting on $\redTopSeq$, we have
\begin{corollary}
	The category $\Galg$ has a cofibrantly generated model structure, with the generating cofibrations (respectively, acyclic cofibrations) $I_{\GG}$ (respectively, $J_{\GG}$). The weak equivalences (respectively, fibrations) are give by the weak equivalence (respectively, fibrations) of the underlying topological sequences.
\end{corollary}
\begin{proof}
First, observe that since the class of inclusions is closed under (transfinite) compositions and pushouts, it follows from Lemma \ref{lempush} that the $\redTopSeq$-maps underlying $I_{\GG}$ and $J_{\GG}$ are inclusions. Now, we have the following fact:
\begin{fact} 
	Every object in $\Top$ is small relative to the inclusions 
\end{fact}
Fix a $\Galg$ $A$. Using the fact above, it follows that $A$ is small relative to  $I_{\GG}$ \emph{in $\redTopSeq$}. Let $\kappa$ be the cardinal guaranteed by this corresponding to $A$. Let $\lambda \geq \kappa$ and $W \colon \lambda \to I_{\GG}$ be a functor. Using Lemma \ref{modproplem1} it follows that 
\[ \varinjlim_{\lambda} Hom^{\redTopSeq}(A, W_{\lambda}) \xrightarrow{\sim} Hom^{\Galg}(A, \varinjlim^{\Galg}_{\lambda} W_{\lambda}) \]
is an isomorphism. Since $\redTopSeq$-maps underlying $I_{\GG}$ are inclusions, we know that the $\redTopSeq$-map underlying 
\[ W_{\lambda} \to  \varinjlim^{\Galg}_{\lambda} W_{\lambda}\]
is an inclusion as well. Then from Lemma \ref{modproplem2}, it follows that we have a pullback diagram
\[	
		\begin{tikzcd}
			Hom^{\Galg}(A, W_{\lambda}) \ar{r}\ar{d} & Hom^{\Galg}(A,\varinjlim^{\Galg}_{\lambda} W_{\lambda}) \ar{d} \\
			Hom^{\TopSeq}(A, W_{\lambda}) \ar{r} & Hom^{\TopSeq}(A,\varinjlim^{\Galg}_{\lambda} W_{\lambda})
		\end{tikzcd}
\]
Using the fact that filtered colimits commute with limits we now conclude that 
\[ \varinjlim_{\lambda} Hom^{\Galg}(A, W_{\lambda}) \xrightarrow{\sim} Hom^{\Galg}(A, \varinjlim_{\lambda}^{\Galg} W_{\lambda})\]
is an isomorphism. This proves that $A$ is small relative to $I_{\GG}$ \emph{in $\Galg$}. Exactly the same argument shows that similar conclusion holds for $J_{\GG}$. This completes the proof of the fact that monad $\GG$ acting on $\redTopSeq$ satisfies the hypothesis of Proposition \ref{modprop}. The conclusion of the corollary now follows from the Proposition.
\end{proof}




\section{Cofibrant resolution via the bar construction}\label{bar_const_red_sec}
In this section we described how to use the (geometric realization of) two sided bar construction to produce cofibrant resolutions of reduced properads. 

Recall that for an algebra $X \in \redTopSeq$ over the triple $(\GG,\mu,\eta)$, the two sided bar construction $B_\bullet(\GG,\GG,X)$ is a simplicial object in $\redTopSeq$ with $n$-simplices $B_n(\GG,\GG,X) = \GG^{n+1}X$ and with the usual face and degeneracy maps.

\subsection{Geomeric realizations}	\ \\
%%%%%%%%%%%%%%%%%%%%%%%%%%%%%%%%%%%%%%%%%%%%%%%%%%%%%%
%%%%%%%%%%%%%%%%                      Subsection              %%%%%%%%%%%%%%%%%%%%%%%
\subsubsection{Tensor and Cotensor over $\Top$}
The category of $\GG$-algebras in $\redTopSeq$ is tensored and cotensored over $\Top$. In other words, along with the bifunctor 
\[  [\ ,\ ] \colon \Galg \times \Galg^{op} \to \Top \]
providing a $\Top$ enrichment of $\Galg$, there are two additional bifunctors 
\begin{align*}
\odot &\colon \Galg \times \Top \to \Galg, \mbox{  and }\\ 
(\_)^{\_} &\colon \Galg \times \Top \to \Top  
\end{align*}
such that 
\[ Hom_{\Galg}(P \odot Z, Q) \simeq Hom_{\Top} ( Z, [P,Q]) \simeq Hom_{\Galg}(P,Q^Z) \]
naturally in all the variables involved. These bifunctors are constructed as follows:

The $\Top$-enrichment of $\Galg$ is given by endowing the $\Galg$ hom-spaces with the natural topology obtained by identifying them as subspaces of the $\redTopSeq$ hom-spaces. In other words, $[P,Q]$ is the space of $\GG$-algebra maps form $P$ to $Q$ with the product compact-open topology. Equivalently, if $[\ ,\ ]^{\redTopSeq}$ is the $\Top$-enrichment of $\redTopSeq$, $[ P,Q ]$ is defined as the equalizer
\[ [ P ,Q ]  \to [ P , Q ]^{\redTopSeq} \rightrightarrows [\GG P , Q ]^{\redTopSeq} \]
where the first arrow is induced by pullback along the $\GG$-algebra structure map $\GG P \to P$ and the second map is the composition $[P,Q]^{\redTopSeq} \to [\GG P, \GG Q]^{\redTopSeq} \to  [\GG P, Q]^{\redTopSeq}$

The operation $\odot$ can be constructed as 
\[ P \odot Z := coeq \Big[ \GG(\GG P \times Z) \rightrightarrows \GG(P \times Z) \Big] \]
Here, for $X \in \redTopSeq, Z \in \Top$, $X \times Z$ denotes the component-wise product $\{ X(n_-,n_+) \times Z\}_{n_-,n_+} \in \redTopSeq$. The coequalizer can be taken in either $\Galg$ or $\redTopSeq$, since they coincide in this case. The first arrow is induced from the $\GG$-algebra structure map $\GG P \to P$, whereas the second arrow is induced from $\GG P \times Z \to \GG( P \times Z)$ using the universal property of free $\GG$-algebras.

The operation $(\_)^{\_}$ can be constructed as follows: The underlying sequence of $P^Z$ is given by $P^Z=\{Hom_{\Top}(Z,P(n_-,n_+))\}_{n_-,n_+}$ and the $\GG$-algebra structure is induced from that of $P$.\\

These operations are compatible with the model structure in the following sense:
\begin{lemma}[\cite{rezk1996spaces}]\label{complem}
	The following equivalent conditions are satisfied:
	\begin{enumerate}
		\item \label{lem1} If $i \colon Y \to Z$ is a cofibration in $\Top$ and $j \colon P \to Q$ is a cofibration in $\Galg$, then $P \odot Z \bigsqcup_{P \odot Y} Q \odot Y \to Q \odot Z$ is a cofibration of $\GG$-algebra. Moreover this map is a weak equivalence whenever $i$ or $j$ is.
		\item  \label{lem2} If $i \colon Y \to Z$ is a cofibration in $\Top$ and $j \colon P \to Q$ is a fibration in $\Galg$, then $P^Z \to P^Y \times_{Q^Y} Q^Z$ is a fibration in $\Galg$. Moreover this map is a weak equivalence whenever $i$ or $j$ is.
		\item If $j \colon P \to Q$ is a cofibration and $k \colon R \to S$ a fibration in $\GG$-algebra, then $[Q,Z] 	\to [P,S] \times_{[Q,S]} [Q,R]$ is a fibration. Again, this map is a weak equivalence whenever $j$ or $k$ is.
	\end{enumerate}
\end{lemma}
	\begin{proof}
	We explain the proof of (\ref{lem2}) and the implication (\ref{lem2}) $\Rightarrow$ (\ref{lem1}), since this is the only part we shall use.\\
	(\ref{lem2}): Note that we need to show $P^Z \to P^Y \times_{Q^Y} Q^Z$ is a fibration of underlying topological sequences. The fibered product $P^Y \times_{Q^Y} Q^Z$ coincides with the fibered product of the underlying sequences. Let $V \to W$ be an acyclic cofibration in $\redTopSeq$. We need to show that any square
\[	\begin{tikzcd}
		V \ar{r}\ar{d} & P^Z \ar{d}\\
		W \ar{r}& P^Y \times_{Q^Y} Q^Z
	\end{tikzcd}
\]
	admits a lift $W \to P^Z$ (in $\redTopSeq$). This is equivalent to the statement that any square
\[
	\begin{tikzcd}
		W \times Y \sqcup_{V \times Y} V \times Z \ar{r}\ar{d} &P \ar{d} \\
		W \times Z \ar{r} &Q
	\end{tikzcd}
\]
	admits a lift $W \times Z \to P$ (in $\redTopSeq$). But this follows from the fact that
	\begin{itemize}
		\item $P \to Q$ is a fibration in $\Galg$, and hence in $\redTopSeq$
		\item $X \to Y$ is cofibration in $\Top$ and $V \to W$ is a acyclic cofibration in $\redTopSeq$, and hence $W \times Y \sqcup_{V \times Y} V \times Z \to W \times Z$ is an acylic cofibration in $\redTopSeq$.
	\end{itemize}
		This completes the proof of (\ref{lem2}).
		
	(\ref{lem2}) $\Rightarrow$ (\ref{lem1}): To show that $Q \odot Y \sqcup_{P \odot Y} P \odot Z \to Q \odot Z$ is a cofibration in $\GG$-algebra, we show that for any acyclic fibration $E \to F$ in $\Galg$ the square
\[
	\begin{tikzcd}
			Q \odot Y \sqcup_{P \odot Y}  P \odot Z \ar{r}\ar{d} & E \ar{d}\\
			Q \odot Z \ar{r} & F
	\end{tikzcd}
\]
admits a lift $Q \odot Z \to E$ (in $\Galg$). This is equivalent to 
\[
	\begin{tikzcd}
	P \ar{r}\ar{d} & E^Z \ar{d}\\
	Q \ar{r} &F^Y \times_{F^Y} F^Z
	\end{tikzcd}
\]
admitting a lift $Q \to E^Z$ (in $\Galg$). But this is a consequence of (\ref{lem2}) combined with the fact that $P \to Q$ is a $\Galg$ cofibration.
	\end{proof}
	
%%%%%%%%%%%%%%%%%%%%%%%%%%%%%%%%%%%%%%%%%%%%%%%%%%%%%%
%%%%%%%%%%%%%%%%                      Subsection              %%%%%%%%%%%%%%%%%%%%%%%

\subsubsection{the Geometric realization}
Now, using this tensor product over $\Top$ on $\Galg$, we can define the geometric realization of a simplcial $\GG$-algebra $\{P_\bullet\}_{\bullet \geq 0}$ by the usual formula:
\begin{equation}\label{geomrelform}
|P_\bullet|^{\Galg} = \int_{\Delta}^{\Galg} P_n \odot \Delta^n := coeq^{\Galg} \Big[ \coprod_{\phi \colon [n] \to [m] \in \Delta}^{\Galg} P_m \odot \Delta^n\rightrightarrows \coprod_{[n] \in \Delta}^{\Galg} P_n \odot \Delta^n \Big]
\end{equation}
Here $\Delta$ denotes the simplex category, and $\Delta^n \in \Top$ is the standard $n$-simplex.	In the coequalizer, the first arrow is induced from the simplcial structure map $P^*(\phi) \colon P_m \to P_n \in \Galg$ corresponding to $\phi \colon [n] \to [m]$ and the second arrow is induced from the map $\Delta(\phi) \colon \Delta^n \to \Delta^m \in \Top$.

There is an alternate description of $|\ \ |^{\Galg}$ as an iterated pushout in terms of the `latching spaces': Define the $n$th \textbf{latching space} of $P_\bullet$ to be 
\[ L_nP_\bullet := coeq^{\Galg} \Big[ \coprod_{0 \leq i < j \leq n}^{\Galg} P_{n-1} \rightrightarrows \coprod_{0 \leq i \leq n}^{\Galg} P_n \Big] \]
The two arrows are defined as follows: on the $(i,j)$-th summand, the first arrow is induced by the degeneracy map $s_i \colon P_{n-1} \to P_n$ into the $j$th summand and the second arrow is induced from $s_{j-1} \colon P_{n-1} \to P_n$ into the $i$th summand. Note that there is a canonical map
\[ L_{n-1}P_\bullet \to P_{n+1} \]
$L_n P_\bullet$ can be thought of intuitively as ``the space of degenerate simplices in $P_{n+1}$''
Define $|\ \ |^{\Galg}_{(n)}$ inductively as:
	\begin{itemize}
	\item $|P_\bullet|^{\Galg}_{(0)} =  P_0$, and
	\item $|P_\bullet|^{\Galg}_{(n+1)}$ is defined to be the pushout
	\[
		\begin{tikzcd}\label{latchingdig}
			L_nP_\bullet \odot \Delta^{n+1} \bigsqcup_{L_nP_\bullet \odot \partial \Delta^{n+1}} P_{n+1} \odot \partial \Delta^{n+1} \ar{r}\ar{d} & P_{n+1} \odot \Delta^{n+1} \ar{d}\\
			{|P_\bullet|}^{\Galg}_{(n)} \ar{r} & {|P_\bullet|}^{\Galg}_{(n+1)}\\
		\end{tikzcd}
	\]
	\end{itemize}

Then $|P_\bullet|^{\Galg} = \varinjlim_n |P_\bullet|^{\Galg}_{(n)}$

\begin{lemma}
	If $L_nP_\bullet \to P_{n+1}$ is a cofibration, then $|P_\bullet|^{\Galg}$ is cofibrant in $\Galg$.
\end{lemma}
\begin{proof}
	It suffice to prove that the upper horizontal arrow in the diagram \ref{latchingdig} is a cofibration for all $n$. This follows from the compatibility statements from Lemma	 \ref{complem} since $\partial \Delta^{n+1} \hookrightarrow \Delta^{n+1} \in cofib(\Top)$ and $P \to Q \in cofib(\Galg)$.
\end{proof}

%%%%%%%%%%%%%%%%%%%%%%%%%%%%%%%%%%%%%%%%%%%%%%%%%%%%%%
%%%%%%%%%%%%%%%%                      Subsection              %%%%%%%%%%%%%%%%%%%%%%%

\subsubsection{Cofibrancy of the bar construction}\label{Cofibrancy of the bar construction}
We now show that the geometric realization of the simplicial $\GG$-algebra $B_\bullet(\GG,\GG,P)$ is cofibrant provided the topological sequence underlying the $\GG$-algebra $P$ is cofibrant.

The latching object for the bar construction are given by 
\[ L_n B(\GG,\GG,P) = coeq^{\Galg} \Big[ \coprod_{0 \leq i < j \leq n}^{\Galg}\GG^{n}P \rightrightarrows \coprod_{0 \leq i \leq n}^{\Galg} \GG^{n+1}P \Big] \]
The free $\GG$-algebra functor, being the left adjoint in the pair $\GG \colon \redTopSeq \leftrightarrows \Galg \colon$ Forget, commutes with coproudcts and coequalizer. Thus,
\[ L_nB(\GG,\GG,P) = \GG \Big( coeq^{\redTopSeq} \Big[ \coprod_{0 \leq i < j \leq n}^{\redTopSeq}\GG^{n-1}P \rightrightarrows \coprod_{0 \leq i \leq n}^{\redTopSeq} \GG^{n}P \Big] \Big) \]
Set 
\[ K_nP := coeq^{\redTopSeq} \Big[ \coprod_{0 \leq i < j \leq n}^{\redTopSeq}\GG^{n-1}P \rightrightarrows \coprod_{0 \leq i \leq n}^{\redTopSeq} \GG^{n}P \Big] \]
Then, to show that $L_nB(\GG,\GG,X) \to B_{n+1}(\GG,\GG,P) = \GG^{n+2}P \in cofib(\Galg)$ is a cofibration, it suffices to show that $K_nP \to \GG^{n+1}P \in cofib(\redTopSeq)$. But, it is not difficult to see that this map identifies $K_nP$ with the subspace of $\GG^{n+1}P$ given by $\bigcup_{0 \leq i \leq n} s_i(\GG^n P)$, where $s_i$ are the simplicial degeneracy maps.\\
Moreover, note that $\GG^{n+1}P$ is a disjoint union of spaces indexed by $(n+1)$-nested graphs and $s_i(\GG^nP)$ is the union corresponding to a subset of this indexing set. Since the topological sequence underlying $P$ is cofibrant, it follows that the inclusion $\bigcup_{0 \leq i \leq n}s_i(\GG^n P) \hookrightarrow \GG^{n+1}P$ a cofibration of topological sequences. Thus we have:
\begin{proposition}
If $P_\bullet$ is cofibrant as an object of $\redTopSeq$, then $|B(\GG,\GG,P)|^{\Galg}$ is cofibrant in $\Galg$.
\end{proposition}

\subsubsection{Comparison of $|P_\bullet|^{\Galg}$ and $|P_\bullet|^{\redTopSeq}$ for a simplicial $\GG$-algebra}
In this section we follow the discussion in \cite[Section 7]{mandell2019operads}.\\
Let $P_\bullet$ be a simplicial $\GG$-algebra. Via the forgetful functor, we may interpret $P_\bullet$ as a simplicial object in $\redTopSeq$ and thus consider the space $|P_\bullet|^{\redTopSeq}$, where $|\ \ |^{\redTopSeq}$ is the geometric realization functor on the category of simplicial objects in $\redTopSeq$. It is defined by a formula analogous to \ref{geomrelform}, with the coproduct and the coequalizer taken in $\redTopSeq$ instead of $\Galg$.


It is a well known that $|\ \ |^{\redTopSeq}$ is a monoidal functor (i.e. satisfies $|X_\bullet \times Y_\bullet| = |X_\bullet | \times |Y_\bullet |$), and hence it follows that $|P_\bullet|^{\redTopSeq}$ has a natural $\GG$-algebra structure:\\
Given $G \in \mathcal{G}$, we have $\prod_{v\in G} |P_\bullet|(in(v), out(v)) \to |P_\bullet|(in(G), out(G))$, induced by 
\[ 
	\begin{tikzcd}
	\prod_{v\in G} |P_\bullet|(in(v), out(v)) \ar[dashrightarrow]{r} \ar{d}{\simeq} &  {|P_\bullet|}(in(G), out(G)) \ar{d}{\simeq}  \\
	{|\prod_{v\in G} P_\bullet(in(v), out(v))|} \ar{r} & {|P_\bullet (in(G),out(G))| }\\
	\end{tikzcd}
\]
A priori it is not clear what the relation between $\GG$-algebras $|P_\bullet|^{\Galg}$ and $|P_\bullet|^{\redTopSeq}$ is. However, we have the following (somewhat surprising) fact:

\begin{proposition}[\cite{mandell2019operads}]
	$|P_\bullet|^{\Galg} \simeq |P_\bullet|^{\redTopSeq}$
\end{proposition}
\begin{proof}
	Recall that $P_n \odot \Delta^n = coeq \Big[ \GG(\GG P_n \times \Delta^n) \rightrightarrows \GG(P_n \times \Delta^n) \Big]$ (with the coequalizer in $\Galg$ and $\redTopSeq$ coinciding). The maps $P_n \times \Delta^n \to |P_\bullet|^{\redTopSeq}$ induce $\GG$-algebra maps $P_n \odot \Delta^n \to |P_\bullet|^{\redTopSeq}$. Moreover, these maps are compatible with the simplicial face and degeneracy maps, and hence we get a $\GG$-algebra map
\[ coeq \Big[ \coprod_{m,n}^{\Galg} \rightrightarrows \coprod_{n}^{\Galg} P_n \odot \Delta^n \Big] = |P_\bullet|^{\Galg} \to |P_\bullet|^{\redTopSeq} \]
Now, notice that for $P \in \Galg$, $|\GG P_\bullet|^{\redTopSeq} \simeq \GG|P_\bullet|^{\redTopSeq}$. This holds since $|\ \  |^{\redTopSeq}$, being a left adjoint ($|\ \ |^{\redTopSeq}\colon \redTopSeq \leftrightarrows \redTopSeq \colon (\_)^{\Delta^\bullet}$), commutes with colimits and as noted previously above it also commutes with finite products.

Furthermore, for any $X \in \redTopSeq$, $|\GG X_\bullet|^{\Galg} \simeq \GG |X_\bullet|^{\redTopSeq}$: Firstly, note that for $X \in \redTopSeq$,
\begin{align*}
 \GG X_n \odot \Delta^n = coeq \Big[ \GG (\GG(\GG X_n) \times \Delta^n ) \rightrightarrows \GG(\GG X_n \times \Delta^n) \Big] \simeq \GG(X_n \times \Delta^n) 
\end{align*}
\begin{align*}
\implies \coprod_{m,n}^{\Galg} \GG X_n \odot \Delta^m &\simeq \GG \Big( \coprod_{m,n}^{\redTopSeq} X_n \times \Delta^m \Big) \mbox{ and } \\
 \coprod_{n}^{\Galg} \GG X_n \odot \Delta^n &\simeq \GG \Big( \coprod_{n}^{\redTopSeq} X_n \times \Delta^m \Big) 
\end{align*}
Since $\GG$ commutes with colimits, this gives us
\begin{equation*}\begin{split}
 coeq^{\Galg} \Big[ \coprod^{\redTopSeq} & X_n \odot \Delta^m \rightrightarrows  \coprod \GG X_n \odot \Delta^n\Big] \simeq\\
 & \GG \Big( coeq^{\redTopSeq} \Big[ \coprod^{\redTopSeq} X_n \times \Delta^m \rightrightarrows  \coprod X_n \times \Delta^n \Big] \Big) 
\end{split}
\end{equation*}

Thus, it follows that for $X_\bullet \in \redTopSeq$, $|\GG X_\bullet|^{\Galg} \simeq \GG |X_\bullet|^{\redTopSeq}$
Finally, for any $P_\bullet \in s\Galg$
\[ \GG(\GG P ) \rightrightarrows \GG P \to P\]
is a coequalizer diagram (in $\Galg$ as well as in $\redTopSeq$). Since $|\ \ |^{\Galg}$ is a left adjoint ($|\ \ |^{\Galg} \colon s\Galg \leftrightarrows \Galg \colon (\_)^{\Delta^\bullet}$), it commutes with the coequalizer. Thus we have the following commutative diagram in $\Galg$:
\[
	\begin{tikzcd}
		{|\GG(\GG P) |}^{\Galg} \ar[shift left]{r}\ar[shift right]{r} \ar{d} & {|\GG P |}^{\Galg} \ar{r}\ar{d} & {|P_\bullet|}^{\Galg} \ar[dashed]{d}\\
		\GG \GG {|P_\bullet|}^{\redTopSeq}\ar[shift left]{r}\ar[shift right]{r} & {\GG|P_\bullet|}^{\redTopSeq} \ar{r} &  {|P_\bullet|}^{\redTopSeq}
	\end{tikzcd}
\]
where both the rows are coequalizer diagrams mentioned  above. Further, the previous discussion shows that the first two vertical maps are isomorphism. It thus follows that the rightmost vertical arrow is also an isomorphism, proving the proposition.
\end{proof}

\begin{corollary}
	For $P_\bullet \in \Galg$, cofibrant as a topological sequence, $|B_\bullet(\GG,\GG,P)|^{\Top} \to P$ is a cofibrant resolution of $P$.
\end{corollary}
\begin{proof}
	By the above discussion it follows that $|B_\bullet(\GG,\GG,P)|^{\Top}  \simeq |B_\bullet(\GG,\GG,P)|^{\Galg} $ and hence is cofibrant as a $\GG$-algebra. By the standard argument (see \cite{} ) it follows that the constant simplicial object $P$ is a simplicial deformation retract of $B_\bullet(\GG,\GG,P)$ in $\redTopSeq$, hence taking the geometric realization we get that $|B_\bullet(\GG,\GG,P)|^{\Top} \to P$ is a weak equivalence.
\end{proof}


\section{Properads of Riemann surface moduli}
For $g,n_-,n_+$ in the stable range $2-2g-n_--n_+<0$, let $\Mmod_{g,n_-,n_+}$ denote the moduli space of genus $g$ riemann surface with $n_-$ input and $n_+$ output marked points. Let $\overline{\Mmod}_{g,n_-,n_+}$ be its Deligne-Mumford compactification, given by the moduli spaces of stable nodal curves with $n_-$ input and $n_+$ output marked points, and arithmetic genus $g$.

Let $\Mmod^{fr}_{g,n_-,n_+}$ denote the moduli space of genus $g$ riemann surface with $n_-$ input and $n_+$ output boundary components, each equipped with an analytic $S^1$-parametrization. $\Mmodbar^{fr}_{g,n_-,n_+}$ is the moduli space of stable nodal riemann surfaces with $n_-$ input and $n_+$ output boundaries which are non-nodal and analytically $S^1$-parameterized. Also, let $\Mmodhat^{fr}_{g,n_-,n_+}$ be subspace of $\Mmodbar^{fr}_{g,n_-,n_+}$ given by the moduli of stable curves in which every irreducible component contains at least one output.

In addition, we also consider the moduli spaces $\Mmod^{fr}_{g,n_-,n_+},\Mmodbar^{fr}_{g,n_-,n_+}, \Mmodhat^{fr}_{g,n_-,n_+}$ for the unstable range $(g,n_-,n_+) \in \{(0,1,1), (0,0,2), (0,2,0)\}$. In these cases, the spaces $\Mmod^{fr}_{g,n_-,n_+}$ are all identified with the space of annuli with $S^1$-parametrized boundary components. Spaces $\Mmodbar^{fr}_{g,n_-,n_+}$ and $\Mmodhat^{fr}_{0,0,2}$ are obtained by compactifying these spaces  by also including nodal annuli, thought of as annuli with modulus $\infty$(with boundaries suitably labeled as inputs and outputs). Consistent with the restriction that each irreducible component has an output, the space $\Mmodhat_{0,1,1}$ coincides with $\Mmod_{0,1,1}$ and $\Mmodhat^{fr}_{0,2,0}$ is empty.
%For $\Mmodhat^{fr}_{g,n_-,n_+}$, consistent with the restriction that each irreducible component contains at least one output, $\Mmodhat^{fr}_{0,2,0}$ is empty and $\Mmodhat^{fr}_{0,1,1}$ does not contain the nodal annulus. Where as the other hand $\Mmodhat^{fr}_{0,0,2}$ coincides with $\Mmodbar^{fr}_{0,0,2}$ and in particular contains the nodal annulus.
%In the case of $\Mmodbar^{fr}_{g,n_-,n_+}, \Mmodhat^{fr}_{g,n_-,n_+}$ these in particular include nodal annuli, thought of as annuli with modulus $\infty$.
Moreover, in the case $(g,n_-,n_+)=(0,1,1)$ for all the three types of moduli mentioned above, we also allow exceptional points corresponding to degenerate annuli with modulus $0$. These points are added to ensure that the reduced properads we consider below are untial.

Finally, denote by $\Mmod^{fr}_{g,n_-,n_+}(i)$ and $\Mmodbar^{fr}_{g,n_-,n_+}(i)$ the moduli of surfaces in $\Mmod^{fr}_{g,n_-,n_+}$ and $\Mmodbar^{fr}_{g,n_-,n_+}$ respectively, with $i$ analytically embedded circles disjoint from the nodes and the boundaries.

%\Mfrfull_{g,n_-,n_+}$, where $\Mfrfull_{g,n_-,n_+} = 
% \Mmodbar_{g,n_-,n_+}$, where $\Mfrfull_{g,n_-,n_+} =
% \Mhat_{g,n_-,n_+}$, where $\Mfrfull_{g,n_-,n_+} = 

We now describe the reduced properads which will appear in our discussion:
	\begin{itemize}
	\item $\Mfrfull$ be the properad defined by $\Mfrfull (n_-,n_+) = \coprod_{g \geq 0} \Mmod^{fr}_{g,n_-,n_+}$. The properad compositions are given by gluing the riemann surfaces along suitable boundaries via the $S^1$-parametrization.
	\item $\Mbar$ be the properad defined by $\Mbar (n_-,n_+) = \coprod_{g \geq 0} \Mmodbar^{fr}_{g,n_-,n_+}$. The properad compositions are again given by gluing the nodal riemann surfaces along suitable boundaries via the $S^1$-parametrization, followed by \emph{collapsing unstable components}
	\item $\Mfr$ be   the subproperad of $\Mfrfull$ consisting of all operations which have at least one output i.e. $\Mfr (n_-,n_+) = \coprod_{g \geq 0} \Mfr_{g,n_-,n_+}$, where
\[\Mfr_{g,n_-,n_+} = \begin{cases} \varnothing &\mbox{ if } n_+=0 \\ \Mmod^{fr}_{g,n_-,n_+} &\mbox{ otherwise }\end{cases} \]
	\item $\Mhat$ be the subproperad of $\Mbar$ defined by $\Mhat (n_-,n_+) = \coprod_{g \geq 0}\Mmodhat^{fr}_{g,n_-,n_+}$. Note that similar to $\Mfr$, this is empty in the components with $n_+=0$
	\item $\Mfrfullunst$ be the subproperad of  $\Mfrfull$ consisting of the moduli in the unstable range i.e. it is empty in all components except the following:
\[
				\Mfrfullunst(1,1) = \Mmod^{fr}_{0,1,1}, \Mfrfullunst(0,2) = \Mmod^{fr}_{0,0,2}, \Mfrfullunst(2,0) = \Mmod^{fr}_{0,2,0} %\Mfrfullunst(0,0) = \Mmod^{fr}_{0,0,0}  \sqcup \Mmod^{fr}_{1,0,0}
\]
	\item $\Mbarunst$ be the subproperad of  $\Mbar$ consisting of the moduli in the unstable range i.e. it is empty in all components except the following:
\[
				\Mbarunst(1,1)=\Mmodbar^{fr}_{0,1,1}, \Mbarunst(0,2) = \Mmodbar^{fr}_{0,0,2}, \Mbarunst(2,0) = \Mmodbar^{fr}_{0,2,0} %\Mfrfullunst(0,0) = \
\]
	The properad operation is defined as in $\Mbar$, i.e. by gluing along boundaries followed by stablilization. $\Mfrunst$ is a subproperad of $\Mbar$.
	\item $\Mhatunst$ be the subproperad of $\Mhat$ consisting of the moduli in the unstable range i.e. it is empty in all components except the following:
\[
				\Mhatunst(1,1)=\Mmodhat^{fr}_{0,1,1} , 	  \Mhatunst(0,2)=\Mmodhat^{fr}_{0,0,2} 
\]
	\item Finally $\Mfrunst$ be the subproperad of $\Mfr$ consisting of the moduli in the unstable range. It is empty in all components except the following
\[
	  \Mfrunst(1,1)=\Mmod^{fr}_{0,1,1} , 	  \Mfrunst(0,2)=\Mmod^{fr}_{0,0,2} 
\]
Clearly $\Mfrunst$ is a subproperad of $\Mfr, \Mhatunst,$ and $\Mbarunst$.
	\end{itemize}
	In the discussion in Section \ref{} below, we will also use the following modifications of properads $\Mfr, \Mfrunst, \Mhatunst$ which we record here for convenience:
	\begin{itemize}
\item $\mmfrfull$ is a subproperad of $\Mfr$ which coincides with $\Mfr$ in all degrees except $(1,1), (0,2)$. In these degrees the genus $0$ components of $\mfrfull(1,1),\mfrfull(0,2)$ are the subspaces of $\mmfrfull(1,1), \mmfrfull(0,2)$ respectively, containing only the exceptional annuli which have modulus $0$. 
\item $\Mfrunstth$ is defined as a subproperad of $\Mfrunst$ in an analogous manner.
\item $\Mhatunstth$ is a subproperad of $\Mhatunst$. In the degree $(1,1)$, similar to the above cases, it only contains the exceptional annuli of modulus $0$. However in this case, we define $\Mhatunstth(0,2)$ to coincide with $\Mhatunst(0,2)$.
\end{itemize}
	
\section{Main theorem}
\begin{theorem}\label{mthm1}
	$\Mhat$ is the homotopy pushout of 
	\begin{equation}\label{thm1dgm}
		\begin{tikzcd}
		\Mfrunst \ar{r}\ar{d} & \Mhatunst\\
		\Mfr & 
		\end{tikzcd}
	\end{equation}
	in $\Galg$.
\end{theorem}

Rest of the Section occupies proof of the Theorem \ref{mthm1}
Below, whenever we use $coeq,\sqcup,$ and $\coprod$ without a superscript $\Galg$, it is understood that they refer to the respective operations in $\redTopSeq$.

\subsection{Pushouts in $\Galg$}
We start by giving a description of pushouts in the category of $\GG$-algebras.

\begin{lemma}\label{pushoutlem}
	$P \leftarrow R \to Q$ be a diagram in $\Galg$. Then the pushout $P \sqcup_{R}^{\Galg} Q$ can be obtained as the following coequalizer in $\redTopSeq$
	\[ coeq\Big[ \GG(\GG P \coprod_{\GG R} \GG Q) \rightrightarrows \GG(P \coprod_{R} Q) \Big] \]
	with the $\GG$-algebra structure induced from that of $\GG(P \sqcup_{R} Q)$. Here 
	\begin{enumerate}
	\item the first arrow is given by applying functor $\GG$ to the map $ \GG P \sqcup_{\GG R} \GG Q \to P \sqcup_{R} Q$ which itself is induced from the $\GG$-algebra structure maps $\GG P \to P, \GG Q \to Q$, and $\GG R \to R$
	\item the second arrow is induced via the universal property of free $\GG$-algebras using the map $ \GG P \sqcup_{\GG R} \GG Q \to \GG(P \sqcup_{R} Q) $ which in turn is the pushout of the maps given by applying functor $\GG$ to $P \to P \sqcup_R Q$, $Q \to P \sqcup_R Q$, and $R \to P \sqcup_R Q$
	\end{enumerate}
\end{lemma}
The Lemma is proved by directly verifying that the coequalizer with the given $\GG$-algebra structure satisfies the universal property of the pushout.

\begin{notation} 
	We say that a subgraph $H$ of a rda-graph $G$ is \textbf{collapsible} if the graph $G/H$, obtained by replacing $H$ with a single vertex, contains no directed cycles.\\	
	More precisely: consider the topological sequence $\mathcal{G}= \{\mathcal{G}_{n_-,n_+}\}$ (with the discrete topology), where, as in Section \ref{GGtripsubsec}, $\mathcal{G}$ is the collection containing one representative from each isomorphism class of rda-graphs . It has a natural $\GG$-algebra structure. We say that $H \subset G$ is collapsible if $\mu_K(\prod_{v \in V(K)} H_v ) = G$ for some rda-graphs $K$ and $H_{v}, v \in V(K)$, such that $H_v$, considered as a subgraph of $G$, coincides with $H$ for some $v$.
\end{notation}

Unraveling the coequalizer in Lemma \ref{pushoutlem}, we get that the $\redTopSeq$ underlying $P \sqcup_{R}^{\Galg} Q$ is given by 
\[ \frac{\GG(P \sqcup Q)}{\sim} \]
where $\sim$ is the equivalence relation generated by the following identifications:
\begin{enumerate}
	\item Let $G$ be a $(P \sqcup Q)$-labeled graph which has a collapsible subgraph $H \subset G$ with all the vertices labeled by $P$ (resp. $Q$). Then $G$ is identified with the graph obtained by collapsing $H$ to a single vertex labeled by the element of $P$(resp. $Q$) obtained by applying $\GG$-algebra composition (indexed by $H$) to the labels of $H$. 
	\item Suppose that the maps $R\to P$ and $R \to Q$ are denoted by $i$ and $j$ respectively. Let $G$ be a $(P \sqcup Q)$-labeled graph with a vertex $v \in G$ labeled by an element $i(r) \in i(R) \subset P$ for some $r \in R$. Then $G$ is identified with the graph obtained by switching the label of $v$ to $j(r) \in j(R) \subset Q$.
\end{enumerate}

\subsection{The homotopy pushout}
In this section we show that the pushout 
\[|B_\bullet \Mfr| \sqcup^{\Galg}_{|B_\bullet \Mfrunst|} |B_\bullet \Mhatunst|\] gives the homotopy pushout of the diagram \eqref{thm1dgm}.

Consider the inclusion 
\begin{equation}\label{inclusionofdiagrams} 
	\begin{tikzcd}
	\mmfrfull \ar{d} & \Mfrunstth\ar{l}\ar{d}\ar{r} & \Mhatunstth \ar{d} \\ 
	\Mfr & \Mfrunst \ar{l}\ar{r}& \Mhatunst
	\end{tikzcd}
\end{equation} 
of pushout diagrams. Note that this is a weak-equivalence of pushout diagrams.\\
 Moreover, $\mmfrfull,\Mfrunstth, \Mhatunstth$ \textcolor{red}{are $\redTopSeq$-cofibrant} and the maps $\Mfrunstth \to \mmfrfull, \Mfrunstth \to \Mhatunstth$ are $\redTopSeq$-cofibrant in the standard model structure on $\redTopSeq$.
Therefore the corresponding bar constructions 
\[|B_\bullet \mmfrfull|,|B_\bullet \Mfrunstth|, \mbox{ and } |B_\bullet \Mhatunstth|\]
 are $\Galg$-cofibrant and the maps 
 \[|B_\bullet \Mfrunstth| \to |B_\bullet \mmfrfull| \mbox{ and } |B_\bullet \Mfrunstth| \to |B_\bullet \Mhatunstth|\]
  are $\Galg$-cofibrations.\\ 
It then follows that the pushout
\[|B_\bullet \mmfrfull| \sqcup^{\Galg}_{|B_\bullet \Mfrunstth|} |B_\bullet \Mhatunstth|\]
computes the homotopy pushout of \eqref{thm1dgm}. Analogous conclusion for the pushout 
\[ |B_\bullet \Mfr| \sqcup^{\Galg}_{|B_\bullet \Mfrunst|} |B_\bullet \Mhatunst|\]
 is not clear from the argument above, since the map $\Mfrunst \to \Mhatunst$ is not a $\redTopSeq$-cofibration. The proposition \ref{compareproposition} below allows us compare these two pushouts.

%We start with an easy preparatory Lemma:
%\begin{lemma}
%	$\GG$ preserves weak equivalence between cofibrant objects.
%\end{lemma}
%Proof follows just from 
%	\begin{proof}
%	This is a consequence of the \emph{Ken Brown's lemma} using the fact that the pair $\GG \colon \redTopSeq \leftrightarrows \Galg \colon Forget$ is a \emph{Quillen adjunction} (cf. \cite{}). More concretely, in the category $\redTopSeq$ this can be seen as follows: \\
%	Let $f \colon X \to Y$ be a weak equivalence between cofibrant $\redTopSeq$-objects. It can factored as $X \xrightarrow{i} M(f) \xrightarrow{r} Y$, where $ M(f)$ is the mapping cylinder of $f$, $i$ is an acyclic cofibration, and $r$ is a retraction such that $r \circ j =1_Y$ for an acyclic cofibration $j \colon X \to Y$. From the adjunction $\GG \colon \redTopSeq \leftrightarrows  \Galg \colon Forget$, it follows that $\GG$ maps $\redTopSeq$-acyclic cofibrations to $\Galg$-acyclic cofibrations. It then follows 
that $\GG(i)$ and $\GG(j)$ are weak equivalences. The latter implies that $\GG(r)$ is a weak equivalence which, combined with the former, in turn implies that $\GG(f)$ is a  weak equivalence.\\
%	\end{proof}

\begin{proposition}\label{compareproposition}
Let 
\begin{equation}\label{comparepropositiondig}
	\begin{tikzcd}
	P \ar{d} & R \ar{l}\ar{r}\ar{d} & Q \ar{d} \\
	P^{\prime}  & R^{\prime} \ar{l}\ar{r} & Q^{\prime} 
	\end{tikzcd}
\end{equation}
be a weak-equivalence of $\Galg$-diagrams.
%be an acyclic cofibration of $\redTopSeq$-diagrams. In other words, suppose that 
%	\begin{enumerate}
%		\item All the vertical maps are weak equivalences
%		\item $R \to R^{\prime}, P \sqcup_{R} R^{\prime} \to P^{\prime}$, and $Q \sqcup_{R} R^{\prime} \to Q^{\prime}$ are $\redTopSeq$-cofibrations.
%	\end{enumerate}
Assume that 
	\begin{enumerate}
		\item $P,Q,R$ and $P^{\prime}, Q^{\prime} , R^{\prime}$ are all $\redTopSeq$-cofibrant
		\item $R \to P$ , $R^{\prime} \to P^{\prime}$ are $\redTopSeq$-cofibrations
	\end{enumerate}
	Then 
	\[ |B_\bullet P| \coprod^{\Galg}_{|B_\bullet R|} |B_\bullet Q| \to |B_\bullet P^{\prime}| \coprod^{\Galg}_{|B_\bullet R^{\prime}|}|B_\bullet P^{\prime}|\]
	is a weak-equivalence.
\end{proposition}
	
	\begin{proof}
		Now,
	\begin{align}
	\label{barpushout}	|B_\bullet P| \coprod_{|B_\bullet R|}^{\Galg}  |B_\bullet Q|  & \simeq \Big|B_\bullet P \coprod_{B_\bullet R}^{\Galg} B_\bullet Q\Big| \simeq  \Big|\GG^{\bullet+1}  P  \coprod_{\GG^{\bullet+1}  R}^{\Galg} \GG^{\bullet+1} Q \Big| \\
	\notag & \simeq  \Big| \GG \big(\GG^{\bullet}  P \coprod_{\GG^{\bullet}  R} \GG^{\bullet} Q \big) \Big|
	\end{align}
	Here the first equality holds since $|\ \ |$ is a left adjoint in the pair $| \ \ | \colon \leftrightarrows \Galg \colon (\_)^{\Delta^\bullet}$ and hence commutes with colimits. Similarly the last equality holds since $\GG$ is a left adjoint in the pair $\GG \colon \redTopSeq \leftrightarrows  \Galg \colon Forget$ and hence commutes with colimits.
	Similarly, 
	\[ |B_\bullet P^{\prime}| \coprod_{|B_\bullet P^{\prime}|}^{\Galg} |B_\bullet Q^{\prime}|  \simeq  \Big| \GG \big(\GG^{\bullet}  P^{\prime} \coprod_{\GG^{\bullet}  R^{\prime}} \GG^{\bullet} Q^{\prime} \big) \Big| \]
	The discussion in Section \ref{bar_const_red_sec} implies the objects $B_\bullet P$, $B_\bullet Q$, and $B_\bullet Q$ are Reedy cofibrant, and the map $B_\bullet R \to B_\bullet P$ is a Reedy cofibration in $s\Galg$. Thus, $B_\bullet P \sqcup^{\Galg}_{B_\bullet R} B_\bullet Q $ is Reedy cofibrant in $s\Galg$. Similarly, $B_\bullet P^{\prime} \sqcup^{\Galg}_{B_\bullet R^{\prime}}B_\bullet P^{\prime}$ is Reedy cofibrant. Thus from \cite[Corollary 10.6]{riehl2013theory} it suffice to show that the map
	\[   \GG \big(\GG^{\bullet}  P \coprod_{\GG^{\bullet}  R} \GG^{\bullet} Q \big) \to  \GG \big(\GG^{\bullet}  P^{\prime} \coprod_{\GG^{\bullet}  R^{\prime}} \GG^{\bullet} Q^{\prime} \big)  \]
	is a Reedy weak-equivalence i.e. a levelwise weak-equivalence.\\
	Applying the functor $\GG^{\bullet}$ to the diagram \eqref{comparepropositiondig}, we get an inclusion of $\redTopSeq$  diagrams
	\begin{equation}\label{comparepropositiondig}
		\begin{tikzcd}
		\GG^{\bullet}P \ar{d} & \GG^{\bullet}R \ar{l}\ar{r}\ar{d} & \GG^{\bullet}Q \ar{d} \\
		\GG^{\bullet}P^{\prime}  & \GG^{\bullet}R^{\prime} \ar{l}\ar{r} & \GG^{\bullet}Q^{\prime} 
		\end{tikzcd}
	\end{equation}
	From the definition of the functor $\GG^{\bullet}$ it follows that 
	\begin{itemize}
		\item $\GG^{\bullet}P,\GG^{\bullet}Q,\GG^{\bullet}R$ and $\GG^{\bullet}P^{\prime}, \GG^{\bullet}Q^{\prime} ,\GG^{\bullet} R^{\prime}$ are all $\redTopSeq$-cofibrant
		\item $\GG^{\bullet}R \to \GG^{\bullet}P$ , $\GG^{\bullet}R^{\prime} \to \GG^{\bullet}P^{\prime}$ are $\redTopSeq$-cofibrations, and that 
		\item The vertical maps are weak equivalences
	\end{itemize}
Now, note that in the model category $\redTopSeq$, pushout of any map along a cofibration coincides with the homotopy pushout (cf. \cite{}). Thus it follows that the pushouts of both the rows in \eqref{comparepropositiondig} coincide with the homotopy pushouts of the underlying diagrams. Since the vertical arrows in \eqref{comparepropositiondig} are weak-equivalences, it then follows that the two pushotus are weak-equivalent.	 This completes the proof of the Proposition.
	\end{proof}
	
	The inclusion of diagrams \eqref{inclusionofdiagrams} satisfies the hypotheses of the proposition. Thus:
	\begin{corollary}
	\[|B_\bullet \mmfrfull| \sqcup^{\Galg}_{|B_\bullet \Mfrunstth|}|B_\bullet \Mhatunstth| \xrightarrow{\sim} |B_\bullet \Mfr| \sqcup^{\Galg}_{|B_\bullet \Mfrunst|} |B_\bullet \Mhatunst|\]
	 is a weak equivalence. In particular, $|B_\bullet \Mfr| \sqcup^{\Galg}_{|B_\bullet \Mfrunst|} |B_\bullet \Mhatunst|$ computes the homotopy pushout of \eqref{thm1dgm}
	\end{corollary}


\subsection{$|B_\bullet \Mfr| \sqcup^{\Galg}_{|B_\bullet \Mfrunst|} |B_\bullet \Mhatunst| \to \Mhat$ is a weak homotopy equivalence}\label{finalhtpyeq}
From the computation \eqref{barpushout}, we have:
\[|B_\bullet \Mfr| \coprod_{|B_\bullet\Mfrunst|}^{\Galg}  |B_\bullet \Mhatunst|  \big)  \Big|  \simeq  \Big|\{ \GG \big(\GG^{\bullet} \Mfr \coprod_{\GG^{\bullet}\Mfrunst} \GG^{\bullet}\Mhatunst \big) \Big|\]

We prove Theorem \ref{mthm1} using the following strategy: We will prove that the map
\[ \pi \colon \Big|\{ \GG \big(\GG^{\bullet} \Mfr \coprod_{\GG^{\bullet}\Mfrunst} \GG^{\bullet}\Mhatunst \big) \Big| \to \Mhat \]
satisfies the property that every $\Sigma \in \Mhat$ has a neighborhood $U_{\Sigma}$ such that for every finite collection $\Sigma_1,...,\Sigma_k$, 
\begin{equation}
	\label{finintwheq} \pi^{-1}(U_{\Sigma_1,...,\Sigma_k}) \to U_{\Sigma_1,...,\Sigma_k} \mbox{ is a weak homotopy equivalence}
\end{equation}
where $U_{\Sigma_1,...,\Sigma_k} = U_{\Sigma_1} \cap ... \cap U_{\Sigma_k}$.

\begin{remark} Note that here we are using a local criterion for weak equivalences (see for example, \cite[Corollary 1.4]{may1990weak}). We have to resort to this since $\pi$ is not a fibration in general and hence just showing the contractibility of the fiber may not be sufficient to imply weak equivalence. However to understand the contracting homotopies constructed below, particularly those in \ref{}, it might be helpful to think of them as ways of extending contracting homotopies of fibers over points to those of the inverse images of neighborhoods of these points.
\end{remark}

In fact, for every $\Sigma \in \Mhat$ we shall construct a neighborhood $U_{\Sigma}$ satisfying the following properties: $\pi^{-1}(U_{\Sigma})$ has a filtration
\[ F^1\pi^{-1}(U_{\Sigma}) \subset ... \subset F^n\pi^{-1}(U_{\Sigma}) ... \subset \pi^{-1}(U_{\Sigma})\]
such that each $F^i \pi^{-1}(U_{\Sigma})$ satisfies the following property:
\begin{equation}\label{filtrantncondition}
	\begin{split}
		&F^n \pi^{-1}(U_{\Sigma})  \mbox{ admits a fiberwise contraction homotopy, \emph{inside} } \pi^{-1}(U_{\Sigma}),\\[-4pt]
		&\mbox{ onto a section } \chi_n \colon U_{\Sigma} \to \pi^{-1}(U_{\Sigma}) \mbox{ of } \pi|_{U_{\Sigma}}
	\end{split}
\end{equation}


%\begin{enumerate}\label{filtrantncondition}
%	\item $F^i \pi^{-1}(U_{\Sigma})$ is fiberwise non-empty
%	\item $F^i \pi^{-1}(U_{\Sigma})$ admits a fiberwise contraction onto a section $\chi_i \colon U_{\Sigma} \to \pi^{-1}(U_{\Sigma})$ of $\pi|_{U_{\Sigma}}$.
%\end{enumerate}

In particular, it will suffice to verify that \eqref{finintwheq} is satisfied for the collection consisting of a single element.


%for any non-empty open $V \subset U_{\Sigma_0}$,  $\pi^{-1}(V)$ has an open cover $\{W_{n,\Sigma_0} \cap \pi^{-1}(V)\}$ such that every finite intersection is non-empty, connected and and contractible in $\pi^{-1}(V)$. It then follows that $\pi^{-1}(V)$ is contractible.
%for any open subset $V \subset U_{\Sigma_0}$, the inclusion $W_{\Sigma_0}^{\epsilon} \hookrightarrow W_{\Sigma_0}^{\epsilon/2}$ is null-homotopic. This will show that $\pi^{-1}(V) \to V$ is a weak homotopy equivalence.

%%%%%%%%%%%%%%%%%%%%%%%%%%%%%%%%%%%%%%%%%%%%%%%%%%%%%%%%
%%%%%%%%%%%%%%%%%%%%%%%%%%%%%%%%%%%%%%%%%%%%%%%%%%%%%%%%

\subsubsection{Construction of $U_{\Sigma}$ and $F^n \pi^{-1}(U_{\Sigma})$}\label{gammasubsubsect}
Let us start by describing the open set $U_{\Sigma}$: Let $\widehat{\Sigma} = \coprod \Theta_i$ be the normalization of $\Sigma$, where $\Theta_i$ are smooth connected riemann surfaces. For every node $\nu$ of $\Sigma$, denote by $\hat{\nu}, \check{\nu}$ its pre-images in $\widehat{\Sigma}$. Let $\un{\nu}$ be the collection of the points $\hat{\nu}, \check{\nu}$ as $\nu$ varies over the nodes of $\Sigma$. Let $\un{\nu}_i$ denote the subcollection of points on $\Theta_i$. Now, let $U_{\Theta_i,\un{\nu}_i}$ be a small neighborhood of $\Theta_{i,\un{\nu}_i}$ in the moduli of riemann surfaces with the same topological type as $\Theta_i$ and with internal marked points indexed by $\un{\nu}_i$. Let $\cfam_{\Theta_i, \un{\nu}_i} \to U_{\Theta_i,\un{\nu}_i}$ be the universal curve over it, and let $s_\nu$ be the sections corresponding to the marked points $\nu \in \un{\nu}_i$.

Consider the product family 
\[ \cfam_{\widehat{\Sigma}, \un{\nu}} = \prod \cfam_{\Theta_i,\un{\nu}_i} \to U_{\widehat{\Sigma}, \un{\nu}} =  \prod \cfam_{\Theta_i, \un{\nu}_i}\]
along with the sections $s_{\nu}, \nu \in \un{\nu}$. Now let $z_\nu$ be a holomorphic function in a neighborhood of Im$(s_{\nu})$ which restricts to a holomorphic co-ordinate on each fiber. Shrinking $U_{\Theta_i,\un{\nu}_i}$ and rescaling the $z_{\nu}$ if necessary, we may assume that $R_{\nu} = \{ |z_{\nu}| <1\}$ contains no marked points other than $\nu$ and are isomorphic to $U_{\Sigma_{\theta,\un{\nu}}} \times \DD$. $t_{\hat{\nu}} = t_{\check{\nu}} \colon \cfam \times \DD^{\un{\nu}} \to \DD$ be the projection onto the $\nu$th component of $\DD^{\nu}$. Now, set $U_{\Sigma}$ to be the  neighborhood of $\Sigma$ in $\Mhat$ given by the 
\[ U_{\Sigma} := U_{\widehat{\Sigma}} \times \DD^{\un{\nu}} \]


The restriction of the universal curve over $U_{\Sigma}$
\[ \cfam_{\Sigma} \to U_{\Sigma} = U_{\widehat{\Sigma} \times \DD^{\un{\nu}}} \]
is obtained as follows: remove the closed subsets 
\[ \{(z,\un{t}) \in R_{\nu} \times \DD^{\un{\nu}} | |z_{\nu}(z)| \leq t_{\nu} \} \]
from $\cfam_{\widehat{\Sigma}} \times \DD^{\un{\nu}}$ and identify rest of $R_{\hat{\nu}}$ and $R_{\check{\nu}}$ via 
\[ z_{\hat{\nu}} \cdot z_{\check{\nu}} = t_{\nu} \]


Let us now turn to the construction of the filtration $F^i\pi^{-1}(U_{\Sigma})$ satisfying conditions \eqref{filtrantncondition}.

Let $\un{x}$ be an element in $\GG^1 \Mfr \coprod_{\GG^1 \Mfrunst} \GG^1 \Mhatunst$. It is an $\left( \Mfr \coprod_{\Mfrunst} \Mhatunst  \right)$-labeled graph, say with underlying rda-graph $G$ and labeling $\prod_{v \in V(G)} \Sigma_{v}$. Now Consider the map
\[ \GG^1 \Mfr \coprod_{\GG^1 \Mfrunst} \GG^1 \Mhatunst \to \Mhat\]
% = \prod_{v \in V(G)} \Sigma_{v}$, an $\Mfr \coprod_{\Mfrunst} \Mhatunst$-labeled graph, 
It maps $\un{x}$ to the riemann surface $\Sigma_{\un{x}} = \#_{\prod_{v \in V(G)}} \Sigma_{v}$ obtained by gluing the $\Sigma_v$'s as prescribed by the edges of $G$. In particular, we get an analytic embedding 
\[ \label{gammacircle} \gamma_{\un{x}} \colon \coprod_{E(G) \cup in(G) \cup out(G)} S^1 \to \Sigma_{\un{x}}\]
whose image lies away from the nodes of $\Sigma_{\un{x}}$.

Note that the maps $ \GG^{n+1} \Mfr \coprod_{\GG^{n+1}  \Mfrunst} \GG^{n+1} \Mhatunst \to \Mhat$ factor as 
\[  \GG^{n+1} \Mfr \coprod_{\GG^{n+1}  \Mfrunst} \GG^{n+1} \Mhatunst \to  \GG^{1} \Mfr \coprod_{\GG^{1}  \Mfrunst} \GG^{1} \to \Mhat \] 
where the first map is the composition of simplicial degeneracy maps $s_1 \circ \hdots \circ s_{n-1} $ and the second map is as described above. In particular, to any 
\[\un{x} \in  \GG^{n+1} \Mfr \coprod_{\GG^{n+1}  \Mfrunst} \GG^{n+1} \Mhatunst,\]
we can assign an embedding $\gamma_{\un{x}}$ as in \ref{gammacircle}. For an example see Figure \ref{gamma_fig}.

%%%%%%%%%%%%%%%%%%%%%%%%%%%%%%%%%%%%%%%%%%%%%%%%%%%%%%
%%%%%%%%%%%%%%%%                      Figure                 %%%%%%%%%%%%%%%%%%%%%%%

\begin{figure}
	\centering
		\includegraphics[width=\textwidth]{../img/gamma_fig.png}
		\caption{}
		\label{gamma_fig}
	\end{figure}

%%%%%%%%%%%%%%%%%%%%%%%%%%%%%%%%%%%%%%%%%%%%%%%%%%%%%%
%%%%%%%%%%%%%%%%%%%%%%%%%%%%%%%%%%%%%%%%%%%%%%%%%%%%%%

For each $n$ and $\nu$, let $\widetilde{A}^n_{\nu} \subset \cfam_{\Sigma}$ be the open set defined by 
\[ \widetilde{A}^n_{\nu}  =  \{ (z, \un{t}) \in R_{\nu} \times \DD^{\Nu} | |t_\nu| \leq \frac{1}{2n^2}, |z_{\hat{\nu}}(z) | \in [ |t_{\nu}| n, \frac{1}{n}] \} \]
Set 
\[\widetilde{A}^n= \bigcup_{\nu \in \Nu} \widetilde{A}^n_{\nu}\]

Define $\WW^n_{\Sigma}$ to be the subcomplex of $\GG^{\bullet+1} \Mfr \coprod_{\GG^{\bullet+1}  \Mfrunst} \GG^{\bullet+1} \Mhatunst$ consisting of the simplices $\un{x}$ such that the associated embedding $\gamma_{\un{x}}$ (as above) is disjoint from $\widetilde{A}^n$.  Define $F^n\pi^{-1}(U_{\Sigma})$ to be the \textcolor{red}{open subset} of $\left| \GG^{\bullet+1} \Mfr \coprod_{\GG^{\bullet+1}  \Mfrunst} \GG^{\bullet+1} \Mhatunst \right|$ associated with this subcomplex.\\
Let $U_{\Sigma}^{sing}$ denote the subset
\[ \{ (z,\un{t})| |t_{\nu}| =0 \mbox{ for some } \nu\} \]
This is the locus of nodal riemann surfaces in $U_{\Sigma}$. Denote by $\nu_n(U_{\Sigma}^{sing})$ the neighborhood of $U_{\Sigma}^{sing}$ given by
\[  \{(z,\un{t})| |t_{\nu}| \leq \frac{1}{2n^2} \mbox{ for all } \nu\} \]
In the constructions below we will also use the subcomplex $\WW^n_{\Sigma} \cap \pi^{-1}(\nu_n(U_{\Sigma}^{sing}))$. Denote this by $\WW^n_{\Sigma}|_{\nu_n(U_{\Sigma}^{sing})}$. The subset of $\pi^{-1}(U_{\Sigma})$ associated to this is given by $F^n (U_{\Sigma}) \cap \pi^{-1}(\nu_n(U_{\Sigma}^{sing}))$.
%%%%%%%%%%%%%%%%%%%%%%%%%%%%%%%%%%%%%%%%%%%%%%%%%%%%%%
%%%%%%%%%%%%%%%%                      Figure                 %%%%%%%%%%%%%%%%%%%%%%%

%\begin{figure} 
%	\centering 
%		\includegraphics[width=\textwidth]{../img/filteration_fig.png}
%		\caption{}
%
%\end{figure}

\begin{figure} 
	\centering
	\begin{subfigure}[b]{0.3\textwidth}
		\includegraphics[width=\textwidth]{../img/filteration_fig.png}
		\caption{Open sets $W_{n,\Sigma}$}
		\label{filtrationfig1}
	\end{subfigure}
	\begin{subfigure}[b]{0.35\textwidth}
		\includegraphics[width=\textwidth]{../img/filteration_fig2.png}
		\caption{Circles $S_{\hat{\nu}}, S_{\check{\nu}}$}
		\label{filtrationfig1}
	\end{subfigure}
	\caption{}
\end{figure}


%%%%%%%%%%%%%%%%%%%%%%%%%%%%%%%%%%%%%%%%%%%%%%%%%%%%%%
%%%%%%%%%%%%%%%%%%%%%%%%%%%%%%%%%%%%%%%%%%%%%%%%%%%%%%
%%%%%  																						  %%%%% 
%%%%% 							Subsection Not needed ?										   %%%%% 
%%%%% 																						   %%%%% 
%%%%%%%%%%%%%%%%%%%%%%%%%%%%%%%%%%%%%%%%%%%%%%%%%%%%%%
%%%%%%%%%%%%%%%%%%%%%%%%%%%%%%%%%%%%%%%%%%%%%%%%%%%%%%
%%%%%%%%%%%%%%%%%%%%%%%%%%%%%%%%%%%%%%%%%%%%%%%%%%%%%%


%\subsubsection{Proof of condition (1) in \ref{Wsigmacond}}\label{cond1}
%Pick any $M \in \NN$, $M> max(I)$ and a $\Sigma \in U_{\Sigma}$ such that the co-ordinates of whose co-ordinates $(z(\Sigma),\un{t}(\Sigma)) \in  U_{\Sigma} = U_{\hat{\Sigma}_0} \times \DD^{\Nu}$ satisfy,
%\[ |t_\nu(\Sigma)| \leq \frac{1}{2n^2} \mbox{ for every } \nu \]
%Consider an embedding $\gamma \colon \bigsqcup_{\Nu} (S^1 \sqcup S^1) \to \Sigma$ given by the circles
%\[ S^{\hat{\nu}}(\Sigma) = \{(z,t_{\nu}(\Sigma)) | |z_{\nu}(z)| = \frac{1}{M} \} \subset \Sigma, S^{\check{\nu}}(\Sigma) = \{(z,t_{\nu}(\Sigma)) | |z_{\nu}(z)| = t_{\nu}(\Sigma)M \} \]
%There is an unique element $\un{x} \in \GG^1 \Mfr \coprod_{\GG^1  \Mfrunst} \GG^1 \Mhatunst$, such that the curve $\gamma$ coincides with the associated embedding $\gamma_{\un{x}}$. Then labels of the vertices of $\un{x}$ are given by the connected components of $\Sigma \setminus Im(\gamma)$, with the edges corresponding to all circles $S^{\hat{\nu}}(\Sigma), S^{\check{\nu}}(\Sigma)$ directed away from the vertex associated to the component $\{(z,t_{\nu}(\Sigma) |  |z_\nu(z)| \in [t_\nu M , \frac{1}{M}] \}$ bounded by $S^{\hat{\nu}}(\Sigma), S^{\check{\nu}}(\Sigma)$. Note that $\un{x} \in \bigcap_{I}W_{n,\Sigma}$. This proves the fiberwise non-emptiness condition. 

%\subsubsection{Construction of sections $\chi_n \colon U_{\Sigma} \to \pi^{-1}(U_{\Sigma})$ as in Condition \eqref{filtrantncondition}}
%For $\Sigma^{\prime} \in U_{\Sigma}$, consider the graph $G_{\Sigma^{\prime}}$ with 
%\begin{itemize}
%	\item The set of vertices given by the connected components of $\Sigma^{\prime} \setminus \bigcup_{\Nu} \partial A_{\nu}(\Sigma^{\prime})$
%	\item The edges correspond to the the connected components of $\bigcup_{\Nu} \partial A_{\Nu}$, with all edges directed away from vertices labeled by the annuli $A_{\Nu}$.
%\end{itemize}
%let $\un{x}_{\Sigma^{\prime}}$ be the $0$-simplex with the underlying graph $G_{\Sigma^\prime}$ and with the labeling of each vertex given by the corresponding connected component of $\Sigma^{\prime} \setminus \partial A_{\nu}$. Set 
%$\chi_i(\Sigma^{\prime}) = \un{x}_{\Sigma^{\prime}}$. It is not difficult to see that this indeed defines a continuous section of $\pi|_{\pi^{-1}}(U_{\Sigma})$.

\subsubsection{Proof of  condition \eqref{filtrantncondition}}\label{prfcondfiltrtop}
Let us now turn to the fiberwise contraction. We begin by describing a number of preliminary notions and maps:
	\paragraph{\textbf{Curves $S^{\check{\nu}}, S^{\hat{\nu}}$ and the annuli $A_{\nu}$}}\label{curveannuluspara}
For any $\Sigma^{\prime} \in \nu_n(U_{\Sigma})$
% whose co-ordinates $(z(\Sigma^{\prime}),\un{t}(\Sigma^{\prime})) \in  U_{\Sigma} = U_{\hat{\Sigma}} \times \DD^{\Nu}$ satisfy
%\[ |t_\nu(\Sigma^{\prime})| \leq \frac{1}{2n^2} \mbox{ for every } \nu \]
Denote by $A_{\nu}(\Sigma^\prime)$ the annulus $ \widetilde{A}^{2n}_{\nu} \cap \Sigma^\prime$, where we think of $\Sigma^\prime$ as sitting inside the universal curve $\cfam_{\Sigma}$ as a fiber. Denote by $S^{\check{\nu}}, S^{\hat{\nu}}$ boundaries of this annuls given by 
\[ S^{\hat{\nu}}(\Sigma^{\prime}) = \{(z,t_{\nu}(\Sigma^{\prime})) | |z_{\nu}(z)| = \frac{1}{2n} \} , S^{\check{\nu}}(\Sigma^{\prime}) = \{(z,t_{\nu}(\Sigma^{\prime})) | |z_{\nu}(z)| =2 t_{\nu}(\Sigma^{\prime}) n \} \subset \Sigma^{\prime}\]
Now consider the graph $G_{\Sigma^{\prime}}$ with 
\begin{itemize}
	\item The set of vertices given by the connected components of $\Sigma^{\prime} \setminus \bigcup_{\Nu} \partial A_{\nu}(\Sigma^{\prime})$
	\item The edges correspond to the the connected components of $\bigcup_{\Nu} \partial A_{\Nu}(\Sigma^{\prime})$, with all edges directed away from vertices labeled by the annuli $A_{\Nu}(\Sigma^{\prime})$.
\end{itemize}
Denote by $\un{x}_{\Sigma^{\prime}}$ the $0$-simplex with the underlying graph $G_{\Sigma^\prime}$ and with the labeling of each vertex given by the corresponding connected component of $\Sigma^{\prime} \setminus \partial A_{\nu}(\Sigma^{\prime})$. 
%Set $\chi_i(\Sigma^{\prime}) = \un{x}_{\Sigma^{\prime}}$. It is not difficult to see that this indeed defines a continuous section of $\pi|_{\pi^{-1}}(U_{\Sigma})$.\\
%There is an unique element $\un{x} \in \GG^1 \Mfr \coprod_{\GG^1  \Mfrunst} \GG^1 \Mhatunst$, such that the curve $\gamma$ coincides with the associated embedding $\gamma_{\un{x}}$. Then labels of the vertices of $\un{x}$ are given by the connected components of $\Sigma \setminus Im(\gamma)$, with the edges corresponding to all circles $S^{\hat{\nu}}(\Sigma), S^{\check{\nu}}(\Sigma)$ directed away from the vertex associated to the component $\{(z,t_{\nu}(\Sigma) |  |z_\nu(z)| \in [t_\nu n , \frac{1}{n}] \}$ bounded by $S^{\hat{\nu}}(\Sigma), S^{\check{\nu}}(\Sigma)$.\\
%Fix any $M \in I$. Let $\Sigma \in \pi(W_{I,\Sigma_0})$. Let $S^{\hat{\nu}}$, and $S^{\check{\nu}}$ be defined as in \ref{cond1}. 
	\paragraph{\textbf{The graphs $H_v$}}\label{hvgrph}\textcolor{red}{Need to explain why result of applying $H_v$ to every vertex if of type `at least one output' cf. correction\_1.png}. Let $\un{x} \in \WW^{n}_{\Sigma}|_{\nu_n(U_{\Sigma}^{sing})}$ with underlying rda-graph $G$. Let $\un{x}_v$ denote the label of $v \in V(G)$ in $\un{x}$ and let $\Sigma^{\prime}_v$ be the surface $\pi(\un{x}_v)$. Since $\un{x}$ lies in $\WW^{n}_{\Sigma}$ it follows that for any $\nu \in \Nu$, $A_{\nu}(\Sigma^{\prime}) \cap \Sigma^{\prime}_v$ is either empty or $A_{\nu}^i(\Sigma^{\prime})$. We can then construct a graph $H_v$ as follows:
		\begin{itemize}
			\item The vertices of $H_v$ are given by the connected components of $\Sigma^{\prime}_v \setminus \bigcup_{\Nu} \partial A_{\nu}$
			\item Edges of $H_v$ are given by the connected components of $\bigcup_{\Nu} \partial A_{\nu}$ lying inside $\Sigma^{\prime}_v$, with each edge directed away from the vertex corresponding to a component of the form $A_{\nu}$ for some $\nu$.
		\end{itemize}
\begin{notation}
In what follows, for the sake of brevity we will use $\PP_{\Sigma}$ to denote the simplicial complex underlying $\pi^{-1}(U_{\Sigma})$. The simplicial complex underlying $\pi^{-1}(\nu_n(U_{\Sigma}^{sing}))$ will be denoted by $\PP_{\Sigma}|_{\nu_n(U_{\Sigma}^{sing})}$
\end{notation}
	\paragraph{\textbf{The cut maps}}\label{cutmaps} We define a simplical map $c \colon \WW^{n}_{\Sigma}|_{\nu_n(U_{\Sigma}^{sing})} \to \PP_{\Sigma}|_{\nu_n(U_{\Sigma}^{sing})}$ which will be used in the definitions of homotopies below. Intuitively, it corresponds to making `cuts' to the underlying riemann surface along the circles $S^{\hat{\nu}}$ and $S^{\check{\nu}}$.
	Maps $c_r \colon (\WW^{n}_{\Sigma}|_{\nu_n(U_{\Sigma}^{sing})})^r \to (\PP_{\Sigma}|_{\nu_n(U_{\Sigma}^{sing})})^r$ are defined inductively as follows:
		\begin{itemize}
			\item Base case: Let $\un{x} \in |_{\nu_n(U_{\Sigma}^{sing})}$. Let $G$ be the graph underlying $\un{x}$. For $v \in G_{\un{x}}$ let $\Sigma_v$ and $H_v$ be as described above in \ref{hvgrph}. Consider an element $\un{y}_v \in \PP^1_{\Sigma_v}$ whose underlying rda-graph is $H_v$ and whose vertex labels are given by the corresponding connected components of $\Sigma_v \setminus \bigcup_{\Nu} \partial A_{\nu}$.  We then define $c_0(\un{x})$ as the graph composition 
			\begin{equation}\label{c0eq} c_0(\un{x}) := \mu_G( \prod \un{y}_v ) \end{equation}
			\item Assume inductively that $c_i, 0 \leq i <r$ has been defined and satisfies the following property: For any $\un{x} \in(\PP_{\Sigma}|_{\nu_n(U_{\Sigma}^{sing})})^i$, with the underlying graph $G$, $c_{i} (\un{x}) = \mu_{G}( \prod \un{y}_v)$ for some $\un{y}_v \in (\PP_{\Sigma}|_{\nu_n(U_{\Sigma}^{sing})})^i$. Notice that the condition is satisfied for $i=0$.\\
		Now suppose $\un{x} \in(\PP_{\Sigma}|_{\nu_n(U_{\Sigma}^{sing})})^r$, having the underlying rda-graph $G$. Consider element $c_{r-1}(d_r \un{x}) \in(\PP_{\Sigma}|_{\nu_n(U_{\Sigma}^{sing})})^{r-1}$. From the induction hypothesis we know that $c_{r-1}(d_r \un{x}) =  \mu_{G_{d_r \un{x}}}(\prod \un{y}^{\prime \prime}_{w})$ for some $\un{y}^{\prime \prime}_w$, $w \in V( G_{d_r\un{x}} )$ where $G_{d_r\un{x}}$ is the rda-graph underlying $d_r\un{x}$. But since  $G_{d_r \un{x}}  = \mu_G( \prod G_v)$ for some graphs $G_v$, we can express $c_{r-1}(d_{r-1} \un{x})$  as a graph composition $\mu_G( \prod \un{y}^{\prime}_v )$ for some $\un{y}^{\prime}_v \in (\PP_{\Sigma}|_{\nu_n(U_{\Sigma}^{sing})})^{r-1}$. For $v \in G$, consider an element $\un{y}_v \in \PP^n_{\Sigma}$ defined as follows: the rda-graph underlying $\un{y}_v$ is $H_v$.The label of a vertex $u \in  H_v$ in $\un{y}_v$ is the labeled subgraph of $\un{y}^{\prime}_v$ consisting of vertices and edges mapping under $\pi$ into the connected component $\Sigma_u$ of $\Sigma \setminus \bigcup_{\Nu}\partial A_{\nu}$,  as in \ref{hvgrph}, corresponding to $u$. Then, we set 
		\[c_r(\un{x}) = \mu_G( \prod \un{y}_v)\]
		\end{itemize}

\paragraph{\textbf{The cut homotopies}} We define a fiber preserving homotopy $\Phi \colon F^n (U_{\Sigma}) \cap \pi^{-1}(\nu_n(U_{\Sigma}^{sing})) \to \pi^{-1}(\nu_n(U_{\Sigma}))$. This is defined as the geometric realization of a simplicial homotopy $\phi \colon \WW^{n}_{\Sigma}|_{\nu_n(U_{\Sigma}^{sing})} \to \PP_{\Sigma}|_{\nu_n(U_{\Sigma}^{sing})}$. This in turn consists of maps $\phi^j \colon(\WW_{\Sigma}|_{\nu_n(U_{\Sigma}^{sing})})^r \to (\PP_{\Sigma}|_{\nu_n(U_{\Sigma}^{sing})})^r$ for $0 \leq j \leq i$, satisfying the usual conditions of a simplicial homotopy.\\
\begin{itemize}
\item Base Case: Let $\un{x} \in(\PP_{\Sigma}|_{\nu_n(U_{\Sigma}^{sing})})^0$ and let $G$ be the graph underlying $\un{x}$. Suppose that $c_0(\un{x})=\mu_G(\prod \un{y}_v)$. Then, we define 
	\[ \phi^0(\un{x}) := \mu_G( \prod \eta(\un{y}_v)) \]
	recall from \ref{GGtripsubsec}, $\eta \colon \mathbbm{1}_{\redTopSeq} \to \GG$ is the unit of the monad $\GG$.
\item Assume inductively that the maps $\phi^j$ have been defined for all degrees $i < r$ and satisfy the following property similar to the one used above in \ref{cutmaps}: For any $\un{x} \in (\WW^{n}_{\Sigma}|_{\nu_n(U_{\Sigma}^{sing})})^{i}$, with the underlying graph $G$, $\phi^{j} (\un{x}) = \mu_{G}( \prod \un{y}_v)$ for some $\un{y}_v \in  (\WW^{n}_{\Sigma}|_{\nu_n(U_{\Sigma}^{sing})})^{i}$. The maps $\phi^j \colon (\WW^{n}_{\Sigma}|_{\nu_n(U_{\Sigma}^{sing})})^{i} \to (\PP_{\Sigma}|_{\nu_n(U_{\Sigma}^{sing})})^{i+1}$ are defined as follows:\\
Let $\un{x} \in  (\WW^{n}_{\Sigma}|_{\nu_n(U_{\Sigma}^{sing})})^{i}$ and $G$ be the rda-graph underlying $\un{x}$. Then define $\phi^j(\un{x})$ as follows:
	\begin{itemize}
		\item $j <r$: From the condition in the inductive hypothesis we have $\phi^j(d_r\un{x}) = \mu_G( \prod \un{y}_v)$, for some $\un{y}_v$. Then define
		 \[\phi^j(\un{x}) := \mu_G ( \prod \eta(\un{y}_v) )\]
		\item $j=r$: For $\un{x} \in  (\WW^{n}_{\Sigma})^{r}$, suppose that $c_r(\un{x}) = \mu_G(\prod \un{y}_v)$. Then define
		\[ \phi^j(\un{x}) := \mu_G( \prod \eta(\un{y}_v)) \]
	\end{itemize}
\end{itemize}

\paragraph{\textbf{The homotopy $\Psi$}}\label{psiproprmain} Let $\PP_{\Sigma}|_{\nu_n(U_{\Sigma}^{sing})}^M$ be the subcomplex of $\PP_{\Sigma}|_{\nu_n(U_{\Sigma}^{sing})}$ consisting of simplices $\un{x}$ such that if $\un{x}_v$  is the label of a vertex in $\un{x}$ and $\nu \in \Nu$, then $\pi(\un{x}_v) \cap A_{\nu} \neq \phi \Rightarrow \un{x}_v = \eta^r(A_{\nu})$. Let $\pi^{-1}(U_{\Sigma})^M$ be th geometric realization of $\PP_{\Sigma}|_{\nu_n(U_{\Sigma}^{sing})}^M$.
Let $\chi^M_n \colon U_{\Sigma}|_{\nu_n(U_\Sigma^{sing})}  \to \PP_{\Sigma}|_{\nu_n(U_{\Sigma}^{sing})}^M$ be a section of $\pi$ over $U_{\Sigma}|_{\nu_n(U_\Sigma^{sing})}$ defined as follows: the image of $\chi^M_n $ is contained inside the $0$-skeleton $(\PP_{\Sigma}|_{\nu_n(U_{\Sigma}^{sing})})^0$ with 
\[ \chi^M_n \colon \Sigma^\prime \mapsto \un{x}_{\Sigma^{\prime}} \mbox{ for every } \Sigma^{\prime} \in \nu_n(U_{\Sigma}^{sing})\]
where $\un{x}_{\Sigma^{\prime}}$ is a $0$-simplex as defined in \ref{curveannuluspara}. Denote the simplicial complex underlying the image of $\chi^M_n$ by $\mathbb{X}^M_n$.

% and $\nu \in \Nu$, then $\pi(\un{x}_v) \cap A_{\nu} \neq \phi \Rightarrow \pi(\un{x}_v) \subset A_{\nu}$. Note that $Im(c) \subset \PP^M_{\Sigma_0}$. Let $\pi^{-1}(U_{\Sigma_0})^M$ denote the geometric realization of $\PP^M_{\Sigma_0}$.\\
%Let $\Ss^M_{\Sigma_0}$ be a further subcomplex contained inside the $0$-skeleton of $\PP^M_{\Sigma_0}$. It consists of $0$-simplies $\un{x}_{\Sigma}, \Sigma \in U_{\Sigma_0}$ defined as follows: the underlying graph $G_{\Sigma}$ has 
%	\begin{itemize}
%		\item The set of vertices given by the connected components of $\Sigma \setminus \bigcup_{\Nu} \partial A_{\nu}$
%		\item The edges correspond to connected components of $\bigcup_{\Nu} \partial A_{\nu}$, with all edges directed away from vertices labeled by the annuli $A_{\nu}$.
%	\end{itemize} $\un{x}_{\Sigma}$ is obtained by labeling each vertex by the corresponding connected component of $\Sigma \setminus \bigcup_{\Nu} \partial A_{\nu}$. $S^M_{\Sigma_0} \subset \pi^{-1}(U_{\Sigma_0})^M$ be its geometric realization. $\Ss^M_{\Sigma_0}$ is the constant simplcial space associated with $S^M_{\Sigma_0}$. Also, notice that $S^M_{\Sigma_0}$ 
defines a (continuous) section of the map $\pi \colon \pi^{-1}(U_{\Sigma_0}) \to U_{\Sigma_0}$.\\

We will construct a fiber-preserving homotopy $\Psi$ contracting $\pi^{-1}(U_{\Sigma})^M$ into $\mathbb{X}^M_n$. This will again be a geometric realization of a simplcial homotopy, which itself will be constructed by specifying extra degeneracies and augmentation on $\PP_{\Sigma}|_{\nu_n(U_{\Sigma}^{sing})}^M$ (cf. \cite[Section 6]{meyer1984bar}). For this we need to define:
	\begin{enumerate}
	 \item Extra degeneracies: $s_{n+1} \colon (\PP_{\Sigma}|_{\nu_n(U_{\Sigma}^{sing})}^M)^n \to (\PP_{\Sigma}|_{\nu_n(U_{\Sigma}^{sing})}^M)^{n+1}$ for all $n \geq -1$, where $(\PP_{\Sigma}|_{\nu_n(U_{\Sigma}^{sing})}^M)^{-1} := \mathbb{X}^M_n$
	 \item Augmentations: $d_0 \colon (\PP_{\Sigma}|_{\nu_n(U_{\Sigma}^{sing})}^M)^0 \to \mathbb{X}^M_n$
	\end{enumerate}
satisfying:
	\begin{enumerate}
		\item $d_{n+1}s_{n+1}$ is identity on $(\PP_{\Sigma}|_{\nu_n(U_{\Sigma}^{sing})}^M)^n$
		\item for $n \geq 1, 0 \leq i,j \leq n$, 
			\begin{enumerate}
			\item $d_{i}s_{n+1} = s_{n}d_i$
			\item $s_{j+1}s_{n+1} = s_{n+2}s_j$
			\end{enumerate}
		on $(\PP_{\Sigma}|_{\nu_n(U_{\Sigma}^{sing})}^M)^n$
	\end{enumerate}
	The contracting homotopy $\psi$ underlying $\Psi$ is then obtained by setting $\psi_i \colon (\PP_{\Sigma}|_{\nu_n(U_{\Sigma}^{sing})}^M)^n \to (\PP_{\Sigma}|_{\nu_n(U_{\Sigma}^{sing})}^M)^{n+1}$, $0 \leq i \leq n$, to be $s_n \hdots s_{i+1}d_{i+1} \hdots d_n s_{n+1}$.
	
	\begin{itemize}\label{extradegpsi}
	\item Extra degeneracies: We define $s_{n+1} \colon (\PP_{\Sigma}|_{\nu_n(U_{\Sigma}^{sing})}^M)^n \to (\PP_{\Sigma}|_{\nu_n(U_{\Sigma}^{sing})}^M)^{n+1}$. Let $\un{x} \in (\PP_{\Sigma}|_{\nu_n(U_{\Sigma}^{sing})}^M)^n$. Let $G_{\pi(\un{x})}$ be the graph for the riemann surface $\Sigma^\prime$ defined in \ref{prfcondfiltrtop}. From the definition of $\PP_{\Sigma}|_{\nu_n(U_{\Sigma}^{sing})}^M$ it follows that the vertices of $\un{x}$ can be regrouped into subgraphs $\un{y}_v \in  (\PP_{\Sigma}|_{\nu_n(U_{\Sigma}^{sing})}^M)^n , v\in G_{\pi(\un{x})}$ so that 
$\un{x} = \mu_{G_{\pi(\un{x})}} ( \prod \un{y}_v)$.
	Then  set 
	\[ s_{n+1}(\un{x}) := \mu_{G_{\un{x}}} ( \prod \eta(\un{y}_v)) \]
	\item Augmentation: Note that for $\un{x} \in  (\PP_{\Sigma}|_{\nu_n(U_{\Sigma}^{sing})}^M)^0 $, $d_0(s_1\un{x}) = \un{x}_{\pi(\un{x})} \in \Ss^M_{\Sigma}$. We define $d_0 \colon (\PP_{\Sigma}|_{\nu_n(U_{\Sigma}^{sing})}^M)^0 \to \mathbb{X}^M_n$ as
	\[ d_0 = d_0 \circ s_1 \]
	\end{itemize} 

\paragraph{\textbf{The homotopy ${\Psi}^{ns}$}}\label{psihtpy} We construct another fiber-preserving homotopy, this time contracting $\pi^{-1}(U^{ns}_{\Sigma})$ onto a section ${\chi}^{ns}$ of $\pi$ over $U^{ns}_{\Sigma}$. Here $U^{ns}_{\Sigma} = \{(z,\un{t})| t_{\nu} \neq 0 \mbox{ for all } \nu\} \subset U_{\hat{\Sigma}} \times \DD^{\Nu} = U_{\Sigma}$ is the subspace of non-singular curves in $U_{\Sigma}$. Let $\PP^{ns}_{\Sigma}$ be the corresponding subcomplex of $\PP_{\Sigma}$.\\
The section ${\chi}_n^{ns}$ is contained in the $0$-skeleton of $\PP_{\Sigma}^{ns}$ and is given by $\Sigma^\prime \mapsto \widetilde{\un{x}}_{\Sigma^\prime}$, where the graph underlying $\tilde{\un{x}}_{\Sigma^\prime}$ has a single vertex, no internal edges, and with the label of the vertex being simply $\Sigma^\prime$. Note that since $\Sigma^\prime$ is non-singular, $\widetilde{\un{x}}_{\Sigma^\prime}\in \GG \Mfr \subset (\PP_{\Sigma})^0$.


Let the subcomplex of  ${\PP}^{ns}_{\Sigma}$ corresponding to $\chi^{ns}$ be denoted by $\mathbb{X}^{ns}_{\Sigma}$. It is a constant simplcial space and is contained inside the $0$-skeleton of $\PP_{\Sigma}$.
Again we shall construct the contraction by specifying extra degeneracies and augmentation as in \ref{psihtpy}. We continue using the same notation for these as in \ref{psiproprmain}.
	\begin{itemize}
		\item Extra degeneracies: In this case the map $s_{n+1} \colon (\PP^{ns})^n \to  (\PP^{ns})^{n+1}$ is defined by 
		\[ s_{n+1}(\un{x}) _= \eta(\un{x}) \]
		\item The augmentation $d_0 \colon  (\PP^{ns})^0 \to \mathbb{X}^{ns}_{\Sigma}$ is defined by 
		\[ d_0 := d_0 \circ s_1 \]
	\end{itemize}
The verification of conditions satisfied by $s_{n+1}$'s and $d_0$ is analogous to \ref{psihtpy} and is omitted. The simplcial homotopy ${\psi}^{ns}$ underlying ${\Psi}^{ns}$ is obtained from $s_{n+1}$'s and $d_0$ by formulas similar to those in \ref{psihtpy}.

\paragraph{\textbf{The contraction for proving condition \eqref{filtrantncondition}}}\label{mthm1finalparag}
We are finally ready to define the homotopy needed to show that the condition \eqref{filtrantncondition} is satisfied:	We construct a homotopy $\Omega$ from the inclusion $F^n\pi^{-1}(U_{\Sigma}) \hookrightarrow \pi^{-1}(U_{\Sigma})$ to a map which takes $F^n\pi^{-1}(U_{\Sigma})$ onto a section of $\pi$ over $U_{\Sigma}$. The section coincides with $\chi_i$ near the singular locus $U_{\Sigma} \setminus U^{ns}_{\Sigma}$, and with $\chi^{ns}$ away from it.

Fix a continuous function $\zeta \colon U_{\Sigma} \to [0,1]$ such that $\zeta \equiv 1$ on $\bigcup_{\Nu} \{|t_{\nu}|\leq \frac{1}{4n^2}\}$ and $\zeta \equiv 0$ outside of $\bigcup_{\Nu} \{|t_{\nu}|\leq \frac{1}{2n^2}\}$. Set $\kappa = \zeta \circ \pi$. The homotopy $\Omega$ is defined as follows:

	\begin{align*}
	\mbox{ On } \pi^{-1}(\bigcup_{\Nu} \{|t_{\nu}|\leq \frac{1}{4n^2} \})&: \Omega(\un{x},t) = 
		\begin{cases} 
			\Phi(\un{x},2t), & t \leq 1/2 \\
			\Psi(\Phi(\un{x},1), 2t-1) & t >1/2
		\end{cases}\\
	\mbox{ On } \pi^{-1}(U^{ns}_{\Sigma})&: \Omega(\un{x},t) = 
		\begin{cases} 
			\Phi(\un{x},2t), & t \leq \kappa(\un{x})/2 \\
			\Phi({\Psi^{ns}}(\un{x},\frac{t-\kappa/2}{1-\kappa/2}), \kappa(\un{x})) & t > \kappa({\un{x}})/2
		\end{cases}
	\end{align*}
	To ensure that $\Omega$ is well-defined it suffices to check that the two definitions agree on points lying over $U^{ns}_{\Sigma} \bigcap \bigcup_{\Nu} \{|t_{\nu}|\leq \frac{1}{4n^2} \}$
	\[ \Psi ( \Phi (\un{x},1),t) = \Phi ({\Psi^{ns}} ( \un{x} , t) ,1 ) \mbox{ for all } t\]
Since $\Phi(\un{x},1)=c(\un{x})$, this amounts to checking $\psi_i c_n = c_{n+1} \psi^{ns}_{i}$ on $n$-simplices lying over $U^{ns}_{\Sigma} \cap \bigcup_{\Nu} \{|t_{\nu}|\leq \frac{1}{4n^2}|\}$. Since $c$ is simplcial, from the definitions of $\psi_i,\psi^{ns}_i$ we are reduced to checking 
\[ \psi_{n+1} c_n =  c_{n+1} \psi^{ns}_{n+1}\]
This can be verified directly from the definition of $\phi_{n+1}$ and $\psi^{ns}_{n+1}$.

\section{From $\Mfr$ to $\Mfrfull$}
The aim of this section is to show that homotopically a (full) reduced CFT can be obtained by adding a suitably compatible pairing and trace operation to a positive boundary reduced CFT


\begin{theorem}\label{mainthm2}
\[
	\begin{tikzcd}
		\Mfrunst \ar{r}\ar{d} & \Mfrfullunst \ar{d} \\
		\Mfr \ar{r} & \Mfrfull
	\end{tikzcd}
\]
\end{theorem}

The proof of the theorem occupies rest of the section. In Sec. \ref{}, we describe properads $\mfrfullunst$ and $\mfrfull$, homotopy equivalent to $\Mfrfullunst$ and $\Mfrfull$ respectively, which are more suitable for computing homotopy pushouts. In Subsections \ref{} we express the homotopy pushout in term of the pushouts of the bar resolutions and prove that the pushout is homotopy equivalent to $\mfrfullunst$ by using the local criteria for weak homotopy equivalences (see Proposition \ref{}). This last step is carried out by making suitable modification to the argument in Subsection \ref{finalhtpyeq}

%The rest of this section we provide the proof of the theorem. An outline of the proof follows: In Sec. \ref{}, we replace $\Mfrfullunst$ and $\Mfrfull$ with homotopy equivalent properads $\mfrfullunst$ and $\mfrfull$ respectively, which would be more suitable for our computation of the homotopy pushout. In Sec. \ref{} we describe the homotopy pushout as the $\Galg$ pushout of the cofibrant replacements of $\mfr,\mfrunst$, and $\mfrfull$. The cofibrant replacements $\mmfr, \mmfrunst, \mmfrfullunst$ are obtained using the monadic bar construction for the monad $\GG$ (see \ref{}). The cofibrant resolutions and their pushout will be modeled on `nested' directed acyclic graphs (da-graphs). Fianlly in Sec. \ref{}, we prove that the pushout is homotopy equivalent to $\mfrfullunst$ by using the local criteria for weak homotopy equivalences (see Proposition \ref{}).

\subsection{Properads $\mfrfullunst$, and $\mfrfull$}
First suppose that $2-2g-n_-n_+<0$.\\
Define $(\widetilde{\mfrfull}_{g,n_-,n_+})_{k}$ to be the moduli space of tuples $(\Sigma, \un{\gamma} , \un{t})$, where
\begin{itemize}
	\item $\Sigma \in \Mmod^{fr}_{g,n_-,n_+}$
	\item $\un{\gamma} \colon \coprod_{i=1}^{k} S^1 \to \Sigma$ is an analytic embedding disjoint from the boundary circles
	\item $\un{t} = (t_1,...,t_k) \in [0,1]^k$, such that $t_i=1$ for at least one $i$ whenever $n_+=0$
\end{itemize}
Define $(\mfrfull_{g,n_-,n_+})_{k}$ to be the quotient $(\widetilde{\mfrfull}_{g,n_-,n_+})_{k}  / \Sigma_k$, where $\Sigma_k$ acts on $(\widetilde{\mfrfull}_{g,n_-,n_+})_{k}$ by permuting components of the domain of $\un{\gamma}$ and by permuting the co-ordinates of $\un{t}$. $(\mfrfull_{g,n_-,n_+})_{k}$ can be thought of as moduli space of riemann surface with $k$ unlabeled, analytically embedded circles each carrying a weight in $[0,1]$, with at least one of the weights being 1 when the surface has no output. Following \cite{oancea2020deligne}, we call such a surface a `\textbf{\splitsurface}'.\\
Now define $\mfrfull_{g,n_-,n_+} = \coprod_{k \geq 0} (\mfrfull_{g,n_-,n_+})_{k}/ \sim$, where $\sim$ identifies a \splitsurface\  with the one obtained by forgetting the embedded circles with weight $0$. More precisely, $\mfrfull_{g,n_-,n_+}$ is obtained as the colimit 
\[ \mfrfull_{g,n_-,n_+} = \varinjlim \big[ F_0 \mfrfull_{g,n_-,n_+} \to F_1 \mfrfull_{g,n_-,n_+} \to \hdots \big] \] 
where the spaces $F_k \mfrfull_{g,n_-,n_+}$ are defined inductively satisfying the following properties:
\begin{enumerate}
	\item $F_0 \mfrfull_{g,n_-,n_+}  = (\mfrfull_{g,n_-,n_+})_{0}$ 
	\item $F_k \mfrfull_{g,n_-,n_+} $ admits a map from $\big[ \mfrfull_{g,n_-,n_+}(k+1) \big]^{-}$, where $\big[ \mfrfull_{g,n_-,n_+}(k+1) \big]^{-}$ is the subspace of $\mfrfull_{g,n_-,n_+}(k+1 )$ of {\splitsurface}s where at least one embedded circle has weight $0$.
	\item $F_{k+1} \mfrfull_{g,n_-,n_+} $ is obtained inductively from $F_k \mfrfull_{g,n_-,n_+} $ as the pushout
	\[ 
		\begin{tikzcd}\label{polkapush}
			\big[ (\widetilde{\mfrfull}_{g,n_-,n_+})_{k+1} \big]^{-} \ar{r}\ar{d} &  (\widetilde{\mfrfull}_{g,n_-,n_+})_{k+1} \ar{d}\\
			F_{k} \mfrfull_{g,n_-,n_+} \ar{r} &  F_{k+1} \mfrfull_{g,n_-,n_+}
		\end{tikzcd}
	\]
\end{enumerate}

\begin{notation}
\begin{enumerate}
	\item Image of a tuple $(\Sigma,\un{\gamma},\un{t})$ in $\mfrfull_{g,n_-,n_+}$ will be denoted by $[\Sigma,,\un{\gamma},\un{t}]$
	\item We fix notation for some surfaces which appear frequently in the later sections: An annulus with both ends outputs (respectively, inputs) and with a single positive weight embedded circle will be called an positive weight output (respectively, input) macaroni. Note that the embedded circle in a positive weight input macaroni necessarily has weight $1$.
%	\begin{enumerate}
%		\item Consider $(\mathbb{P}^1 \in \Mmod^{fr}_{0,1,0}, (0), (1))$, the \splitsurface with the underlying surface a disk with an input, and a single embedded circles at the origin carrying weight $1$. This will be called the full-weight cup and denoted $\mathbb{P}_-$.
%		\item Similarly, the \splitsurface $(\mathbb{P} \in \Mmod_{0,0,1}, (0), (1))$ will be called full-weight cap denoted $\mathbb{P}_+$. 
%	\end{enumerate}
\end{enumerate}
\end{notation}

\begin{proposition}\label{polkaproposition}
	The forgetful map $\mfrfull_{g,n_-,n_+} \to \Mmod^{fr}_{g,n_-,n_+}$ is a weak homotopy equivalence.
\end{proposition}
\begin{proof}
	When $n_+>0$, consider the section of $\mfrfull_{g,n_-,n_+} \to \Mmod^{fr}_{g,n_-,n_+}$ given by mapping a Riemann surface $\Sigma$ to a weighted split surface with underlying riemann surface $\Sigma$ and with no embedded circles. Then we have a homotopy contracting $\Mmod^{fr}_{g,n_-,n_+}$ to the image of this section given by shrinking the weights of all the embedded circles to $0$.\\
	Let's now consider the general case. Note that the map $\mfrfull_{g,n_-,n_+} \to \Mmod^{fr}_{g,n_-,n_+}$ is a fibration and hence it suffices to show that the fiber of this map is contractible. Fix a surface $\Sigma \in \Mmod^{fr}_{g,n_-,n_+}$. We show that there exists a filtration $F_1 \subset F_2 \subset ... \subset F_n \subset ... $ of the fiber over $\Sigma$, such that the inclusion of each $F_n$ inside the fiber is null-homotopic. 
	The filtration is defined as follows: Recall that since $(n_-,n_+) \neq (0,0)$, $\Sigma$ has at least one boundary component. Let $\gamma_0$ be the $S^1$-parametrization of this component. Choose a conformal embedding $\widetilde{\gamma}_0 \colon [0,\epsilon) \times S^1 \to \Sigma$ extending $\gamma_0$. For $n \in \NN$ large enough define $F_n$ to be the subspace of the fiber consisting of {\splitsurface}s with all embedded circles with weight $\geq 1/2$ contained inside $\Sigma \setminus \widetilde{\gamma}_0([0,\frac{1}{n}] \times S^1)$. $F_n$ gives an increasing sequence of open subsets of the fiber. Moreover, since in any \splitsurface\  all the embedded circles lie away from the boundaries, it follows that $F_n$ is also exhausting.\\
	We now prove that each $F_n$ is contractible \emph{inside} the fiber over $\Sigma$. Consider the embedded circle $\gamma = \widetilde{\gamma}_0|_{\{\frac{1}{2n}\} \times S^1}$. Consider the \splitsurface\  in the fiber over $\Sigma$ which has a single weight $1$ embedded circle given by $\gamma$ viz $[\Sigma,\gamma,(1)]$. We show that $F_n$ can be contracted to this {\splitsurface}  \emph{inside the fiber}.\\
	The contraction homotopy is given by concatenation of $3$ homotopies $H_1,H_2,H_3$:
	\begin{enumerate}
		\item $H_1$: Consider a function $\tau \colon [0,1] \to [0,1]$ which is identically $0$ on $[0,1/2]$ and identically $1$ on $[3/4,1]$. Let $f_{1}$ be the map $[(\Sigma,\un{\gamma}, (t_1,...,t_k))] \mapsto [ (\Sigma,\un{\gamma}, ( \tau(t_1),...,\tau(t_k))) ]$. Then $H_1$ is a homotopy from identity to $f_{1}$. The homotopy is simply given by linearly interpolating weights of the embedded circles.
		\item $H_2$: Now consider map $f_{2}$ which takes any surface $[(\Sigma,\un{\gamma},\un{t})]$ in $F_n$ to the surface obtained by adding the embedded circle $\gamma = \widetilde{\gamma}_0|_{\{\frac{1}{2n}\} \times S^1}$ to $f_1([\Sigma,\un{\gamma},\un{t}])$. (In particular, note that $f_2$ takes values outside $F_n$.) $H_2$ is a homotopy between $f_1$ and $f_2$, given by interpolating the weight of $\gamma_0$ from $0$ to $1$. The well-definedness (and continuity) of $f_2$ and $H_2$ follows from the restriction in the definition of the subset $F_n$ that any embedded circle intersecting $\widetilde{\gamma}_0([0,\frac{1}{n}] \times S^1)$ has weight $< 1/2$
		\item $H_3$: Finally, $H_3$ is a homotopy from $f_2$ to the constant map $f_3$ with value $[\Sigma,\gamma,(1)]$. The homotopy $H_3$ is given by decreasing the weights of embedded circles outside $\widetilde{\gamma}_0([0,\frac{1}{n}) \times S^1)$
	\end{enumerate}
\end{proof}	
	

The properads $\mfrfullunst$ and $\mfrfull$ are now defined as follows:
\begin{itemize}
	\item $\mfrfull$ is defined in a manner similar to $\Mfrfull$ with the space $\Mmod^{fr}_{g,n_-,n_+}$ replaced by the corresponding spaces $\mfrfull_{g,n_-,n_+}$.
	
	\item $\mfrfullunst$ is the subproperad of $\mfrfull$ which has operations only in the unstable range :
	\[\mfrfullunst(1,1) =\mfrfullunst_{0,1,1}, \mfrfullunst(0,2) = \Mmod^{fr}_{0,0,2},\mfrfullunst(2,0)=\mfrfullunst_{0,2,0} \]
\end{itemize}		



The properad compositions are given by the maps induced by the gluing of riemann surfaces underlying the {\splitsurface}s.

The following is a consequence of Proposition \ref{polkaproposition}:
\begin{corollary}
The maps $\mfrfullunst \to \Mfrfullunst$, and $\mfrfull \to \Mfrfull$ are weak homotopy equivalences.
\end{corollary}

%except for the following changes:\\
%	\begin{enumerate}
%	\item For $\mfrunst$, 
%	\begin{itemize}
%		\item $\mfrunst(1,1), \mfrunst(0,2)$ consist is a single point in $\mfr_{0,1,1}, \mfr_{0,1,1}$ resp. corresponding to surfaces with no embedded disks.
%		\item $\mfunst(1,0)$ is a single point corresponding to the surface with disk
%	\end{enumerate}
%\begin{enumerate}


\subsection{Homotopy pushout of $\mfr \leftarrow \mfrunst \to \mfrfullunst$}

\textcolor{red}{$\redTopSeq$-cofibrancy for $\mfr,...$}\\
Using results from Section \ref{bar_const_red_sec}, that $|B_{\bullet}\mfr |, |B_{\bullet} \mfrunst| $, and $|B_{\bullet}\mfrfullunst |$ are cofibrant replacements of $\mfr, \mfrunst,$ and $\mfrfullunst$ respectively and that maps $|B_{\bullet} \mfrunst | \to |B_{\bullet}\mfr|$ and $|B_{\bullet} \mfrunst | \to |B_{\bullet}\mfrfullunst|$ are cofibrations of reduced properads. Thus the homotopy pushout of the diagram $\mfr \leftarrow \mfrunst \to \mfrfullunst$ is given by,
\[ | B_{\bullet} \mfr | \coprod_{|B_{\bullet} \mmfrunst |}^{\Galg}|B_{\bullet}\mmfrfullunst| \]
As in the proof of Proposition \ref{compareproposition}, we have
\[ | B_{\bullet} \mfr | \coprod_{|B_{\bullet} \mfrunst |}^{\Galg}|B_{\bullet}\mfrfullunst| = \big| \GG( \GG^{\bullet} \mfr \sqcup_{\GG^{\bullet} \mfrunst} \GG^{\bullet} \mfrfullunst) \big| \]

The following then completes the proof of Theorem \ref{mainthm2}:

\subsection{$\pi \colon \big| \GG\big( \GG^{\bullet} \mfr \sqcup_{\GG^{\bullet} \mfrunst} \GG^{\bullet} \mfrfullunst\big) \big| \to \mfrfull$ is a $\Galg$-weak equivalence}
	The strategy of proof is similar to that of Subsection \ref{finalhtpyeq}. In this case instead of cuts made around nodal points, we will construct cuts around the positive weight embedded circles. This gives a canonical way of decomposing an element of $\mfrfull$ into a \splitsurface \ in $\mfr$ with no positive-weight embedded circles and a collection positive weight macaronis. As in Subsection \ref{finalhtpyeq}, some care will be needed to extend this decomposition continuously to a neighborhood of the \splitsurface. Cut homotopies, analogous to those in Subsection \ref{finalhtpyeq}, will then be constructed to homotope the points in the pushout to the zero simplices corresponding to such decompositions.
%	We start by describing map 
%	\[\phi \colon \mfrfullunst \to \big| \GG\big( \GG^{\bullet} \mfr \sqcup_{\GG^{\bullet} \mfrunst} \GG^{\bullet} \mfrfullunst\big) \big|\]
%	 Below we will show that $\chi$ is a homotopy inverse of $\pi$. 
%	 The image of $\chi$ is contained inside the $0$-simplices of $\big| \GG\big( \GG^{\bullet} \mfr \sqcup_{\GG^{\bullet} \mfrunst} \GG^{\bullet} \mfrfullunst\big) \big|$. Consider the point of $\mfrfullunst$ represented by a \splitsurface $(\Sigma,\phi,\un{t})$. It is mapped under $\chi$ to the point in $\GG \big(  \mfr \sqcup_{\mfrunst} \mfrfullunst \big)$ 
%	 \splitsurface represented by scaling each of the embedded disk via the map $z \mapsto z/2$. More precisely, it is mapped to the  point represented by $(\Sigma,\phi \circ s, \un{t})$, where $s \colon \coprod \DD \to \coprod \DD$ is the map which on each disk acts by $z \mapsto z/2$.

%To highlight the parallels between the formal structure of the argument here and in Subsec \ref{finalhtpyeq}, we will use the same notation as in Subsection \ref{} for the analogous notions here. As we shall see, having suitably adopted various definitions to the current context, many of the maps and homotopies will be given by the same expressions as in Subsection \ref{finalhtpyeq}.

%Let us first restrict to the case $(n_-,n_+) \neq (0,0)$. The case $(n_-,n_+)=(0,0)$ will require some additional argument having to do with the fact that $\mfrfull_{g,n_-,n_+}$ is a topological stack rather than a topological space in this case.
%
%\textbf{{\large Case I $(n_-,n_+) \neq (0,0)$}}
%
%%As before, we will a cover $\mfrfull_{g,n_-,n_+}$ a neighborhood $U_{\un{\Sigma}}$ such that for any finite collection $U_{\un{\Sigma_1}}, ..., U_{\un{\Sigma_n}}$ in $\mfrfull_{g,n_-,n_+}$, 
%%\[\pi^{-1}|_{U_{\un{\Sigma_1}}\cap ... \cap U_{\un{\Sigma_n}}} \colon \pi^{-1}(U_{\un{\Sigma_1}}\cap ... \cap  U_{\un{\Sigma_n}}) \to U_{\un{\Sigma_1}} \cap ... \cap U_{\un{\Sigma_n}}\]
%%is a weak-equivalence.
%
%As before, we will find a cover $\{U_{\lambda}\}$ of $\mfrfull_{g,n_-,n_+}$ such that for any finite collection $U_{\lambda_1}, ..., U_{\lambda_n}$ in $\mfrfull_{g,n_-,n_+}$, 
%\[\pi^{-1}|_{U_{\lambda_1,...,\lambda_k}} \colon \pi^{-1}(U_{\lambda_1,...,\lambda_k}) \to U_{\lambda_1}, ..., U_{\lambda_n}\]
%is a weak-equivalence, where $U_{\lambda_1,...,\lambda_k}:= U_{\lambda_1} \cap ... \cap U_{\lambda_n}$.
%
%Moreover, each $U_{\lambda_1,...,\lambda_k}$ has a filtration 
%\[ F^1\pi^{-1}(U_{\lambda_1,...,\lambda_k}) \subset \hdots \subset F^n\pi^{-1}(U_{\lambda_1,...,\lambda_k}) \subset \hdots \pi^{-1}(U_{\lambda_1,...,\lambda_k})\]
%such that 
%\begin{equation}\label{filtrfulltcft}
%	\begin{split}
%		 &\mbox{each filtered piece } F^n\pi^{-1}(U_{\lambda_1,...,\lambda_k}) \mbox{ contracts, \emph{inside} } \pi^{-1}(U_{\lambda_1,...,\lambda_k}),\\[-4pt]
%		 & \mbox{ onto a section } \chi_n \colon U_{\lambda_1,...,\lambda_k} \to F^n\pi^{-1}(U_{\lambda_1,...,\lambda_k})
%	\end{split}
%\end{equation}
%
%
%\subsubsection{Construction of $U_{\lambda}$ and $F^n \pi^{-1}(U_{\lambda_1,...,\lambda_n})$}
%We start by describing the open subsets $U_{\lambda}$ for a $\un{\Sigma} \in \mfrfull_{g,n_-,n_+}(n)$. Let $U_{\alpha, n}$ be the open subset of $\mfrfull_{g,n_-,n_+}$ consisting of \splitsurface s containing exactly $n$ embedded circles with weight $>\alpha$ and no embedded circle with weight $=\alpha$. $Z_{\alpha,i} \subset U_{\alpha,i}$ be the subspace of \splitsurface s containing $n$ positive weight circles all of which have weight $1$. There exists a deformation retraction $r_{\alpha} \colon U_{\alpha,i} \to Z_{\alpha,i}$ which homotopes weights of weight $<\alpha$ embedded circles to $0$ and the weights of weight $>\alpha$ embedded circles to $1$. Also, we have a covering map $\Mmod^{fr}_{n_-,n_+}(n) \to Z_{\alpha,i}$. The fiber over a \splitsurface in $Z_{\alpha,i}$ is given by different ways of labeling the embedded circles. Consider an open cover $\{V_{\alpha,i}^{\lambda}\}$ of $Z_{\alpha,i}$ such that there exist open sets $W_{\alpha,i}^{\lambda}$ of $\Mmod^{fr}_{n_-,n_+}(n)$ satisfying the following conditions:
%\begin{itemize}
%	\item the covering map maps $W_{\alpha,i}^{\lambda}$ homeomorphically onto $Z_{\alpha,i}^{\lambda}$. Let $\phi_{\alpha,i}^{\lambda}$ denote the inverse.
%	\item there exists a choice of continuous charts around embedded circles over $W_{\alpha,i}^{\lambda}$ in the following sense: Let $\mathcal{C}_{g,n_-,n_+}(n) \to \mfrfull_{g,n_-,n_+}(n)$ denote the universal curve. Then, there exist maps
%	there exists a maps 
%%	\[ 
%%		\begin{tikzcd}
%%			\sigma_i \colon W_{\alpha,n,\lambda} \times \DD \ar[rr]\ar[rd] && \mathcal{C}_{g,n_-,n_+}(n) \ar[ld]\\
%%			& W_{\alpha,n,\lambda} & 
%%		\end{tikzcd}
%%	\]
%	\[ \sigma_{\alpha,i}^{\lambda,i} \colon W_{\alpha,i}^{\lambda} \times \sqcup_{n} \DD \to \mathcal{C}_{g,n_-,n_+}(n)|_{W_{\alpha,i}^{\lambda}} \]
%	over $W_{\alpha,i}^{\lambda}$, such each component of $\sigma_{\alpha,i}^{\lambda,i}$ is a fiberwise chart taking $0$ to the $i$th embedded circle. Moreover we assume that the charts given by $\sigma_{\alpha,i}^{\lambda,i}$ are pairwise disjoint.	
%\end{itemize}

%Define $U_{\alpha,i}^{\lambda} = r^{-1} ( W_{\alpha,i}^{\lambda} )$. $U_{\alpha,i}^{\lambda}, \lambda \in \Lambda_{\alpha,i}$ is a cover of $U_{\alpha,i}$. Putting all the $\{U_{\alpha,i}^{\lambda}\}$'s together as $\alpha, n , \lambda$ vary in $(0,1), \ZZ_{\geq 0}, \Lambda_{\alpha,i}$.
%
%Let $\alpha \in (0,1)$ be such that $\un{\Sigma}$ has no embedded circle with weight $\alpha$. Let $\un{\gamma}_{> \alpha}$ denote the subset of embedded circles with weight $>\alpha$. For every $\gamma \in \un{\gamma}_{> \alpha}$ select an open neighborhood $V_{\gamma}$ around $\gamma$ in $\Sigma$, so that $V_{\gamma}$ are disjoint from each other. Define $U_{\Sigma}$ to be the open subset of $\mfrfull_{g,n_-,n_+}$ given by points represented by weighted marked surfaces with exactly one point with weight $>\alpha$ in $V_{\gamma}$ for every $\gamma \in \un{\gamma}_{> \alpha}$ and having no embedded circles with weight $\alpha$. Note that on $U_{\Sigma}$ there are well defined hodge bundles $\phi_{\gamma}$ with fiber at a point give by the tangent space to the underlying riemann surface at unique embedded circle in $V_{\gamma}$ with weight $> \alpha$. Shrinking $U_{\Sigma}$ by restricting the underlying riemann surfaces to a neighborhood of $\Sigma$ if necessary, we may assume that the line bundles $\phi_{\gamma}, \gamma \in \un{\gamma}_{> \alpha}$ admit non-vanishing sections over $U_{\Sigma}$. Moreover, we may assume that these sections are independent of the weights of the embedded circle points provided they are larger than $\alpha$. Moreover $V_{\gamma}$ are disjoint implies that over $U_{\un{\Sigma}}$ we have sections $\gamma$ giving the embedded circles over the corresponding surface \textcolor{red}{ (a slice of the $\Sigma_{\gamma}$-action. Orbifold points ?)}
%

%Choose $U_{\un{\Sigma}}$ to be a neighborhood of $\un{\Sigma}$ satisfying the following conditions:
%\begin{itemize}
%\item  $U_{\un{\Sigma}}$ contains no surface with a embedded circle of weight $\alpha$
%\item  The hodge bundles corresponding to embedded circles with weight $>\alpha$ admit non-vanishing sections over this neighborhood. Fix such a section. We may further assume that the section is independent of the choice of the weights as long as the correponding
%\end{itemize}

To highlight the parallels between the formal structure of the argument here and in Subsec \ref{finalhtpyeq}, we will use the same notation as in Subsection \ref{} for the analogous notions here. As we shall see, having suitably adopted various definitions to the current context, many of the maps and homotopies will be given by the same expressions as in Subsection \ref{finalhtpyeq}.

As before, we will find a cover $\{U_{\lambda}\}_{\lambda}$ of $\mfrfull_{g,n_-,n_+}$ such that for any finite collection $U_{\lambda_1}, ..., U_{\lambda_k}$ in $\mfrfull_{g,n_-,n_+}$, 
\[\pi^{-1}|_{U_{\lambda_1,...,\lambda_k}} \colon \pi^{-1}(U_{\lambda_1,...,\lambda_k}) \to U_{\lambda_1}, ..., U_{\lambda_k}\]
is a weak-equivalence, where $U_{\lambda_1,...,\lambda_k}:= U_{\lambda_1} \cap ... \cap U_{\lambda_k}$.

Moreover, each $U_{\lambda_1,...,\lambda_k}$ has a filtration 
\[ F^1\pi^{-1}(U_{\lambda_1,...,\lambda_k}) \subset \hdots \subset F^n\pi^{-1}(U_{\lambda_1,...,\lambda_k}) \subset \hdots \pi^{-1}(U_{\lambda_1,...,\lambda_k})\]
such that 
\begin{equation}\label{filtrfulltcft}
	\begin{split}
		 &\mbox{each filtered piece } F^n\pi^{-1}(U_{\lambda_1,...,\lambda_k}) \mbox{ has a fiberwise contraction, \emph{inside} } \pi^{-1}(U_{\lambda_1,...,\lambda_k}),\\[-4pt]
		 & \mbox{ onto a section } \chi_n \colon U_{\lambda_1,...,\lambda_k} \to F^n\pi^{-1}(U_{\lambda_1,...,\lambda_k})
	\end{split}
\end{equation}

\begin{notation}
Denote the simplicial complex underlying $\pi^{-1}(U_{\lambda_1 ...,\lambda_k})$ by $\PP_{\lambda_1,...,\lambda_k}$. The simplicial complex underlying $F^n\pi^{-1}(U_{\lambda_1 ...,\lambda_k})$ is denoted $F^n\PP_{\lambda_1,...,\lambda_k}$
\end{notation}

\subsubsection{Construction of $U_{\lambda}$ and $F^n \PP_{\lambda_1,...,\lambda_k}$}
We start by describing the spaces $U_{\lambda}$. Let $U_{\alpha, n}$ be the open subset of $\mfrfull_{g,n_-,n_+}$ consisting of \splitsurface s containing exactly $i$ embedded circles with weight $>\alpha$ and no embedded circle with weight $=\alpha$. $Z_{\alpha,i} \subset U_{\alpha,i}$ be the subspace of \splitsurface s containing $i$ positive weight circles all of which have weight $1$. There exists a deformation retraction $r_{\alpha} \colon U_{\alpha,i} \to Z_{\alpha,i}$ which homotopes weights of weight $<\alpha$ embedded circles to $0$ and the weights of weight $>\alpha$ embedded circles to $1$. Also, we have a covering map $\Mmod^{fr}_{n_-,n_+}(i) \to Z_{\alpha,i}$. The fiber over a \splitsurface \ in $Z_{\alpha,i}$ is given by different ways of labeling the embedded circles. Consider a covering $\{W^{\lambda}_{\alpha,i}\}
_{\lambda \in \Lambda_{\alpha,i}}$  of $\mfrfull_{g,n_-,n_+}$, where $ \Lambda_{\alpha,i}$ is some indexing set, satisfying the following conditions:
	\begin{itemize}
		\item $W^{\lambda}_{\alpha,i}$ maps homeomorphically onto an open subset of $Z_{\alpha,i}$.
		\item \begin{equation} 
		\begin{minipage}{0.9\textwidth}
	 There exists a continuous choices of embedded annuli around the embedded circles over $W_{\alpha,i}^{\lambda}$ in the following sense: Let $\mathcal{C}_{g,n_-,n_+}(i) \to \mfrfull_{g,n_-,n_+}(i)$ denote the universal curve. Then, there exists $r^{\lambda}_{\alpha,i}>0$ and maps
	\[ \sigma_{\alpha, i}^{\lambda} \colon W_{\alpha,i}^{\lambda} \times \sqcup_{j=1}^i \{1/r^{\lambda}_{\alpha,i} \leq |z| \leq r^{\lambda}_{\alpha,i}\} \to \mathcal{C}_{g,n_-,n_+}(n)|_{W_{\alpha,i}^{\lambda}} \]
	over $W_{\alpha,i}^{\lambda}$, such $\sigma_{\alpha,i}^{\lambda}$ is a fiberwise embedding taking $\{|z|=1\}$ in the $i$th annulus to the $i$th embedded circle. Moreover we assume that the embeddings given by $\sigma_{\alpha,i}^{\lambda}$ are disjoint from the boundaries.
		\end{minipage}
	\end{equation}
	\end{itemize}

%Note that $W = \coprod_{\lambda} W_{\alpha,i}^{\lambda} \to Z_{\alpha,i}$ defines an atlas for $Z_{\alpha,i}$ as well.

Let $\widetilde{W}^{\lambda}_{\alpha,i}$ be the image of $W^{\lambda}_{\alpha,i}$ in $Z_{\alpha,i}$. $\{\widetilde{W}^{\lambda}_{\alpha,i}\}$ forms a cover of $Z_{\alpha,i}$. Define $U_{\alpha,i}^{\lambda} = r^{-1} (\widetilde{W}_{\alpha,i}^{\lambda}) $, where $r \colon U_{\alpha,i} \to Z_{\alpha,i}$ is the retraction mentioned above. $U_{\alpha,i}^{\lambda}, \lambda \in \Lambda_{\alpha,i}$ gives a covering of $U_{\alpha,i}$. Putting all the $\{U_{\alpha,i}^{\lambda}\}$'s together as $\alpha, n ,$ and $\lambda$ vary in $(0,1), \ZZ_{\geq 0}$, and $\Lambda_{\alpha,i}$ respectively, we obtained the desired cover $\{U_{\lambda}\}$ of $\mfrfull_{g,n_-,n_+}$.

 
For $\lambda_1,...,\lambda_k \in \Lambda$, we now turn to showing that $\pi^{-1}(U_{\lambda_1, ..., \lambda_n}) \to U_{\lambda_1, ..., \lambda_n}$ is a weak equivalence. Let $\alpha_1, ... , \alpha_k$ be the constants for $\lambda_1,...,\lambda_k$ as above. Without loss of generality, assume that $0 = \alpha_{0} < \alpha_{1} \leq ... \leq \alpha_{k} < \alpha_{k+1} = 1$. From the definitions of the sets $U_{\lambda_i}$ it is clear that for any surface in the intersection $U_{\lambda_1,...,\lambda_k}$ the number of embedded circles with weight in $>\alpha_r$ stays constant for every $r$, viz. $i_r$.\\ %Denote these by $\gamma_{1,k},...,\gamma_{n_k,k}$. \\


%
%For a point $\un{\Sigma}^{\prime}$ in $U_{\un{\Sigma}}$, denote by $\gamma(\un{\Sigma}^{\prime})$ the unique embedded circle in $B_{\gamma}$ with weight $>1/2$. $\gamma(\un{\Sigma}^{\prime})$ varies continuously as $\un{\Sigma}$ varies in $U_{{\lambda_1,...,\lambda_n}}$. Note that $U_{\un{\Sigma}}$ is contractible. Hence we can find disc embeddings $B_{\gamma}^1$ contained inside $B_{\gamma}$ centered at the unique embedded circle in $B_{\gamma}$ with weight $>1/2$ varying continuously over $U_{{\lambda_1,...,\lambda_n}}$. By scaling these embeddings we can obtained a sequence of (continuously varying) sequence of disc embeddings $B_{\gamma}^i(\un{\Sigma}^{\prime})$ with centers at $\gamma(\un{\Sigma}^{\prime})$ and radii going to $0$.
We now construct filtration $F^n\PP_{\lambda_1,...,\lambda_k}$ satisfying \eqref{filtrfulltcft}.
The spaces $F^n \PP_{\lambda_1,...,\lambda_k}$ are defined as follows:\\
Consider the fiberwise embeddings 
\[(\sigma_{\alpha_{1},i_{1}}^{\lambda_{1}})^n \colon W_{\alpha_{1},i_{1}}^{\lambda_{1}} \times  \{1/(r^{\lambda})^{1/n}_{\alpha,i} \leq |z| \leq (r^{\lambda}_{\alpha,i})^{1/n}\}  \to {\mathcal{C}_{g,n_-,n_+}(i)}|_{W_{\alpha,i}^{\lambda}} \]
obtained by restricting the maps $\sigma_{\alpha_{1},i_{1}}^{\lambda_{1}}$ from \ref{}. Let $\phi^{\lambda_1}_{\alpha_{1},i_{1}} \colon U_{\alpha_{1},i_{1}}^{\lambda_{1}} \to W_{\alpha_{1},i_{1}}^{\lambda_{1}}$ denote the inverse of the homeomorphism $W_{\alpha_{1},i_{1}}^{\lambda_{1}} \to U_{\alpha_{1},i_{1}}^{\lambda_{1}}$. Pulling back the maps $(\sigma_{\alpha_{1},i_{1}}^{\lambda_{1}})^n$ by $\phi^{\lambda_1}_{\alpha_{1},i_{1}}$ we obtain
\[ (\phi_{\alpha_{1},i_{1}}^{\lambda_{1}})^*(\sigma_{\alpha_{1},i_{1}}^{\lambda_{1}})^n \colon U_{\alpha_{1},i_{1}}^{\lambda_{1}} \times  \{1/(r^{\lambda})^{1/n}_{\alpha,i} \leq |z| \leq (r^{\lambda}_{\alpha,i})^{1/n}\}  \to (\phi_{\alpha_{1},i_{1}}^{\lambda_{1}})^* {\mathcal{C}_{g,n_-,n_+}(n)}|_{W_{\alpha,i}^{\lambda}} \]
Denote the image of this map by $A^n_j$.\\
Analogous to Subsubsection \ref{gammasubsubsect}, given an $\un{x} \in   \GG^{n+1}\mfr \coprod_{\GG^{n+1}\mfrunst} \GG^{n+1}\mfrfullunst$, with underlying da-graph $G$, we can define
\[ \gamma_{\un{x}} \colon \coprod_{E(G) \cup in(G) \cup out(G)} S^1 \to \GG^{n+1}\mfr \coprod_{\GG^{n+1}\mfrunst} \GG^{n+1}\mfrfullunst \to \Sigma_{\un{x}}\]
Let $\mathbb{V}_{\lambda_1,...,\lambda_k}$ denote the subcomplex of $\GG^{\bullet}\mfr \coprod_{\GG^{\bullet}\mfrunst} \GG^{\bullet}\mfrfullunst$ consisting of simplices $\un{x}$ such that the corresponding embedding $\gamma(\un{x})$ is disjoint from $\bigcup_{j} A^n_j$. We then set $F^n \PP_{\lambda_1,...,\lambda_k}$ to be the open subset of $|\GG^{n+1}\mfr \coprod_{\GG^{n+1}\mfrunst} \GG^{n+1}\mfrfullunst|$ corresponding to this subcomplex.

\subsubsection{Construction of sections $\chi_n \colon U_{\lambda_1,...,\lambda_k} \to F^n \PP_{\lambda_1,...,\lambda_k}$}
We now construct sections $\chi_n$ as in \ref{}. The image of the map lies in the subspace of $0$-simplices in $\pi^{-1}U_{\lambda_1,...,\lambda_k}$.

%Let $\un{\Sigma}$ be a point in $U_{{\lambda_1,...,\lambda_n}}$. Consider the da-graph $G_{\un{\Sigma}}$ defined as follows: the set of vertices is given by the set connected components of $\un{\Sigma} \setminus \partial \bigcup_{\un{\gamma}_{=1}} B^i_{\gamma}(\un{\Sigma})$. For every $\gamma$ in $\un{\gamma}_{=1}$, there is a directed edge from the vertex corresponding to the connected component $\un{\Sigma} \setminus \bigcup_{\un{\gamma}_{=1}} B^i_{\gamma}(\un{\Sigma})$ to the one corresponding to $B^i_{\gamma}(\un{\Sigma})$. Now we define a point $\un{\Sigma}_{\GG}$ in $\GG \big( \mfr \sqcup_{\mfrunst} \mfrfullunst \big)$ with the underlying da-graph $G_{\un{\Sigma}}$: the connected component $\Sigma^{\prime} \setminus \bigcup_{\un{\gamma}_{=1}} B^i_{\gamma}(\un{\Sigma})$ with boundaries corresponding to discs $B^i_{\gamma}(\un{\Sigma}), \gamma \in \un{\gamma}_{=1},$ treated as outputs, defines an element in $\mfrfull_{g,n_-,n_+ + n}$, having no weight $1$-embedded circles. We label the vertex given by $\Sigma \setminus \bigcup_{\un{\gamma}_{=1}} B^i_{\gamma}(\un{\Sigma})$ by this label. The vertices corresponding to $\gamma \in \un{\gamma}_{=1}$ are labeled by discs $B^i_{\gamma}(\un{\Sigma_i})$ with boundary treated as output, and with a single weight $1$ embedded circle at $\gamma(\un{\Sigma_i})$.\\ 
% Set $\chi_{n}(\un{\Sigma})= x_{\un{\Sigma}}$. 


Consider an $(n_-,n_+)$- da-graph $G_{\lambda_1,...,\lambda_k}$ defined as follows: $G_{\lambda_1,...,\lambda_k}$ has $(n_1+1)$-vertices labeled by $0,1,...,n_1$ and $2n_1$ edges, two for every $1\leq i \leq n_1$, both directed from $0$ to $i$. It has $n_-$ input and $n_+$ output legs, all of which are attached to the vertex $0$. The image of $\chi_n$ consist of $0$-simplices whose labels have underlying graphs $G_{\lambda_1,...,\lambda_k}	$. Let $\un{\Sigma} \in U_{\lambda_1,...,\lambda_k}$ and $\Sigma$ the Riemann surface underlying $\un{\Sigma}$. $\Sigma$ is identified with the fiber over $\un{\Sigma}$ of the pullback of $\mathcal{C}_{g,n_-,n_+}(n_1)$ to $U_{\lambda_1,...,\lambda_k}$. Let $A_j^n(\un{\Sigma}) \subset \Sigma$ be the image of the restriction of the (pullback of) the map $(\sigma_{\lambda_1,n_1}^{\lambda_1})^n$ to the fiber. Let $w_j(\Sigma)$ denote the weight of the embedded circle $i$th, $0 \leq j \leq i_1$. 
Set 
\[ \un{A^n_j}(\un{\Sigma}) :=
	\begin{cases}
		[A^n_j,\gamma_i^N,\rho(w_j)] \in \mfrfull_{0,0,2} &\mbox{ for } w_j < \alpha_n \\ 
		[A^n_j,\gamma_i^N,\rho(w_j)=1] \in \mfrfull_{0,2,0} &\mbox{ for } w_j > \alpha_n 
	\end{cases}
\]
where $\rho$ is the step function defined on $[0,1] \setminus \{\alpha_1,...,\alpha_k\}$, given by:
\begin{equation}\label{rhofunction} \rho(x) := \begin{cases}
				0 & \mbox{ for } x \in (\alpha_0,\alpha_1) \\
				\frac{\alpha_{r} + \alpha_{r+1}}{2} &\mbox{ for } x \in (\alpha_{r}, \alpha_{r+1}), 0<r<n \\ 
				1 & \mbox{ for } x \in (\alpha_n,\alpha_{n+1}) 
			\end{cases} 
\end{equation}
(the parametrization of the boundary $\partial A^n_j$ is given by the embeddings $\sigma^{\lambda_1}_{\alpha_1,n_1}$).

Now, note that the number of embedded circles with weight $> \alpha_k$ is $i_k$, and those with weight $< \alpha_k$ is $i-i_k$.\\
Consider the surface $\Sigma \setminus \bigcup_j A^n_j(\Sigma)$. Treating the boundary components $\partial A^n_j(\un{\Sigma})$ corresponding to $j$ with $w_j> \alpha$ as outputs and those with $w_j < \alpha$ as inputs, we get a surface in $\Mmod^{fr}_{g,n_-+2(i-i_1), n_+ + i_1}$. Define a \splitsurface \ $ \un{\Sigma} \setminus \bigcup_j \un{A^n_j}(\un{\Sigma}) \in \mfrfull_{g,n_-+2(i-i_1), n_+ + i_1}$ with this as the underlying surface and with no embedded circles with positive weight. \\
Thus, for any $\un{\Sigma}$ we have an element $G_{\lambda_1,...,\lambda_k}(\un{\Sigma})$ of $\big(\GG^{n}\mfr \coprod_{\GG^{n}\mfrunst} \GG^{n}\mfrfullunst \big)$ with the underlying da-graph $G_{\lambda_1,...,\lambda_k}$ where the label of $0$ is given by $\un{\Sigma} \setminus \bigcup_j \un{A^n_j}(\un{\Sigma}) (\un{\Sigma}) $ and the label of vertex $j$ is given by $\un{A^n_j}(\un{\Sigma})$.

The map $\chi_n \colon U_{\lambda_1,...,\lambda_k} \to \PP_{\lambda_1,...,\lambda_k}$ is then defined by 
\[\Sigma \mapsto G(\un{\Sigma}) \]




\subsubsection{Verifying condition \eqref{filtrfulltcft}}
We construct on $\PP_{\lambda_1,...,\lambda_k}$ a homotopy between the identity and $\chi_{n} \circ \pi$. We start by describing notions analogous to those in \ref{prfcondfiltrtop}.

	\paragraph{\textbf{The graphs $K_v$}}
	This is analogous to graphs $H_v$ defined in \ref{}.\\
	If $\un{x}$ is an element in $\mathbb{V}_{\lambda_1,...,\lambda_k}$ with underlying graph $G$, and $\un{x}_v$ be the label of $v \in V(G)$ in $\un{x}$, Then define $K_v$ to be the graph with:
		\begin{itemize}
			\item the set of vertics given by the connected components of $\pi(\un{x}_v) \setminus \bigcup_{n} \partial A^n_j (\pi(\un{x}))$ 
			\item the edges are given by connected components of $\bigcup_{n} \partial A^n_j (\pi(\un{x}))$.  
		\end{itemize}
		
	\paragraph{\textbf{Cut map $c \colon \PP_{\lambda_1,...,\lambda_k} \to \PP_{\lambda_1,...,\lambda_k}$}}
	As before $c$ is defined as the geometric realization of a map of simplcial spaces given by $c_r \colon \PP^r_{\lambda_1,...,\lambda_n} \to\PP^r_{\lambda_1,...,\lambda_n}$. Maps $c_r$ are defined inductively as follows:
		\begin{itemize}
			\item $r=0$: Let $\un{x} \in \PP^0_{\lambda_1,...,\lambda_n}$ be an element with underlying graph $G$. Then $c_0(\un{x})$ is defined by the same expression as in \ref{c0eq}, where now $\un{y}_v, v \in V(G)$ are defined as follows: let $\un{x}_v$ denote the label of $v$ in $\un{x}$. Then $\un{y}_v$ is defined to be the element of $\PP^0_{\lambda_1,...,\lambda_n}$ with underlying da-graph $K_v$ as above and with the labels of vertices in $K_v$ given by the corresponding connected components of $\un{x}\setminus \partial \bigcup_i A^n_j$.
			\item The inductive step constructing map $c_r$ assuming the existence of $c_i$ for $i <r$ is defined exactly as in \ref{}, with $\PP^i_{\Sigma}, H_v$, and $A_\gamma$ there replaced by $\PP^i_{\lambda_1,...,\lambda_k}, K_v$, and $A^n_j(\un{\Sigma})$ respectively.
		\end{itemize}
	\paragraph{\textbf{Cut homotopies $\Phi \colon F^n \PP_{\lambda_1,...,\lambda_k}\to U_{\lambda_1,...,\lambda_k}$}}
	We now construct a homotopy from identity to the cut map $c$. The homotopy is a concatenation of $2$ homotopies.\\
\begin{itemize}
	\item $\Phi_{\rho}$: Recall the step function $\rho$ from \ref{rhofunction}. We have an induced map $f_{\rho} \colon \PP_{\lambda_1,...,\lambda_k}\to \PP_{\lambda_1,...,\lambda_k}$ given by the geometric realization of the map of simplicial spaces which changes the label of a simplex by applying $\rho$ to the weights of the embedded circles in the label. $\rho$ is continuous on its domain and thus the induced map $f_{\rho} \colon \PP_{\lambda_1,...,\lambda_k}\to \PP_{\lambda_1,...,\lambda_k}$ is continuous as well. Let $\Phi_{\rho}$ be the homotopy from the identity to $f_{\rho}$ obtained by linearly interpolating the weights of the embedded circles in the labels.
%	In particular the image of $\Phi_1$ lies in the fibers of $\pi$ above the \splitsurface s in $U_{n}$ which contained a unique weight $1$ embedded circle at $\gamma(\un{\Sigma}), \gamma \in \gamma_{=1}(\un{\Sigma})$ and no other point with positive weight.
	\item Next we have the cut homotopy $\Phi$. This is defined as a geometric realization of a simplicial homotopy $\Phi \colon \PP_{{\lambda_1,...,\lambda_k}} \to \PP_{{\lambda_1,...,\lambda_k}}$. This in turn is given by maps $\Phi^j \colon \PP_{{\lambda_1,...,\lambda_k}}^r \to \PP_{{\lambda_1,...,\lambda_k}}^{r+1}, 0 \leq j \leq r $, defined inductively as follows:
	Assume that $\Phi^j$ have been defined inductively for $0 \leq r \leq s-1$, with the property that for every $\un{x}_{G} \in \PP_{{\lambda_1,...,\lambda_k}}^r$ with underlying da-graph $G$, there exists $\un{y}_v, v \in V(G)$ such that $\Phi^j(\un{x}_{G}) = \mu_{G} ( \prod_{v \in V(G)} \un{y}_v )$. We define them for $r=s$.
	\begin{itemize}
		\item for $j<s$: let $\un{x}$ be an element in $\PP_{{\lambda_1,...,\lambda_k}}^s$ with the underlying graph $G$. From induction hypothesis we know that $\Phi^j(d_n \un{x}_G) = \mu_G \left( \prod_{v \in V(G)} \un{y}_v \right)$ for some $\un{y}_v$. Then set
		\[ \Phi^j(\un{x}) = \mu_G \left( \prod_{v \in V(G)} \eta \un{y}_v \right) \]
%		\item For $v \in V(G)$, let $\un{y}_v \in \PP_{{\lambda_1,...,\lambda_k}
%		}^{n-1}$ be the label of vertex $v$ in $\un{x}_G$. The define $\phi^i(\un{x}_G)$ as 
%		\[ \phi^i(\un{x}_G) : = \mu_{G} \left( \prod_{v \in V(G)} \eta \un{y}_v  \right) \]
		\item for $i=k$: Suppose that $c_n(\un{x}) = \mu_G \left( \prod \un{x}_v \right)$. Then define 
		\[ \Phi^j(\un{x}_G) := \mu_G \left( \prod_{v \in V(G)} \eta(\un{x})\right)\]
	\end{itemize}
\end{itemize}
%
%Here $c \colon \PP_{\lambda_1,...,\lambda_k} \to \PP_{\lambda_1,...,\lambda_k}$ is the cut map. As before it is the geometric realization of a map of simplcial spaces given by $c_r \colon \PP^r_{\lambda_1,...,\lambda_n} \to\PP^r_{\lambda_1,...,\lambda_n}$. Maps $c_r$ are defined inductively as follows:
%	\begin{itemize}
%		\item $r=0$: Let $\un{x} \in \PP^0_{\lambda_1,...,\lambda_n}$ be an element with underlying graph $G$. Then $c_0(\un{x})$ is defined by the same expression as in \ref{c0eq}, where now $\un{y}_v, v \in V(G)$ are defined as follows: let $\un{x}_v$ denote the label of $v$ in $\un{x}$. Then $\un{y}_v$ is defined to be the element of $\PP^0_{\lambda_1,...,\lambda_n}$ with underlying da-graph $K_v$ as above, with the labels of vertices in $K_v$ given by the corresponding connected components of $\un{x}\setminus \partial \bigcup_i B^i_N$.
%		\item The inductive step constructing map $c_r$ assuming the existence of $c_i$ for $i <r$ is defined exactly as in \ref{}, with $\PP^i_{\Sigma}, H_v$, and $A_\gamma$ there replaced by $\PP^i_{\lambda_1,...,\lambda_k}, K_v$, and $A^n_j(\un{\Sigma})$.
%	\end{itemize}

	\paragraph{\textbf{The Homotopy $\Psi$}} 
	Finally we have a homotopy $\Psi$ form $c$ to the map $\chi_n \circ \pi$. The homotopy is again defined exactly as in \ref{psiproprmain}. Subcomplex $\PP_{\Sigma}^M$ is replaced by the subcomplex $\PP_{\lambda_1,...,\lambda_k}^n$ given by simplices with labels $\un{x}$, such that for a vertex $v$ with label $\un{x}_v$ in $\un{x}$, $\pi(\un{x}_v) \cap A^n_j(\un{\Sigma}) \neq \varnothing \Rightarrow \un{x}_v = \eta^r( A^n_j(\un{\Sigma}))$. The section $\chi^M_n$ is replaced by $\chi_n$, and the graph $G_{\pi(\un{x})}$ is replaced by $G_{\lambda_1,....,\lambda_k}(\un{x})$. Homotopy $\Psi$ is then defined using extra degeneracies and augmentations as in \ref{extradegpsi}.

	\paragraph{Contraction homotopy proving condition \ref{}}  This is just the concatenation of the cut homotopy $\Phi$, from identity to the cut map, and homotopy $\Psi$, from $c$ to $\chi_n\circ \pi$.
	
	This completes the proof of the condition \ref{} and hence of the proof of the Theorem \ref{}.
\section{From $\Mfrfull$ to $\Mbar$}
We now prove that trivilizing the circle action in $\mfrfull$ in a suitable manner gives the Deligne-Mumford properad $\Mfrfull$.

\begin{theorem}\label{mbarmainthm}
	\[ 
		\begin{tikzcd}
			\Mfrfull^{unst} \ar{r} \ar{d} & \Mbar^{unst} \ar{d}\\
			\Mfrfull \ar{r} & \Mbar
		\end{tikzcd}
	\]
	Is a homotopy pushout diagram.
\end{theorem}

%\begin{remark}
%Here we deal with the reduced versions of properads $\Mfrfull, \Mfrfull^{unst}$ which contain no degree $(0,0)$ components. There is a corresponding theorem for \emph{unreduced} properads. The proof provided here carries over to the unreduced case, with suitable modifications to account for the fact that the moduli $\Mmod^{fr}_{g,0,0}$ contains stacky points (see \cite{}).
%\end{remark}

The strategy of the proof is again very similar to that of the proof of \ref{mthm1}. As outlined in Remark \ref{} to account for the appearance of composition of a sequence of nodal cylinders, some changes are needed when constructing the open sets and filtrations as in \ref{}. In the rest of the section carry out this proof.

Using Proposition \ref{compareproposition}, and argument analogous to those in the discussion following it, we know that the homotopy pushout of 
	\[ 
		\begin{tikzcd}
			\Mfrfull^{unst} \ar{r} \ar{d} & \Mbar^{unst} \\
			\Mfrfull &
		\end{tikzcd}
	\]
coincides with 
\[ \left| \GG \left( \GG^{\bullet}\Mfrfull \coprod_{\GG^{\bullet}\Mfrfull^{unst}} \GG^{\bullet} \Mbar^{unst} \right) \right| \]
In the next section we show that this has the same homotopy type as $\Mbar$

\subsection{$\left| \GG \left( \GG^{\bullet}\Mfrfull \coprod_{\GG^{\bullet}\Mfrfull^{unst}} \GG^{\bullet} \Mbar^{unst} \right) \right| \to \Mbar$ is a weak-equivalence}
We prove that as in \ref{finalhtpyeq},  every $\Sigma \in \Mbar$ has a neighborhood $U_{\Sigma}$ such that for every finite collection $\Sigma_1,...,\Sigma_k$, 
\begin{equation}
	\label{finintwheq} \pi^{-1}(U_{\Sigma_1,...,\Sigma_k}) \to U_{\Sigma_1,...,\Sigma_k} \mbox{ is a weak homotopy equivalence}
\end{equation}
where $U_{\Sigma_1,...,\Sigma_k} = U_{\Sigma_1} \cap ... \cap U_{\Sigma_k}$.

Moreover, we shall show that open sets $\pi^{-1}(U_{\Sigma})$ have a covering $\{W^{n}_{\Sigma}\}_{n \in \NN}$ such that for every $I \subset \NN$ non-empty, $W_{I,\Sigma} = \bigcap_{n \in I} W_{n,\Sigma}$ satisfies the following:
\begin{enumerate}\label{Wsigmacond_red}
	\item $W_{I,\Sigma}$ is fiberwise non-empty. 
	\item $W_{I,\Sigma} \hookrightarrow \pi^{-1}(U)$ admits a fiberwise null homotopy. 
\end{enumerate}
In particular, for any finite collection $\Sigma_1,...,\Sigma_k$, for $U_{\Sigma_1,...,\Sigma_k} \subset U_{\Sigma_1}$,  $\pi^{-1}(U_{\Sigma_1,...,\Sigma_k})$ has an open cover $\{W_{n,\Sigma_1} \cap \pi^{-1}(U_{\Sigma_1,...,\Sigma_k})\}$ such that every finite intersection is non-empty, connected and and contractible in $\pi^{-1}(U_{\Sigma_1,...,\Sigma_k})$. A filtration as in the Condition \ref{filtrantncondition} can then be obtained by setting: 
\[F^n \pi^{-1}(U_{\Sigma_1,...,\Sigma_k}) = \bigcup_{i=1}^n W_{n,\Sigma} \cap \pi^{-1}(U_{\Sigma_1,...,\Sigma_k})\]
with the corresponding section $\chi_n$ as defined in \ref{mthm1finalparag}.	

\subsubsection{Construction of $U_{\Sigma}$ and $F^i \pi^{-1}(U_{\Sigma})$}\label{gammasubsubsect_red}
The construction of $U_{\Sigma}$ is exactly as in \ref{gammasubsubsect}, the only difference being this time we work with $\Sigma$ in $\Mmodbar_{g,n_-,n_+}$ instead of $\Mmodhat_{g,n_-,n_+}$.

The construction of the open sets $W_{i,\Sigma}$ is as follows:
For $\un{x} \in \GG^{n+1} \Mfr \coprod_{\GG^{n+1}  \Mfrunst} \GG^{n+1} \Mhatunst $ define $\gamma_{\un{x}}$ as in \ref{gammasubsubsect}. For each $n$ let $K_n \subset \cfam_{\Sigma}$ be the open set defined by 
\[ K_n = \bigcup_{\nu \in \Nu} \{ (z, \un{t}) \in W_{\nu} \times \DD^{\Nu} | |t_\nu| \leq \frac{1}{2(n+1)^2}, |z_{\hat{\nu}}(z) | \in [ t_{\nu} n, t_{\nu} (n+1) ] \cup [\frac{1}{n+1}, \frac{1}{n}] \} \]
Define $\WW_{n,\Sigma}$ to be the subcomplex of $\GG^{n+1} \Mfr \coprod_{\GG^{n+1}  \Mfrunst} \GG^{n+1} \Mhatunst$ consisting of the simplices $\un{x}$ such that the associated embedding $\gamma_{\un{x}}$ is disjoint from $K_i$. Define $W_{n,\Sigma}$ to be the open subset of $\left| \GG^{n+1} \Mfr \coprod_{\GG^{n+1}  \Mfrunst} \GG^{n+1} \Mhatunst \right|$ associated with this subcomplex.






%%%%%%%%%%%%%%%%%%%%%%%%%%%%%%%%%%%%%%%%%%%%%%%%%%%%%%
%%%%%%%%%%%%%%%%                      Figure                 %%%%%%%%%%%%%%%%%%%%%%%

%\begin{figure} 
%	\centering
%		\includegraphics[width=\textwidth]{../img/gamma_fig.png}
%		\caption{}
%	\end{figure}

%%%%%%%%%%%%%%%%%%%%%%%%%%%%%%%%%%%%%%%%%%%%%%%%%%%%%%
%%%%%%%%%%%%%%%%%%%%%%%%%%%%%%%%%%%%%%%%%%%%%%%%%%%%%%


%%%%%%%%%%%%%%%%%%%%%%%%%%%%%%%%%%%%%%%%%%%%%%%%%%%%%%
%%%%%%%%%%%%%%%%                      Figure                 %%%%%%%%%%%%%%%%%%%%%%%

%\begin{figure} 
%	\centering 
%		\includegraphics[width=\textwidth]{../img/filteration_fig.png}
%		\caption{}
%
%\end{figure}

\begin{figure} 
	\centering
	\begin{subfigure}[b]{0.3\textwidth}
		\includegraphics[width=\textwidth]{../img/filteration_fig.png}
		\caption{Open sets $W_{n,\Sigma}$}
		\label{filtrationfig1}
	\end{subfigure}
	\begin{subfigure}[b]{0.35\textwidth}
		\includegraphics[width=\textwidth]{../img/filteration_fig2.png}
		\caption{Circles $S_{\hat{\nu}}, S_{\check{\nu}}$}
		\label{filtrationfig1}
	\end{subfigure}
	\caption{}
\end{figure}


%%%%%%%%%%%%%%%%%%%%%%%%%%%%%%%%%%%%%%%%%%%%%%%%%%%%%%
%%%%%%%%%%%%%%%%%%%%%%%%%%%%%%%%%%%%%%%%%%%%%%%%%%%%%%
%%%%%  																						  %%%%% 
%%%%% 							Subsection Not needed ?										   %%%%% 
%%%%% 																						   %%%%% 
%%%%%%%%%%%%%%%%%%%%%%%%%%%%%%%%%%%%%%%%%%%%%%%%%%%%%%
%%%%%%%%%%%%%%%%%%%%%%%%%%%%%%%%%%%%%%%%%%%%%%%%%%%%%%
%%%%%%%%%%%%%%%%%%%%%%%%%%%%%%%%%%%%%%%%%%%%%%%%%%%%%%


\subsubsection{Proof of condition (1) in \ref{Wsigmacond}}\label{cond1}
$I \subset \NN$ be a finite set. Pick any $M \in \NN$, $M> max(I)$ and a $\Sigma \in U_{\Sigma}$ whose co-ordinates $(z(\Sigma),\un{t}(\Sigma)) \in  U_{\Sigma} = U_{\hat{\Sigma}} \times \DD^{\Nu}$ satisfy,
\[ |t_\nu(\Sigma)| \leq \frac{1}{2n^2} \mbox{ for every } \nu \]
Consider an embedding $\gamma \colon \bigsqcup_{\Nu} (S^1 \sqcup S^1) \to \Sigma$ given by the circles
\[ S^{\hat{\nu}}(\Sigma) = \{(z,t_{\nu}(\Sigma)) | |z_{\nu}(z)| = \frac{1}{M} \} \subset \Sigma, S^{\check{\nu}}(\Sigma) = \{(z,t_{\nu}(\Sigma))  |z_{\nu}(z)| = t_{\nu}(\Sigma)M \} \]
There is an unique element $\un{x} \in \GG^1 \Mfr \coprod_{\GG^1  \Mfrunst} \GG^1 \Mbarunst$, such that the curve $\gamma$ coincides with the associated embedding $\gamma_{\un{x}}$. Then labels of the vertices of $\un{x}$ are given by the connected components of $\Sigma \setminus Im(\gamma)$, with the edges corresponding to all circles $S^{\hat{\nu}}(\Sigma), S^{\check{\nu}}(\Sigma)$ directed away from the vertex associated to the component $\{(z,t_{\nu}(\Sigma) |  |z_\nu(z)| \in [t_\nu M , \frac{1}{M}] \}$ bounded by $S^{\hat{\nu}}(\Sigma), S^{\check{\nu}}(\Sigma)$. Note that $\un{x} \in \bigcap_{I}W_{n,\Sigma}$. This proves the fiberwise non-emptiness condition. 


We now turn to the proof of condition \ref{filtrprop}. Let us start by defining the section $\chi$ appearing in the condition:\\
Fix any $M \in I$. Let $\Sigma \in \pi(W_{I,\Sigma_0})$. Let $S^{\hat{\nu}}$, and $S^{\check{\nu}}$ be defined as in \ref{cond1}. Denote by $A_{\nu}(\Sigma)$ the annulus bounded by $S^{\hat{\nu}}$ and $S^{\check{\nu}}$ in $\Sigma$.


\subsubsection{Construction of section $\chi_n \colon U_{\Sigma} \to \pi^{-1}(U_{\Sigma})$ from Condition \ref{filtrprop}}
For $\Sigma^{\prime} \in U_{\Sigma}$, consider the graph $G_{\Sigma^{\prime}}$ with 
\begin{itemize}
	\item The set of vertices given by the connected components of $\Sigma^{\prime} \setminus \bigcup_{\Nu} \partial A_{\nu}(\Sigma^{\prime})$
	\item The edges correspond to the the connected components of $\bigcup_{\Nu} \partial A_{\Nu}$, with all edges directed away from vertices labeled by the annuli $A_{\Nu}$.
\end{itemize}
let $\un{x}_{\Sigma^{\prime}}$ be the $0$-simplex with the underlying graph $G_{\Sigma}$ given by labeling each vertex by the corresponding connected component of $\Sigma^{\prime} \setminus \partial A_{\nu}$. Set 
$\chi_n(\Sigma^{\prime}) = \un{x}_{\Sigma^{\prime}}$. It is not difficult to see that this indeed defines a continuous section of $\pi|_{\pi^{-1}}(U_{\Sigma})$.

\subsubsection{Proof of  condition \eqref{filtrprop}}\label{prfcondfiltrtop}
Let us now turn to the fiberwise contraction. We begin by describing a number of preliminary notions and maps:\\
For any $\Sigma^{\prime} \in U_{\Sigma_0}$ whose co-ordinates $(z(\Sigma^{\prime}),\un{t}(\Sigma^{\prime})) \in  U_{\Sigma} = U_{\hat{\Sigma}} \times \DD^{\Nu}$ satisfy,
\[ |t_\nu(\Sigma^{\prime})| \leq \frac{1}{2n^2} \mbox{ for every } \nu \]
Consider an embedding $\gamma \colon \bigsqcup_{\Nu} (S^1 \sqcup S^1) \to \Sigma$ given by the circles
\[ S^{\hat{\nu}}(\Sigma^{\prime}) = \{(z,t_{\nu}(\Sigma^{\prime})) | |z_{\nu}(z)| = \frac{1}{n} \} \subset \Sigma^{\prime}, S^{\check{\nu}}(\Sigma^{\prime}) = \{(z,t_{\nu}(\Sigma^{\prime})) | |z_{\nu}(z)| = t_{\nu}(\Sigma^{\prime}) n \} \]
There is an unique element $\un{x} \in \GG^1 \Mfr \coprod_{\GG^1  \Mfrunst} \GG^1 \Mhatunst$, such that the curve $\gamma$ coincides with the associated embedding $\gamma_{\un{x}}$. Then labels of the vertices of $\un{x}$ are given by the connected components of $\Sigma \setminus Im(\gamma)$, with the edges corresponding to all circles $S^{\hat{\nu}}(\Sigma), S^{\check{\nu}}(\Sigma)$ directed away from the vertex associated to the component $\{(z,t_{\nu}(\Sigma) |  |z_\nu(z)| \in [t_\nu n , \frac{1}{n}] \}$ bounded by $S^{\hat{\nu}}(\Sigma), S^{\check{\nu}}(\Sigma)$.\\
	\paragraph{\textbf{The graphs $H_v$}}\label{hvgrphbar}\textcolor{red}{Need to explain why result of applying $H_v$ to every vertex if of type `at least one output' cf. correction\_1.png}.\\
	These are analogous to the graphs defined in \ref{hvgrph}. Let $\un{x} \in \WW^{i}_{\Sigma}$ with underlying da-graph $G$. If $\un{x}_v$ is the label of $v \in V(G)$ in $\un{x}$, graph $H_v$ as follows: The vertice and edges of $H_v$ as defined exactly as in \ref{hvgrph}, except that in this case the edges are directed towards the vertex corresponding to a component of the form $\cap A_{\nu}$ for some $\nu$.
%		\begin{itemize}
%			\item The vertices of $H_v$ are given as in \ref{hvgrph} i.e. by the connected components of $\Sigma^{\prime}_V \setminus \bigcup_{\Nu} \partial A_{\nu}$
%			\item Edges of $H_v$ are given by the connected components of $\bigcup_{\Nu} \partial A_{\nu}$ lying inside $\Sigma^{\prime}_V$, with each edge directed towards the 
%		\end{itemize}

The construction of cut maps, cut homotopies, and the contraction proving Condition \ref{} carries over verbatim to this case once the definition of graphs $H_v$ is adopted as above.\\
This concludes the proof of condition Condition \eqref{filtrprop}, and hence of Theorem \ref{mbarmainthm}
%		
%	\paragraph{\textbf{The cut maps}}\label{cutmaps} We define a simplical map $c \colon \WW^{i}_{\Sigma} \to \PP_{\Sigma}$ which will be used in the definitions of homotopies below. Intuitively, it is corresponds to making `cuts' to the underlying riemann surface along the cirlces $S^{\hat{\nu}}$ and $S^{\check{\nu}}$.
%	Maps $c_r \colon (\WW^{i}_{\Sigma_0})^r \to \PP_{\Sigma_0}^r$ are defined inductively as follows:
%		\begin{itemize}
%			\item Base case: Let $\un{x} \in \PP_{\Sigma}^0$. Let $G$ be the graph underlying $\un{x}$. For $v \in G$ let $\Sigma_v$ and $H_v$ be as described above in \ref{hvgrph}. Consider an element $\un{y}_v \in \PP^1_{\Sigma_v}$ whose underlying da-graph is $H_v$ and whose vertex labels are given by the corresponding connected components of $\Sigma_v \setminus \bigcup_{\Nu} \partial A_{\nu}$ as in \ref{hvgrp}.  We then define $c_0(\un{x})$ as the graph composition 
%			\begin{equation}\label{c0eq} c_0(\un{x}) := \mu_G( \prod \un{y}_v ) \end{equation}
%			\item Assume inductively that $c_i, 0 \leq i <r$ has been defined and satisfies the following property: For any $\un{x} \in \PP^i_{\Sigma}$, with the underlying graph $G$, $c_{i} (\un{x}) = \mu_{G}( \prod \un{y}_v)$ for some $\un{y}_v \in \PP^i_{\Sigma_v}$. Notice that the condition is satisfied for $i=0$.\\
%		Now suppose $\un{x} \in \PP^r_{\Sigma}$, having the underlying da-graph $G$. Consider element $c_{r-1}(d_r \un{x}) \in \PP^{r-1}_{\Sigma}$. From the induction hypothesis we know that $c_{r-1}(d_r \un{x}) =  \mu_{G_{d_r \un{x}}}(\prod \un{y}^{\prime \prime}_{w})$ for some $\un{y}^{\prime \prime}_w$, $w \in V( G_{d_r\un{x}} )$ where $G_{d_r\un{x}}$ is the da-graph underlying $d_r\un{x}$. But since  $G_{d_r \un{x}}  = \mu_G( \prod G_v)$ for some graphs $G_v$, we can express $c_{r-1}(d_{r-1} \un{x})$  as a graph composition $\mu_G( \prod \un{y}^{\prime}_v )$ for some $\un{y}^{\prime}_v \in \PP_{\Sigma}^{r-1}$. For $v \in G$, consider an element $\un{y}_v \in \PP^n_{\Sigma}$ defined as follows: the da-graph underlying $\un{y}_v$ is $H_v$..The label of a vertex $u \in  H_v$ in $\un{y}_v$ is the labeled subgraph of $\un{y}^{\prime}_v$ consisting of vertices and edges mapping under $\pi$ into the connected component $\Sigma_u$ of $\Sigma \setminus \bigcup_{\Nu}\partial A_{\nu}$ corresponding to $u$ as in \ref{hvgrph}. Then, we set 
%		\[c_r(\un{x}) = \mu_G( \prod \un{x}_v)\]
%		\end{itemize}
%
%\paragraph{\textbf{The cut homotopies}} We define a fiber preserving homotopy $\Phi \colon F^i \pi^{-1}(U_{\Sigma})\to \pi^{-1}(U_{\Sigma})$. This is defined as the geometric realization of a simplicial homotopy $\phi \colon \WW^{i}_{\Sigma} \to \PP_{\Sigma}$. This in turn consists of maps $\phi^i \colon (\WW^{i}_{\Sigma})^{s} \to \PP^{s+1}_{\Sigma}$ for $0 \leq i \leq s$, satisfying the usual conditions of a simplicial homotopy.\\
%Assume inductively that the maps $\phi^i$ have been defined for all degrees $s \leq r-1$ and satisfy the following property similar to the one used above in \ref{cutmaps}: For any $\un{x} \in (\WW^{i}_{\Sigma})^{s}$, with the underlying graph $G$, $\phi^{i} (\un{x}) = \mu_{G_{\un{z}}}( \prod \un{y}_v)$ for some $\un{y}_v \in  (\WW^{i}_{\Sigma})^{s}$. The maps $\phi^i \colon (\WW^{i}_{\Sigma})^{s} \to \PP^{s+1}_{\Sigma}$ are defined 	as follows:\\
%Let $\un{x} \in  (\WW^{i}_{\Sigma})^{r}$ and $G$ be the da-graph underlying $\un{x}$.
%	\begin{itemize}
%		\item $i <r$: From the condition in the inductive hypothesis we have $\phi^i(d_r\un{x}) = \mu_G( \prod \un{y}_v)$, for some $\un{y}_v$. Then define
%		 \[\phi^i(\un{x}) := \mu_G ( \prod \eta(\un{y}_v) )\]
%		recall from \ref{GGtripsubsec}, $\eta \colon \mathbbm{1}_{\TopSeq} \to \GG$ is the unit of the monad $\GG$.
%		\item $i=r$: For $\un{x} \in  (\WW^{i}_{\Sigma})^{r}$, suppose that $c_r(\un{x}) = \mu_G(\prod \un{y}_v)$. Then define
%		\[ \phi^i(\un{x}) := \mu_G( \prod \eta(\un{y}_v)) \]
%	\end{itemize}
%
%\paragraph{\textbf{The homotopy $\Psi$}}\label{psiproprmain} Let $\PP^M_{\Sigma}$ be the subcomplex of $\PP_{\Sigma}$ consisting of simplices $\un{x}$ such that if $\un{x}_v$ is the label of a vertex in $\un{x}$, $\pi(\un{x}_v) \cap A_{\nu} \neq \phi \Rightarrow \un{x}_v = \eta^r(A_{\nu})$ for some $\nu \in \Nu$.
%
%% and $\nu \in \Nu$, then $\pi(\un{x}_v) \cap A_{\nu} \neq \phi \Rightarrow \pi(\un{x}_v) \subset A_{\nu}$. Note that $Im(c) \subset \PP^M_{\Sigma_0}$. Let $\pi^{-1}(U_{\Sigma_0})^M$ denote the geometric realization of $\PP^M_{\Sigma_0}$.\\
%%Let $\Ss^M_{\Sigma_0}$ be a further subcomplex contained inside the $0$-skeleton of $\PP^M_{\Sigma_0}$. It consists of $0$-simplies $\un{x}_{\Sigma}, \Sigma \in U_{\Sigma_0}$ defined as follows: the underlying graph $G_{\Sigma}$ has 
%%	\begin{itemize}
%%		\item The set of vertices given by the connected components of $\Sigma \setminus \bigcup_{\Nu} \partial A_{\nu}$
%%		\item The edges correspond to connected components of $\bigcup_{\Nu} \partial A_{\nu}$, with all edges directed away from vertices labeled by the annuli $A_{\nu}$.
%%	\end{itemize} $\un{x}_{\Sigma}$ is obtained by labeling each vertex by the corresponding connected component of $\Sigma \setminus \bigcup_{\Nu} \partial A_{\nu}$. $S^M_{\Sigma_0} \subset \pi^{-1}(U_{\Sigma_0})^M$ be its geometric realization. $\Ss^M_{\Sigma_0}$ is the constant simplcial space associated with $S^M_{\Sigma_0}$. Also, notice that $S^M_{\Sigma_0}$ defines a (continuous) section of the map $\pi \colon \pi^{-1}(U_{\Sigma_0}) \to U_{\Sigma_0}$.\\
%
%We will construct a fiber-preserving homotopy $\Psi$ contracting $\pi^{-1}(U_{\Sigma_0})^M$ into $S^M_{\Sigma_0}$. This will again be a geometric realization of a simplcial homotopy, which itself will be constructed by specifying extra degeneracies and augmentation on $\PP^M_{\Sigma_0}$ (cf. \cite{J.-P. Meyer. Bar and cobar constructions. I, Sec. 6}). For this we need to define:
%	\begin{enumerate}
%	 \item Extra degeneracies: $s_{n+1} \colon (\PP^M_{\Sigma_0})^n \to (\PP^M_{\Sigma_0})^{n+1}$ for all $n \geq -1$, where $(\PP^M_{\Sigma_0})^{-1} := S^M_{\Sigma_0}$
%	 \item Augmentations: $d_0 \colon (\PP^M_{\Sigma_0})^0 \to S^M_{\Sigma_0}$
%	\end{enumerate}
%satisfying:
%	\begin{enumerate}
%		\item $d_{n+1}s_{n+1}$ is identity on $(\PP^M_{\Sigma_0})^n$
%		\item for $n \geq 1, 0 \leq i,j \leq n$, 
%			\begin{enumerate}
%			\item $d_{i}s_{n+1} = s_{n}d_i$
%			\item $s_{j+1}s_{n+1} = s_{n+2}s_j$
%			\end{enumerate}
%		on $(\PP^M_{\Sigma_0})^n$
%	\end{enumerate}
%	The contracting homotopy $\psi$ underlying $\Psi$ is then obtained by setting $\psi_i \colon (\PP^M_{\Sigma_0})^n \to (\PP^M_{\Sigma_0})^{n+1}$, $i \leq n$, to be $s_n \hdots s_{i+1}d_{i+1} \hdots d_n s_{n+1}$.
%	
%	\begin{itemize}\label{extradegpsi}
%	\item Extra degeneracies: We define $s_{n+1} \colon (\PP^M_{\Sigma_0})^n \to (\PP^M_{\Sigma_0})^{n+1}$. Let $\un{x} \in (\PP^M_{\Sigma_0})^n$. Let $G_{\pi(\un{x})}$ be the graph defined in \ref{}. From the definition of $\PP^M_{\Sigma_0}$ it follows that the vertices of $\un{x}$ can be regrouped into subgraphs $\un{y}_v \in  (\PP^M_{\Sigma_0})^n , v\in G_{\pi(\un{x})}$ so that 
%$\un{x} = \mu_{G_{\pi(\un{x})}} ( \prod \un{y}_v)$.
%	Then  set 
%	\[ s_{n+1}(\un{x}) := \mu_{G_{\un{x}}} ( \prod \eta(\un{y}_v)) \]
%	\item Augmentation: Note that for $\un{x} \in  (\PP^M_{\Sigma_0})^0 $, $d_0(s_1\un{x}) = \un{x}_{\pi(\un{x})} \in \Ss^M_{\Sigma_0}$. We define $d_0 \colon (\PP^M_{\Sigma_0})^0 \to S^M_{\Sigma_0}$ as
%	\[ d_0 = d_0 \circ s_1 \]
%	\end{itemize} 
%
%\paragraph{\textbf{The homotopy $\widetilde{\Psi}$}}\label{psihtpy} We construct another fiber-preserving homotopy, this time contracting $\pi^{-1}(U^{ns}_{\Sigma})$ onto a section $\widetilde{S}_{\Sigma}$ of $\pi$ over $U^{ns}_{\Sigma}$. Here $U^{ns}_{\Sigma} = \{(z,\un{t})| t_{\nu} \neq 0 \mbox{ for all } \nu\} \subset U_{\hat{\Sigma}} \times \DD^{\Nu} = U_{\Sigma}$ is the subspace of non-singular curves in $U_{\Sigma}$. Let $\PP^{ns}_{\Sigma}$ be the corresponding subcomplex of $\PP_{\Sigma}$.\\
%The section $\widetilde{S}_{\Sigma}$ is contained in the $0$-skeleton of $\PP_{\Sigma}$ and is given by $\Sigma \mapsto \tilde{\un{x}}_{\Sigma}$, where the graph underlying $\tilde{\un{x}}_{\Sigma}$ has a single vertex, no internal edges, and with the label of the vertex being simply $\Sigma$. Note that since $\Sigma$ is non-singular, $\tilde{\un{x}}_{\Sigma}\in \GG \Mfr \subset (\PP_{\Sigma})^0$.
%
%
%Let the subcomplex of  $\widetilde{\PP}_{\Sigma}$ corresponding to $\widetilde{S}_{\Sigma}$ be denoted by $\widetilde{\Ss}_{\Sigma}$. It is a constant simplcial space and is contained inside the $0$-skeleton of $\PP_{\Sigma}$.
%Again we shall construct the contraction by specifying extra degeneracies and augmentation as in \ref{psihtpy}. We continue using the same notation for these as in \ref{}.
%	\begin{itemize}
%		\item Extra degeneracies: In this case the map $s_{n+1} \colon (\PP^{ns})^n \to  (\PP^{ns})^{n+1}$ is defined by 
%		\[ s_{n+1}(\un{x}) _= \eta(\un{x}) \]
%		\item The augmentation $d_0 \colon  (	\PP^{ns})^0 \to \widetilde{\Ss}_{\Sigma}$ is defined by 
%		\[ d_0 := d_0 \circ s_1 \]
%	\end{itemize}
%The verification of conditions satisfied by $s_{n+1}$'s and $d_0$ is analogous to \ref{psihtpy} and is omitted. The simplcial homotopy $\widetilde{\psi}$ underlying $\tilde{\Psi}$ is obtained from $s_{n+1}$'s and $d_0$ by formulas similar to those in \ref{psihtpy}.
%
%\paragraph{\textbf{The contraction for proving Condition \ref{filtrprop}}}
%We are finally ready to define the homotopy needed to show that Condition \ref{filtrprop} is satisfied:	We construct a homotopy $\Omega$ from the inclusion $F^n\pi^{-1}(U_{\Sigma}) \hookrightarrow \pi^{-1}(U_{\Sigma})$ to a map which takes $F^n\pi^{-1}(U_{\Sigma})$ onto a section of $\pi$ over $U_{\Sigma}$. The section coincides with $\chi_i$ near the singular locus $U_{\Sigma} \setminus U^{ns}_{\Sigma}$, and with $\widetilde{S}_{\Sigma}$ away from it.
%
%Fix a continuous function $\zeta \colon U_{\Sigma} \to [0,1]$ such that $\zeta \equiv 1$ on $\bigcup_{\Nu} \{|t_{\nu}|\leq \frac{1}{4n^2}\}$ and $\zeta \equiv 0$ outside of $\bigcup_{\Nu} \{|t_{\nu}|\leq \frac{1}{2n^2}\}$. Set $\kappa = \zeta \circ \pi$. The homotopy $\Omega$ is defined as follows:
%
%	\begin{align*}
%	\mbox{ On } \pi^{-1}(\bigcup_{\Nu} \{|t_{\nu}|\leq \frac{1}{4n^2} \})&: \Omega(\un{x},t) = 
%		\begin{cases} 
%			\Phi(\un{x},2t), & t \leq 1/2 \\
%			\Psi(\Phi(\un{x},1), 2t-1) & t >1/2
%		\end{cases}\\
%	\mbox{ On } \pi^{-1}(U^{ns}_{\Sigma})&: \Omega(\un{x},t) = 
%		\begin{cases} 
%			\Phi(\un{x},2t), & t \leq \kappa(\un{x})/2 \\
%			\Phi(\widetilde{\Psi}(\un{x},\frac{t-\kappa/2}{1-\kappa/2}), \kappa(\un{x})) & t > \kappa({\un{x}})/2
%		\end{cases}
%	\end{align*}
%	To ensure that $\Omega$ is well-defined it suffices to check that the two definitions agree on points lying over $U^{ns}_{\Sigma} \bigcap \bigcup_{\Nu} \{|t_{\nu}|\leq \frac{1}{4n^2} \}$
%	\[ \Psi ( \Phi (\un{x},1),t) = \Phi (\widetilde{\Psi} ( \un{x} , t) ,1 ) \mbox{ for all } t\]
%Since $\Phi(\un{x},1)=g(\un{x})$, this amounts to checking $\psi_i c_n = c_{n+1} \widetilde{\psi}_{i}$ on $n$-simplices lying over $U^{ns}_{\Sigma} \cap \bigcup_{\Nu} \{|t_{\nu}|\leq \frac{1}{4n^2}|\}$. Since $g$ is simplcial, from the definitions of $\psi_i,\widetilde{\psi}_i$ we are reduced to checking 
%\[ \psi_n c_n =  c_{n+1} \widetilde{\psi}_{n}\]
%
%

\printbibliography

\end{document}